\documentclass[a4paper,11pt]{ltjsarticle}
\usepackage[margin=25mm,headheight=16pt]{geometry}
\usepackage{amsmath,amssymb,amsthm,mathtools,bm}
\usepackage{booktabs,longtable,array,enumitem}
\usepackage{xcolor}
\usepackage[unicode,colorlinks=true,linkcolor=blue!40!black,citecolor=blue!40!black,urlcolor=blue!45!black]{hyperref}
\usepackage{fancyhdr}
\usepackage{microtype}
\setlength{\parindent}{1em}
\setlength{\parskip}{0.25em}
\setlist{itemsep=0.2em,topsep=0.4em}
\allowdisplaybreaks[2]
\pagestyle{fancy}
\fancyhf{}
\fancyhead[L]{\small 3次超曲面のPoisson構造と非可換化}
\fancyhead[R]{\small 2026-11-08 / v2}
\fancyfoot[C]{\thepage}
\newcommand{\CC}{\mathbb C}
\newcommand{\ZZ}{\mathbb Z}
\newcommand{\PP}{\mathbb P}
\newcommand{\OO}{\mathcal O}
\newcommand{\cK}{\mathcal K}
\newcommand{\cP}{\mathcal P}
\newcommand{\cD}{\mathcal D}
\newcommand{\cC}{\mathcal C}
\newcommand{\cE}{\mathcal E}
\newcommand{\cF}{\mathcal F}
\newcommand{\cM}{\mathcal M}
\newcommand{\fm}{\mathfrak m}
\newcommand{\Perf}{\operatorname{Perf}}
\newcommand{\Ku}{\operatorname{Ku}}
\newcommand{\MF}{\operatorname{MF}}
\newcommand{\HH}{\operatorname{HH}}
\newcommand{\QGr}{\operatorname{QGr}}
\newcommand{\qgr}{\operatorname{qgr}}
\newcommand{\Dsg}{D_{\mathrm{sg}}}
\newcommand{\RHom}{\operatorname{RHom}}
\newcommand{\REnd}{\operatorname{REnd}}
\newcommand{\Hom}{\operatorname{Hom}}
\newcommand{\Ext}{\operatorname{Ext}}
\newcommand{\Ind}{\operatorname{Ind}}
\newcommand{\Def}{\operatorname{Def}}
\newcommand{\thick}{\operatorname{thick}}
\newcommand{\Spec}{\operatorname{Spec}}
\newcommand{\Spf}{\operatorname{Spf}}
\newcommand{\Id}{\operatorname{Id}}
\newcommand{\sym}{\operatorname{sym}}
\newcommand{\Cone}{\operatorname{Cone}}
\newcommand{\rank}{\operatorname{rank}}
\newcommand{\Pic}{\operatorname{Pic}}
\newcommand{\otL}{\mathbin{\otimes^{\mathbf L}}}
\newcommand{\DocDate}{2026年11月8日}
\newcommand{\DocVersion}{統合第2版}
\theoremstyle{plain}
\newtheorem{theorem}{定理}[section]
\newtheorem{proposition}[theorem]{命題}
\newtheorem{lemma}[theorem]{補題}
\newtheorem{corollary}[theorem]{系}
\theoremstyle{definition}
\newtheorem{definition}[theorem]{定義}
\newtheorem{problem}[theorem]{問題}
\theoremstyle{remark}
\newtheorem{remark}[theorem]{注意}
\renewcommand{\proofname}{証明}
\hypersetup{pdfauthor={},pdfsubject={3次超曲面のPoisson構造・非可換量子化・rotation付き残余圏}}

\hypersetup{pdftitle={3次超曲面のPoisson構造と非可換量子化}}
\title{3次超曲面のPoisson構造と非可換量子化\\
\large 正規形・退化スキーム・rotation付き残余圏}
\author{}
\date{\DocDate・\DocVersion\\\small 公表用草稿}
\begin{document}
\maketitle
\begin{abstract}
滑らかな複素3次超曲面のPoisson構造について、一般次元の正規形、退化スキーム、
および三次の中心元を持つ次数付き量子化を一つの枠組みで扱う。
Poisson正規形は、Hessianに関する線形条件と不変対称形式を持つLie代数の構造から得られる。
全退化スキームの方程式を示し、六次元・七次元の場合の完全交叉、有限被覆、非被約構造を記述する。
さらに滑らかなHesse型ブロックでは代数的な三次中心元を多項式の閉公式で求める。
量子化から可換幾何を復元する問題では、残余dg圏と標準次数シフトに由来するrotationを同時に記録する。
可換cubicの形式近傍ではmarkedな可換モデルが存在し、適切なsmooth proper代数的族では
その存在をétale近傍へ移せる。独立な楕円Sklyaninブロックについては、
高々一つを除く平行移動が有限位数である場合に明示的な可換モデルを与える。

一方、任意のAS-regularな中心三次商に同じ存在結論を要求することはできない。
Fermat cubicのZhang twistにより、射影圏が滑らかな可換cubicと同値であり、
周囲の代数もCalabi--Yauであるにもかかわらず、rotationを保つ可換モデルを持たない四次元の例を構成する。
その障害はrotation不変Hochschild homologyの次元で検出され、
相対三次周期性を持つsmooth properな連結族を通じても消えない。
これにより、一般のAS-regular中心三次商、可換点からの量子化成分、
形式的復元および有限パラメータでの存在を明確に区別する。
\end{abstract}
\tableofcontents

\section{序論}\label{sec:global-intro}

\subsection{射影完全交叉と斉次多様体の中での位置づけ}
次元三以上の滑らかな通常の射影完全交叉では、
非零の正則Poisson構造を持つものは、
射影空間、二次超曲面、三次超曲面、二次二個の完全交叉に限る
（定理\ref{thm:ci-poisson-existence}）。
非存在側には古典的なtwisted differential formsの消滅を用いる
\cite[Lemma 4.1]{FuHwangCI}。
曲面にはK3完全交叉が加わるので次元仮定が不可欠である。
次元二以上のトーリック多様体と滑らかな射影斉次多様体にも
非零のPoisson構造が存在することをSection
\ref{sec:general-poisson-existence}で証明する。
量子化の二次関係式は明示的だが、三次中心元の一般的な
閉公式が既に得られたという意味ではない。
ただし滑らかなHesse型ブロックでは、既知のSklyanin中心元の公式を
量子化パラメータに翻訳することで明示的な中心元と除外すべき
パラメータの判別式を得る（第\ref{sec:explicit-hesse-center}節）。

\subsection{二つの問題と一つの橋渡し}

\(X=V(F)\subset\PP(V)\)、\(\dim V=N=n+2\)、\(n\ge3\) を滑らかな複素3次超曲面とする。
本稿の第一の問題は \(X\) 上のPoisson構造を分類し、その葉と退化を理解することである。
第二の問題は、その量子化 \(B=A/(\Omega)\) から可換cubicの幾何を回復することである。
ここで \(A\) は多項式環と同じHilbert級数を持つ次数付きAS-regular代数、
\(\Omega\) は次数三の中心正則元である。
両者の間には、\(B\) の残余dg圏と、その次数一のシフトから得られる自己同値という橋渡しがある。

可換cubicに対しては
\[
 D^b(X)=\langle\Ku(X),\OO_X,\OO_X(1),\ldots,\OO_X(n-2)\rangle,
 \qquad
 \rho_X=\mathbf L_{\OO_X}\circ(-\otimes\OO_X(1))
\]
と定める。体上の適切な非可換中心三次商でも、graded matrix factorizationとOrlov windowを通じて
残余成分 \(\mathcal K_B\) とrotation \(\rho_B\) を考えられる。
Artin底上では有界coherent圏全体でなく、\(D(\QGr B)^c\) のcompact残余成分を用いる。
本稿の圏同値はdg enhancementを保つMorita同値であり、rotationとの自然同型を含む。
三角圏の同値だけからHochschild不変量の同一視を結論することはしない。

\subsection{Poisson分類と量子化}

\(F\) を固定すると、任意の大域bivectorは一意な定数trivector \(P\in\Lambda^3V\) により
\(\iota_{dF}P\) と表示される。射影空間への延長も一意である。
非零のPoisson構造の場合、線形座標を選べば
\[
 P=\sum_{s=1}^k\partial_{x_{s1}}\wedge\partial_{x_{s2}}\wedge\partial_{x_{s3}},
 \qquad F=\sum_{s=1}^k f_s(x_s,u),\qquad N=3k+m
\]
となり、異なる三変数ブロック間は可換となる。
中心変数 \(u\) と各ブロックとの混合項は許される。
この正規形の証明は、三重ベクトルの次元別軌道分類を用いない。

第I部では、Casimir、Poisson超曲面、一般のsymplectic leafを記述し、
Euler場と \(q^j\) の外積から全退化スキームのイデアルを与える。
平面三次制限が滑らかな場合には有限平坦なcritical coverを用い、
六次元の階数四で現れる192点、七次元の階数四で現れる種数577の曲線、
七次元の階数六でのprimary分解を求める。
特殊な三変数制限が零となる場合とFermat型の非被約構造も扱う。
量子化では、potential代数のPBW性、形式的中心元の存在、有限次数線形代数による代数化を分離する。
この構成だけから \(\qgr B\) のsmoothnessを結論するものではない。
相対例外列の存在を仮定すれば、全超曲面圏と残余圏のsmoothnessが同値となること、
相対Orlov比較の下でrotationの三次周期性が形式的に成立することを
第II部の最後で証明する。

\subsection{可換モデルの存在とその限界}

存在問題を無条件に全ての非可換cubicへ広げると、次の障害がある。

\begin{theorem}[一般的なrotation付き可換復元の反例]\label{thm:overview-negative}
 多項式環と同じHilbert級数を持つnoetherian Koszul AS-regular Calabi--Yau代数 \(A\) と、
 三次中心正則元 \(\Omega\) であって、次を満たすものが存在する。
 \(B=A/(\Omega)\) の射影圏は滑らかな可換cubic fourfold \(X\) のcoherent圏と同値であり、
 残余圏はsmooth properで、\(\rho_B^3\simeq[2]\) を満たす。
 しかし、いかなる滑らかな可換cubic fourfold \(Y\) に対しても
 \[
 (\mathcal K_B,\rho_B)\not\simeq(\Ku(Y),\rho_Y).
 \]
\end{theorem}
\begin{proof}
 \(F=\sum_{i=0}^5x_i^3\)、\(\tau=\operatorname{diag}(\zeta,\zeta^2,1,1,1,1)\)、
 \(\zeta^3=1\ne\zeta\) とし、\(A=\CC[x_0,\ldots,x_5]^\tau\) を右Zhang twistとする。
 中心元には \(\Omega=F\) を取る。
 定理\ref{thm:zhang-cy-fourfold-counterexample}で全ての仮定を検証し、
 \(\dim\HH_0(\mathcal K_B)^{\rho_B}=8\) と
 \(\dim\HH_0(\Ku(Y))^{\rho_Y}=20\) の不一致を証明する。
\end{proof}

さらに、この反例と標準的cubicのpairは、相対的な \(\rho^3\simeq[2]\) を備えた
連結smooth proper代数的族では結べない。
rotation不変Hochschild複体のEuler標数がそれぞれ10と22となり、
この値が連結底上で一定だからである（系\ref{cor:zhang-cy-component-obstruction}）。
従って、以下の肯定的結果は、一般の中心三次商の全成分を覆う主張として解釈してはならない。

\begin{enumerate}
 \item 滑らかなcubicの斉次座標環の次数ごとに有限自由なArtin変形には、
 rotationを保つ一意なmarked可換cubicモデルが存在する
 （定理\ref{thm:main}、\ref{thm:pair}）。
 \item smooth properな代数的dg圏の族とrotationが既に与えられていれば、
 可換cubicのpairとなる点のétale近傍で、cubicの族とpair同値を構成できる
 （定理\ref{thm:etale}）。この結論は一般族の相対saturationを仮定する。
 \item 独立なSklyaninブロックの和では、全ての \(h_j^6\ne-27\) と、
 高々一つを除く平行移動の有限位数性の下で、明示的な可換モデル
 \[
 V\!\left(\sum_j(x_j^3+y_j^3+z_j^3+h_j^2x_jy_jz_j)+P(u)\right)
 \]
 を得る。\(P\) は原点で孤立臨界点を持つ任意の斉次三次式でよい
 （定理\ref{thm:global}、\ref{thm:multiblock-model}）。
 \item 二つ以上の無限位数因子では全生成は未解決だが、
 properness、Serre互換性、\(\HH^0\) の連結性から完全忠実比較を同値へ強める
 十分条件がある（命題\ref{prop:serre-connected-generation}）。
\end{enumerate}

可換点から出発する特定のPoisson量子化族について、全ての良い有限パラメータに
可換pairモデルが存在するか、という問題は残る。
形式的な存在、étale局所的な存在、元の底への降下、全ファイバーでの存在を区別する。
反例を得た一般AS-regularモジュライの問題と、この特定成分の問題を混同しない。

\subsection{先行研究との関係}

三次元のPoisson構造のGrassmann多様体による記述と、potential代数・PBW量子化の
既知の構成を、第I部の文献比較で位置づける。
本稿の一般次元の正規形では、基点をもたないHessian線形系から得るquadratic Lie代数の
可解根基と高階数単純因子を排除することが中心となる。
scheme-theoreticな退化の記述は、この正規形を全外積へ適用することで得られる。

Orlov、Etingof--Ginzburg、Cassidy--Conner--Kirkman--Mooreによる
特異圏・matrix factorization・中心商の比較は既知の基礎結果である
\cite{Orlov,EG,CCKM}。
楕円Sklyaninブロックの残余圏と楕円曲線の関係には、
Cho--Hong--LauおよびLeeの非可換homological mirror symmetryからの先行研究もある
\cite{CHL,Lee}。
本稿ではこの基礎同値と、rotationのtranslation成分を共役で消去する計算、
高次元へ移す外部積の全生成を分けて述べる。

形式復元はHochschild計算と幾何学的Kodaira--Spencer写像
\cite{BFK,HR,JLLZ}、actual変形とformal moduli problemの比較
\cite{Lurie,BKP}、および相対生成子の持上げ\cite{DDL}を組み合わせる。
étale局所実現にはperfect objectの代数的moduli\cite{TV}を用いる。
一方有限中心の外部積生成については、既知の双方有限中心の結果\cite{ELS}との
仮定の差を明示して証明する。

Lin--Zhangのpolarized categorical Torelli\cite[Theorem 1.7]{LZ}は、
二つの可換モデルの一意性を保証する。
任意の非可換pairが可換cubicから生じるという存在命題を保証するものではない。
反例に用いるZhang twistとNakayama公式は既知である\cite{ZhangTwist,RRZ}。
ここではそれらを標準rotationの作用\cite[Proposition 3.12]{Huy}と組み合わせ、
普遍的な可換pair復元を妨げる具体的な不変量を計算する。
これらの組合せを含む本稿全体の先行研究上の新規性は確認継続中である。

\section{射影完全交叉と斉次多様体のPoisson構造}
\label{sec:general-poisson-existence}
本節は三次超曲面のPoisson構造を一般の射影多様体の中に位置づける。
完全交叉に対する非存在の部分はBr\"uckmannの古典的消滅定理
（Fu--Hwang \cite[Lemma 4.1]{FuHwangCI} に引用）による。
本節の既知結果を新規な分類とみなさない。

\subsection{滑らかな射影完全交叉の存在判定}
\begin{theorem}\label{thm:ci-poisson-existence}
複素数体上の滑らかな通常の射影完全交叉
\[
 X=X_{d_1,\ldots,d_c}\subset\PP^{n+c},
 \qquad n\ge3,\quad d_j\ge2
\]
について、以下は同値である。
\begin{enumerate}[label=(\arabic*)]
\item $H^0(X,\bigwedge^2T_X)\ne0$。
\item $X$は非零の正則Poisson構造を持つ。
\item 多重次数は空列、$(2)$、$(3)$、$(2,2)$のいずれかである。
\end{enumerate}
空列は$\PP^n$を表す。一次式の方程式はあらかじめ除いている。
\end{theorem}
\begin{proof}
$s=\sum_{j=1}^c(d_j-1)$とおく。随伴公式から
\[
 \omega_X\simeq\OO_X(-(n+1-s)),\qquad
 \bigwedge^2T_X\simeq\Omega_X^{n-2}(n+1-s)
 \tag{2.1}\label{eq:ci-forms}
\]
である。Br\"uckmannの消滅\cite[Lemma 4.1]{FuHwangCI}は
\[
 H^0(X,\Omega_X^p(k))=0
 \qquad(1\le p\le n-1,\ k\le p)
 \tag{2.2}\label{eq:ci-vanishing}
\]
を与える。$n\ge3$なので$p=n-2\ge1$を取り、
$s\ge3$なら$n+1-s\le n-2$ゆえ
$H^0(X,\bigwedge^2T_X)=0$となる。
したがって(1)なら$s\le2$であり、$d_j\ge2$より列挙した四種類に限る。

逆に$s\le2$の場合、一般の独立な一次形式
$L_0,\ldots,L_{n-2}$を選び、射影空間のEuler場
$E=\sum_i z_i\partial_{z_i}$を用いて
\[
 \alpha=\iota_E(dL_0\wedge\cdots\wedge dL_{n-2})
       \in H^0(X,\Omega_X^{n-2}(n-1))
\]
を定める。一般の選択では$\alpha|_X$は非零である。
$L_0\ne0$上では
\[
 \alpha=L_0^{n-1}\,
 d(L_1/L_0)\wedge\cdots\wedge d(L_{n-2}/L_0)
\]
となり、一般点でその核は有理写像
$X\dashrightarrow\PP^{n-2}$のファイバーに接する
階数二の可積分分布である。
$s=2$なら\eqref{eq:ci-forms}によって$\alpha$を
二重ベクトルに対応させる。
$s=1$なら非零の一次切断、$s=0$なら非零の二次切断を掛け、
同じ同型によって対応させる。
いずれも一般階数二で非零であり、
可積分分布を定める分解可能二重ベクトルなので
Schouten条件$[\pi,\pi]=0$を満たす。
これで(3)から(2)、(2)から(1)が従う。
\end{proof}

\begin{remark}[曲線と曲面]\label{rem:ci-small-dimension}
曲線では$\bigwedge^2T_X=0$。
曲面では$\bigwedge^2T_X=\omega_X^{-1}$で
Schouten条件は自動的に成立する。
滑らかな完全交叉曲面の非零Poisson構造の存在条件は
$s\le3$であり、上記に加えて$(4)$、$(2,3)$、
$(2,2,2)$のK3完全交叉が現れる。
従って$n\ge3$は本質的である。
特異完全交叉には上の消滅をそのまま適用しない。
\end{remark}

\subsection{トーリック多様体}
\begin{proposition}\label{prop:toric-poisson}
次元二以上の複素正規トーリック多様体には
非零の代数的Poisson括弧がある。
滑らかなら$\pi\in H^0(X,\bigwedge^2T_X)$で与えられる。
\end{proposition}
\begin{proof}
稠密トーラス$T\simeq(\CC^*)^n$のLie代数の
一次独立な二元が定める大域導分を$D_1,D_2$とする。
$T$が可換なので$[D_1,D_2]=0$であり、
$\pi=D_1\wedge D_2$はJacobi条件を満たす。
正規トーリックの各不変アフィンチャートで
$\{f,g\}=D_1(f)D_2(g)-D_2(f)D_1(g)$
は正則で、貼り合わさる。
稠密トーラス上では二導分は独立なので非零である。
\end{proof}

\subsection{滑らかな射影斉次多様体}
\begin{theorem}\label{thm:homogeneous-poisson}
$\CC$上の滑らかな連結射影斉次多様体$X$で
$\dim X\ge2$ならば、$X$は非零の正則Poisson構造を持つ。
\end{theorem}
\begin{proof}
Borel--Remmertの定理\cite{BorelRemmert}により
$X\simeq A\times G/P$と分解できる。
$A$はアーベル多様体、$G/P$は半単純群の旗多様体の直積である。
$A$の並進ベクトル場は可換であり、
$G/P$の最大トーラスから来る基本ベクトル場も可換である。
後者で非自明な単純因子の階数が二以上ならば、
$P$に含まれない単純根$\alpha$と、
それに隣接する単純根$\beta$に対し、
$\alpha$および$\alpha+\beta$は接空間の
一次独立な二つのトーラスウェイトを与える。
従って最大トーラスの一般軌道は二次元以上である。
階数一の非自明因子は$\PP^1$であり、
二つ以上のそのような因子があれば同様に二方向がある。
残る場合、$\dim X\ge2$から$A$の並進方向を
必要な数だけ加えられる。
以上より$X$上に一般点で一次独立な二つの
可換な大域ベクトル場$D_1,D_2$が存在する。
その外積は非零で
$[D_1\wedge D_2,D_1\wedge D_2]=0$を満たす。
\end{proof}
\begin{remark}
非射影の任意の斉次空間に対して同じことを
示したわけではない。
次元一では二重ベクトルは常に零となる。
\end{remark}

\subsection{Poisson量子化の明示性}
\begin{proposition}\label{prop:poisson-quantization-explicit}
本論文のPoisson正規形
$F=\sum_{s=1}^k f_s(x_s,u)$に対し、
量子化を定める二次関係式は明示的に
\[
 [x_{s2},x_{s3}]
 =\hbar\,\operatorname{sym}_R(\partial_{x_{s1}}f_s),
 \quad
 [x_{s3},x_{s1}]
 =\hbar\,\operatorname{sym}_R(\partial_{x_{s2}}f_s),
 \quad
 [x_{s1},x_{s2}]
 =\hbar\,\operatorname{sym}_R(\partial_{x_{s3}}f_s)
\]
と書ける。異なるブロックと中心変数の間の交換子は零である。
三次中心元$\Omega_\hbar$は零のZariski近傍で存在し、
有限自由加群間の中心性方程式
$z\mapsto([z,x_{si}])_{s,i}$の核を解いて選べる。
しかし任意の入力に一様な閉公式が得られたとは主張しない。
\end{proposition}
\begin{proof}
二次関係式は各ブロックのcyclic potential
$W_s=x_{s1}x_{s2}x_{s3}-x_{s2}x_{s1}x_{s3}
-\hbar\,\operatorname{sym}_R(f_s)$のcyclic derivativeである。
中心元の存在と、核の計算を用いた代数化は
本論文の量子化節で証明する。
\end{proof}
\begin{example}[Fermat型]
$[x,y]=\hbar z^2$、$[y,z]=\hbar x^2$、
$[z,x]=\hbar y^2$なら
\[
 \Omega=x^3+y^3+z^3
      -2\hbar^2\operatorname{sym}(xyz)
\]
は実際に中心元となる。直接の交換子計算は
本論文のFermat型の補題にある。
\end{example}


\part{Poisson構造・退化スキーム・量子化}
\section{大域的双ベクトル場と射影空間への延長}\label{sec:pg-extension}
本節から\S\ref{sec:pg-special}までは $\CC$ 上で議論する。
$V$ を次元 $N=n+2\geq5$ のベクトル空間、$S=\operatorname{Sym}(V^*)$ とし、
$F\in S_3$ が滑らかな超曲面 $X=V(F)\subset\PP(V)=\PP^{N-1}$ を定めるとする。
$\PP(V)$ は $V$ の一次元部分空間をパラメータ化する。
座標を $z_1,\ldots,z_N$、Euler 場を $E=\sum_i z_i\partial_{z_i}$ と書き、
$F_i=\partial_iF$、$F_{ij}=\partial_i\partial_jF$ と略記する。
縮約の符号は
\[
\iota_\eta(v_1\wedge\cdots\wedge v_r)
 =\sum_{a=1}^r(-1)^{a-1}\eta(v_a)
 v_1\wedge\cdots\wedge\widehat v_a\wedge\cdots\wedge v_r
\]
とする。Schouten 括弧に関して $[\pi,\pi]=0$ を満たす
$\pi\in H^0(X,\bigwedge^2T_X)$ を Poisson 構造と呼ぶ。
射影空間の場 $\pi$、その制限 $\pi_X$、アフィン錐上の代表 $q$ を区別する。
射影スキームに関する斉次イデアルの等式は、不要イデアルで飽和したイデアル層の等式を意味する。
滑らかさから $F$ は既約である。実際、正次数の二因子があればその射影零点は交わり、交点は $X$ の特異点となる。

\subsection{Gradient の Koszul 複体}
\begin{lemma}\label{lem:pg-gradient}
$F_1,\ldots,F_N$ は $S$ の正則列である。
特に $\sum_i a_iF_i=0$ で、各 $a_i$ が定数または一次式ならば全ての $a_i$ は零である。
\end{lemma}
\begin{proof}
Euler 等式 $3F=\sum_i z_iF_i$ と射影的な滑らかさにより、gradient の共通零点はアフィン原点だけである。
従ってそのイデアルの高さは $N$ である。
$S$ は Cohen--Macaulay なので、$N$ 個の生成元は正則列となる。
第一 syzygy は $F_je_i-F_ie_j$ という Koszul 関係で生成され、その係数次数は二である。
よって係数次数零または一の syzygy は存在しない。
\end{proof}
\begin{lemma}\label{lem:pg-euler}
$\mathbb P=\PP(V)$ とすると
\[
H^0(\mathbb P,\textstyle\bigwedge^2 T_{\mathbb P}(-3))
=H^1(\mathbb P,\textstyle\bigwedge^2 T_{\mathbb P}(-3))=0,
\qquad
H^0(\mathbb P,\textstyle\bigwedge^3 T_{\mathbb P}(-3))\simeq\bigwedge^3V.
\]
\end{lemma}
\begin{proof}
$\dim\mathbb P\geq4$ における線束の中間コホモロジーの消滅\cite[Chapter III]{PHartshorne}を使う。
Euler 列とその外積は
\begin{align*}
0&\longrightarrow\OO(-3)\longrightarrow V\otimes\OO(-2)
 \longrightarrow T_{\mathbb P}(-3)\longrightarrow0,\\
0&\longrightarrow T_{\mathbb P}(-3)\longrightarrow\bigwedge^2V\otimes\OO(-1)
 \longrightarrow\bigwedge^2T_{\mathbb P}(-3)\longrightarrow0,\\
0&\longrightarrow\bigwedge^2T_{\mathbb P}(-3)\longrightarrow\bigwedge^3V\otimes\OO
 \longrightarrow\bigwedge^3T_{\mathbb P}(-3)\longrightarrow0.
\end{align*}
初めの二列が消滅を与え、最後の列の大域切断が同型を与える。
\end{proof}
\begin{proposition}\label{prop:pg-represent}
$F$ を固定すると線形同型
\[
\bigwedge^3V\xrightarrow{\sim} H^0(X,\bigwedge^2T_X),
\qquad P\longmapsto(\iota_{dF}P)|_X
\]
がある。右辺は二次係数の場を射影化して定める。
任意の双ベクトル場は $\PP(V)$ 上に一意に延長される。
\end{proposition}
\begin{proof}
補題\ref{lem:pg-euler}と制限列から
$H^0(\mathbb P,\bigwedge^2T_{\mathbb P})\simeq
H^0(X,(\bigwedge^2T_{\mathbb P})|_X)$ である。
従って延長は存在して一意であり、外積 Euler 列により
$B\in S_2\otimes\bigwedge^2V$ で代表できる。
接性により、$\iota_{dF}B$ の射影接束への像は $X$ 上零である。
Euler 列の核と $H^0(X,\OO_X(3))$ の多項式への持上げを用いると
\[
\iota_{dF}B=FA+hE,\qquad A\in S_1\otimes V,\quad h\in S_3.
\]
再び $dF$ で縮約すれば $0=FA(F)+3Fh$、従って $h=-A(F)/3$ である。
そこで $q=B-\frac13E\wedge A$ とおくと $\iota_{dF}q=0$ となる。
補題\ref{lem:pg-gradient}の Koszul 複体の次数部分
\[
S_0\otimes\bigwedge^3V\xrightarrow{\iota_{dF}}
S_2\otimes\bigwedge^2V\xrightarrow{\iota_{dF}}S_4\otimes V
\]
は完全で、左からの寄与は負次数なので零である。
よって $q=\iota_{dF}P$ を満たす定数 $P$ が一意に存在する。

代表の選択にも依存しない。
射影化が零である二次係数の場は $E\wedge Z$（$Z$ は線形）と書ける。
これが $F$ を Casimir にもつなら $3FZ=Z(F)E$ である。
既約性と各座標が $F$ と互いに素であることから $F\mid Z(F)$、
従って $Z(F)=cF$、$Z=(c/3)E$、$E\wedge Z=0$ となる。
逆に任意の $P$ について $\iota_{dF}^2P=0$ なので、射影化は $X$ に接する。
\end{proof}

\subsection{Jacobi 恒等式の延長}
\begin{proposition}\label{prop:pg-jacobi-extension}
$q=\iota_{dF}P$ の $X$ 上の射影化が Poisson ならば $[q,q]=0$ である。
従って射影空間上の一意な延長も Poisson である。
\end{proposition}
\begin{proof}
標準体積形式の divergence を
\[
\delta(a\partial_{i_1}\wedge\cdots\wedge\partial_{i_r})
=\sum_{b=1}^r(-1)^{b-1}\partial_{i_b}a\,
\partial_{i_1}\wedge\cdots\wedge\widehat{\partial_{i_b}}\wedge\cdots\wedge\partial_{i_r}
\]
と定める。交代テンソルと Hessian の縮約なので $\delta q=0$。
積の微分則から $\delta$ は Schouten 括弧の導分であり、
$J=[q,q]$ は $\delta J=0$ を満たす。
$F$ が Casimir であることから $\iota_{dF}J=0$ でもある。
$J$ の射影化は $X$ 上零なので、補題\ref{lem:pg-euler}より
\[
J=FT+E\wedge B_2,\qquad T\in\bigwedge^3V,
\quad B_2\in S_2\otimes\bigwedge^2V.
\]
定義から $\delta(FT)=\iota_{dF}T$、
$\delta(E\wedge B_2)=NB_2-E\wedge\delta B_2$。
$\delta J=0$ に $E$ を外積すると
$E\wedge B_2=-N^{-1}E\wedge\iota_{dF}T$ となる。
これを代入して $dF$ で縮約すると
\[
0=(1-3/N)F\,\iota_{dF}T.
\]
$N\geq5$ より $\iota_{dF}T=0$。
gradient の定数係数での独立性から $T=0$、従って $J=0$ である。
\end{proof}

\section{Hessian による Poisson 構造の正規形}\label{sec:pg-normal}
\subsection{Jacobi 条件の線形化}
$P\in\bigwedge^3V$、$\alpha,\beta\in V^*$ に対し、
$\gamma(P(\alpha,\beta))=P(\alpha,\beta,\gamma)$ により
$P(\alpha,\beta)\in V$ を定める。次の線形空間を導入する：
\begin{align}
K(P)=\{H\in\operatorname{Sym}^2V^*\mid{}
&H(P(\alpha,\beta),P(\gamma,\eta))
 +H(P(\beta,\gamma),P(\alpha,\eta))\notag\\
&+H(P(\gamma,\alpha),P(\beta,\eta))=0
\quad(\alpha,\beta,\gamma,\eta\in V^*)\}.\label{eq:pg-K}
\end{align}
\begin{lemma}\label{lem:pg-linearize}
$q=\iota_{dF}P$ が Poisson であることと、全 $p\in V$ について
$\operatorname{Hess}_F(p)\in K(P)$ であることは同値である。
また $F$ の滑らかさの仮定の下で、Hessian の線形系は $\PP(V)$ 上基点をもたない。
\end{lemma}
\begin{proof}
$P^{ija}=P(dz_i,dz_j,dz_a)$ とすれば $q^{ij}=\sum_aP^{ija}F_a$ である。
Jacobiator の係数は非零定数倍を除き
\[
\sum_aF_aL_a^{ijk},\qquad
L_a^{ijk}=\sum_{r,b}(P^{ira}P^{jkb}+P^{jra}P^{kib}+P^{kra}P^{ijb})F_{br}.
\]
$L_a^{ijk}$ は一次式なので、補題\ref{lem:pg-gradient}より
Jacobi 条件は全 $L_a^{ijk}=0$ と同値である。
これは\eqref{eq:pg-K}への座標共変ベクトルの代入と符号を除いて一致する。
多重線形性が第一の主張を与える。
三次微分の対称性により
$\operatorname{Hess}_F(p)(w,w)=2dF_w(p)$。
非零の $w$ が全 Hessian の零点なら $dF_w=0$、Euler 等式より $F(w)=0$ となり、
$[w]$ が $X$ の特異点となるので矛盾する。
\end{proof}

\subsection{基点をもたない二次形式と三重ベクトル}
\begin{lemma}\label{lem:pg-quadratic-lie}
$K(P)$ が $\PP(V)$ 上基点をもたなければ、直和をなす三次元部分空間
$W_1,\ldots,W_k\subset V$ と $0\neq P_s\in\bigwedge^3W_s$ が存在して
$P=\sum_sP_s$ と書ける。$P=0$ では $k=0$ とする。
\end{lemma}
\begin{proof}
$P\neq0$ とし、その support を
$W=\operatorname{span}\{P(\alpha,\beta)\mid\alpha,\beta\in V^*\}$ とおく。
これは $P\in\bigwedge^3W$ を満たす最小の部分空間である。
$K(P)$ の $\PP(W)$ への制限も基点をもたない。
標数零の Bertini 定理により、それぞれの一般の二次超曲面は滑らかであり、
従って $V$ 上および $W$ 上での非退化条件は非空開条件である。
両者を満たす $B\in K(P)$ を選ぶ。

$B^\flat:V\to V^*$ により $[x,y]=P(B^\flat x,B^\flat y)$ と定める。
$B([x,y],z)=P(B^\flat x,B^\flat y,B^\flat z)$ は交代的なので $B$ は不変である。
Jacobiator を $B$ で任意の $t$ と対にすると、\eqref{eq:pg-K}の $H=B$ が得られる。
よってこれは Lie 代数 $\mathfrak g$ である。
$[\mathfrak g,\mathfrak g]=W$、$Z(\mathfrak g)=W^\perp$ であり、$B|_W$ の非退化性から
\[
\mathfrak g=W\oplus Z(\mathfrak g),\qquad [W,W]=W,\qquad Z(W)=0.
\]
さらに任意の $H\in K(P)$ について
\begin{equation}
H([a,b],[c,d])+H([b,c],[a,d])+H([c,a],[b,d])=0.\label{eq:pg-lie-H}
\end{equation}

$\mathfrak r=\operatorname{rad}(W)\neq0$ と仮定する。
その導来列の最後の非零項から極小の非零 $W$-イデアル $I$ を取る。
$I$ は可換で、既約な $W$-加群である。
複素有限次元 Lie 代数の既約表現では可解根基はスカラーで作用する
\cite[Proposition 16.9]{PEtingofLie}。
$W=[W,W]$ のため全作用素の trace は零であり、$\mathfrak r$ は $I$ に零で作用する。
従って $I$ は $W/\mathfrak r$ の既約加群であり、$Z(W)=0$ から自明表現ではない。
最高ウェイトベクトル $0\neq v\in I$ と単純根 $\alpha$ を
$\lambda(h_\alpha)>0$ となるように選び、
$e_\alpha,f_\alpha,h_\alpha$ を $W$ に持ち上げる。
$w=\lambda(h_\alpha)^{-1}[f_\alpha,v]\in I$ とおくと
\[
[e_\alpha,w]=v,\quad [e_\alpha,v]=0,\quad
[h_\alpha,v]=\lambda(h_\alpha)v,\quad[v,w]=0.
\]
商の $\mathfrak{sl}_2$ 関係が作用するのは、根基の作用が零だからである。
\eqref{eq:pg-lie-H}に $(a,b,c,d)=(h_\alpha,v,e_\alpha,w)$ を代入すると
$\lambda(h_\alpha)H(v,v)=0$ が全 $H$ に対して成立する。
これは基点自由性に反する。従って $W$ は半単純である。

階数二以上の単純因子があるとする。
その最高根 $\theta$ は二つの正根 $\alpha,\beta$ の和に書ける。
$a=e_\alpha$、$b=e_\beta$、$v=[a,b]\neq0$ とすれば
$[a,v]=[b,v]=0$ である。
$\theta(h)\neq0$ となる Cartan 元を選び、
\eqref{eq:pg-lie-H}に $(a,b,h,v)$ を代入すると
$\theta(h)H(v,v)=0$ となり、再び基点を得て矛盾する。
従って $W\simeq\mathfrak{sl}_2^{\oplus k}$ である。
用いた根分解・最高ウェイトの事実は
\cite[\S19.4, Theorem 25.17, \S32.2]{PEtingofLie}を参照。

異なる単純イデアルは $B$ に関して直交する。
実際、$x=[a,b]$ が一方に、$y$ が他方に属せば
$B(x,y)=B(a,[b,y])=0$ である。中心も $W$ に直交する。
よって Cartan 三形式 $B([x,y],z)$ は各 $\mathfrak{sl}_2$ 因子上の三形式の和であり、
$B$ で添字を上げれば所要の $P=\sum_sP_s$ を得る。
\end{proof}

\subsection{三次式の分解と一般階数}
\begin{theorem}[Poisson 正規形]\label{thm:pg-normal}\label{pthm:normal-form}
任意の非零 Poisson 構造に対して、$N=3k+m$ となる $k\geq1$ と線形座標
$x_s=(x_{s1},x_{s2},x_{s3})$、$u=(u_1,\ldots,u_m)$ が存在し、
\begin{equation}
P=\sum_{s=1}^k\partial_{x_{s1}}\wedge\partial_{x_{s2}}\wedge\partial_{x_{s3}},
\qquad F=\sum_{s=1}^k f_s(x_s,u)\label{eq:pg-normal}
\end{equation}
となる。括弧は
\begin{equation}
\{x_{s1},x_{s2}\}=(f_s)_{x_{s3}},\quad
\{x_{s2},x_{s3}\}=(f_s)_{x_{s1}},\quad
\{x_{s3},x_{s1}\}=(f_s)_{x_{s2}}\label{eq:pg-bracket}
\end{equation}
と、異なるブロック間および $\{u_\alpha,-\}$ が零であることで定まる。
逆に\eqref{eq:pg-normal}を満たす滑らかな $F$ はこの括弧によって Poisson 構造を定める。
$\pi$ と $\pi_X$ の一般階数はともに $2k$ である。
\end{theorem}
\begin{proof}
補題\ref{lem:pg-linearize}、\ref{lem:pg-quadratic-lie}を適用し、各三次元因子の基底の行列式を調整する。
異なる因子からそれぞれ $a,b$ と $c,d$ を取ると、
\eqref{eq:pg-lie-H}は $H([a,b],[c,d])=0$ となる。
$[W_s,W_s]=W_s$ より、全 Hessian の $W_s\times W_t$ 成分は $s\neq t$ のとき零である。
従って $\partial_{x_{si}}\partial_{x_{tj}}F=0$。
$F$ の各単項式は高々一つの active ブロックを含み、中心変数のみの項をいずれかに配分すれば所要の分解となる。
逆向きは三変数 Jacobian 括弧の Jacobi 恒等式
$\nabla f\cdot(\nabla\times\nabla f)=0$ と、異なるブロックが互いの係数を微分しないことから従う。

各 active gradient は恒等的に零ではない。
そうなら $F$ は対応する三変数に依存せず、その座標平面に特異点をもつ。
従ってアフィン空間の一般階数は $2k$。
非零な二次偏微分は既約な三次式 $F$ で割り切れないので、$X$ の一般点でも全ブロックが階数二をもつ。
階数 $2r$ の交代双ベクトルを直線 $L$ で割ると、
$L$ がその像に入らなければ階数は $2r$、入れば $2r-2$ である。
これは像上の Darboux 基底により分かる。
ここで $L=\CC E$。
$m>0$ なら一般点で $E$ は中心方向成分をもち、像に入らない。
$m=0$ なら第 $s$ ブロックの像は $\ker df_s$ で、Euler 方向が入る条件は $3f_s=0$。
$N\geq5$ の下で $m=0$ なら $k\geq2$ であり、全 $f_s=0$ は $X$ の真の閉集合である。
よって一般点では射影化で階数は低下しない。
\end{proof}
座標、ブロックの順序、中心三次式の配分の一意性は主張しない。
また $F$ を $aF$ に替えると命題\ref{prop:pg-represent}の同型は $a$ 倍される。

\begin{corollary}\label{cor:pg-generic}
固定した滑らかな $X$ 上の一般階数二の Poisson 構造の射影化は
$\operatorname{Gr}(3,V)\subset\PP(\bigwedge^3V)$ である。
一般の滑らかな三次超曲面上では、全ての非零 Poisson 構造が一般階数二をもつ。
\end{corollary}
\begin{proof}
$k=1$ と非零 $P$ の分解可能性が同値なので第一の主張を得る。
第二の主張には $k\geq2$ の族の次元を評価する。
直和をなす順序付き $W_1,\ldots,W_k$ の空間は
$\operatorname{Gr}(3,V)^k$ の開部分で、次元 $3k(N-3)$。
この組を固定したときの三次式の空間の次元は
\[
\binom{m+2}{3}+k\left(10+6m+3\binom{m+1}{2}\right).
\]
従って全三次式の空間における余次元は少なくとも
\begin{equation}
\binom{N+2}{3}-\binom{m+2}{3}
-k\left(10+6m+3\binom{m+1}{2}\right)-3k(N-3)
=\frac{3k}{2}\bigl((3k-5)m+3k(k-1)\bigr)>0.\label{eq:pg-codim}
\end{equation}
可能な $k$ は有限個なので一般の三次式はこれらを避ける。
\end{proof}
\begin{remark}
\eqref{eq:pg-codim}は余次元の下界である。
$(n,k)=(6,2),(7,2),(7,3)$ ではそれぞれ $24,27,81$ となる。
正確な余次元や正規形の一意性を結論するものではない。
三次元の場合の Grassmann 記述は\cite[Theorem 1]{PMayanskiy}の既知の結果である。
\end{remark}

\section{Casimir、対数形式および symplectic leaf}\label{sec:pg-leaves}
以下、\eqref{eq:pg-normal}を固定する。
多項式 $G$ が Casimir であるとは $\{G,S\}=0$ を意味する。
射影超曲面が Poisson であるとは、そのイデアル層が括弧で保たれることをいう。

\subsection{Poisson 超曲面}
\begin{proposition}\label{prop:pg-casimir}
既約な斉次式 $G\in S$ に対し、$V(G)\subset\PP(V)$ が Poisson であることと、
$G$ がアフィン錐上の Casimir であることは同値である。
一次 Casimir の空間は $\operatorname{span}(u_1,\ldots,u_m)$ である。
従って Poisson 超平面のパラメータ空間は $m>0$ なら $\PP^{m-1}$、$m=0$ なら空である。
\end{proposition}
\begin{proof}
Casimir なら $(G)$ は括弧で保たれる。
逆向きは $G$ が $F$ の定数倍の場合は既知なので、$\gcd(F,G)=1$ とする。
$d=\deg G$ とおくと、射影的接性と Euler 列より
\[
\iota_{dG}q=GA+hE,\qquad A\in S_1\otimes V,\quad h\in S_d.
\]
$dG$ で縮約すれば $h=-A(G)/d$。
$dF$ で縮約し $\iota_{dF}q=0$ を使うと
$GA(F)-(3/d)FA(G)=0$ である。
互いに素であることと次数から $A(F)=cF$、$A(G)=cdG/3$。
$A-(c/3)E$ は gradient の一次 syzygy なので零である。
従って $\iota_{dG}q=0$。
一次式 $\ell$ については gradient の定数係数での独立性から
$\iota_{d\ell}q=0\iff\iota_{d\ell}P=0$ であり、正規形では
$\ell$ が全 active 変数に依存しないことと同値である。
\end{proof}
各 $u_\alpha,f_s,F$ とその多項式結合は Casimir であるが、
多項式中心全体がそれらで生成されるとは主張しない。

\subsection{最高外積の対数表示}
\begin{proposition}\label{prop:pg-log}
体積形式を固定すると、$\pi^k$ に対応する微分形式は非零定数倍を除き
\begin{equation}
\beta=\iota_E(du_1\wedge\cdots\wedge du_m\wedge df_1\wedge\cdots\wedge df_k)
\in H^0(\PP^{N-1},\Omega^{m+k-1}(N)).\label{eq:pg-log-main}
\end{equation}
$m>0$ では
\[
\beta=\left(\prod_{\alpha=1}^m u_\alpha\prod_{s=1}^kf_s\right)
\bigwedge_{\alpha=2}^m d\log(u_\alpha/u_1)
\wedge\bigwedge_{s=1}^k d\log(f_s/u_1^3),
\]
$m=0$ では
\[
\beta=3\left(\prod_{s=1}^kf_s\right)
\bigwedge_{s=2}^k d\log(f_s/f_1).
\]
$X$ 上の $\pi_X^k$ に対応する形式は
\begin{equation}
\alpha=\iota_E(du_1\wedge\cdots\wedge du_m\wedge df_1\wedge\cdots\wedge df_{k-1})|_X
\in H^0(X,\Omega_X^{m+k-2}(N-3)).\label{eq:pg-log-X}
\end{equation}
\end{proposition}
\begin{proof}
第 $s$ ブロックの場を $q_s$ と書く。
$q_s^2=0$ より $q^k=k!q_1\wedge\cdots\wedge q_k$。
アフィン体積形式への縮約は、非零定数倍を除き
$du_1\wedge\cdots\wedge du_m\wedge df_1\wedge\cdots\wedge df_k$ である。
$df_s$ の中心方向成分は先頭の外積で消える。
さらに $E$ で縮約すれば射影化に対応する水平形式となり、
全スケーリングの重み $m+3k=N$ が twist を与える。
$m>0$ では $\iota_Ed\log u_1=1$、各比の対数微分の縮約は零であり、外積を展開して所要の式を得る。
$m=0$ では $\iota_Ed\log f_1=3$ を使う。
計算は分母が消えない開集合で行うが、左辺が大域的正則なので大域表示として解釈できる。

最後に $f_k$ を $F=\sum_sf_s$ に替えても微分の外積は変わらない。
$X$ 上で法線方向 $dF$ を除き adjunction を適用すると\eqref{eq:pg-log-X}を得る。
局所的に $dF$ が座標微分となる表示で縮約を行えばよく、法線束 $\OO_X(3)$ の遷移関数が twist を三減らす。
形式次数は $m+k-2=n-2k$ となる。
\end{proof}
対数表示の因子の次数の総和は $m+3k=N$ であり、
$V(u_1\cdots u_mf_1\cdots f_{k-1}F)$ を反標準因子にできる。
これは最高外積の表示であって、$k>1$ では $\pi^k$ だけから $\pi$ は復元できず、
この因子の全体が退化軌跡であるとも限らない。

\subsection{葉とその閉包}
\begin{proposition}\label{prop:pg-first-integrals}
有理写像
\[
\Phi:\PP^{N-1}\dashrightarrow\PP(1^m,3^k),\qquad
[x,u]\longmapsto[u_1:\cdots:u_m:f_1:\cdots:f_k]
\]
は一般点で独立な第一積分を与える。
一般 fiber の滑らかな最大階数部分の連結成分は symplectic leaf である。
$X$ 上の底は $m>0$ なら $\PP(1^m,3^{k-1})$、$m=0$ なら $\PP^{k-2}$ であり、次元は $n-2k$ である。
\end{proposition}
\begin{proof}
全 active gradient が非零の開集合では、$du_\alpha$ を法とした $df_s$ は独立なブロックに非零成分をもつ。
従って $(u,f_1,\ldots,f_k)$ の微分の階数は $m+k$。
$F=0$ への制限と射影化は、それぞれ独立な条件 $\sum_sf_s=0$ とスケーリングを一つ除く。
Casimir 性から Hamilton 分布は $\ker d\Phi$ に含まれ、
定理\ref{thm:pg-normal}によって両者の次元は $2k$ なので一致する。
最大階数部分での逆二形式は Jacobi 恒等式により閉じている。
よって滑らかな fiber の連結成分が葉となる。
不定点では $u=0$ かつ全 $f_s=0$。
active gradient が一つでも零ならアフィン階数は高々 $2k-2$、
全て非零なら Euler 方向が像に入るため射影階数は $2k-2$ である。
よって最大階数の葉はこの不定点集合を横切らない。
\end{proof}
\begin{proposition}\label{prop:pg-leaf-closure}
各 $f_{s,0}=f_s(x_s,0)$ が滑らかな平面三次曲線を定めるとする。
$m>0$ なら一般の葉の閉包は $\PP^{3k}$ 内の型 $(3^k)$ の既約な完全交叉であり、
$m=0$ なら $\PP^{3k-1}$ 内の型 $(3^{k-1})$ の既約な完全交叉である。
いずれも次元は $2k$ である。
\end{proposition}
\begin{proof}
$m>0$ では底の一般点で $u_\alpha=c_\alpha u_1$、$c_1=1$ とし、$t=u_1$ とおく。
候補は $f_s(x_s,ct)-d_st^3=0$（$1\leq s\leq k$）であり、$X$ 上では $\sum_sd_s=0$。
$t=0$ の部分は独立な式 $f_{s,0}=0$ で定まり次元 $2k-1$。
$t\neq0$ では $t=1$ と正規化すると $k$ 個のアフィン曲面の積となる。
従って候補は純次元 $2k$ の完全交叉で、無限遠に余分な成分をもたない。
標数零の generic smoothness により、一般の底点上のアフィン積は滑らかである。
各因子の射影閉包は無限遠でも $\nabla f_{s,0}\neq0$ により滑らかであり、
連結な射影超曲面なので既約である。
よってアフィン fiber は既約。
完全交叉は Cohen--Macaulay で埋め込み成分をもたず、全最小成分がこの滑らかな開集合と交わる。
従って候補は被約かつ既約で、葉の閉包である。

$m=0$ では $f_1\neq0$ の一般底点に対して $f_s=c_sf_1$（$s\geq2$）と書ける。
$f_1=0$ の部分は全 $f_s=0$ で定まり次元 $2k-1$。
開部分は $\{f_1=1\}\times\prod_{s=2}^k\{f_s=c_s\}$ の同時 $\mu_3$ 作用による像である。
$c_s\neq0$ の一般点では、各因子は滑らかな三次曲面
$f_s(x_s)=c_st^3$ のアフィン開集合なので既約である。
全 active gradient が非零なので方程式の微分は独立であり、fiber は滑らか。
前段と同じ純次元性・Cohen--Macaulay 性により結論を得る。
\end{proof}
\begin{remark}\label{rem:pg-zero-restriction-leaf}
平面三次制限の滑らかさを仮定しなくても命題\ref{prop:pg-first-integrals}は成立する。
ただし上の方程式から閉包を得るには、$m>0$ では $t$ に関して飽和し、必要なら fiber の既約成分を選ぶ。
例えば $k=1$、$f_{1,0}=0$ では $F(x,ct)=tQ_c(x,t)$ となる。
$Q_\alpha(x)=F_{u_\alpha}(x,0)$ は $\PP_x^2$ 上共通零点をもたない。
そうした点は $X$ の特異点となるからである。
Bertini により一般の $\sum_\alpha c_\alpha Q_\alpha$ は滑らかな平面二次曲線を定める。
従って $Q_c(x,t)$ の階数は少なくとも三で、$t$ で割り切れず既約である。
葉の閉包は固定平面 $t=0$ を除いた residual quadric $Q_c=0$ となる。
\end{remark}

\section{全退化スキームの方程式}\label{sec:pg-degeneracy}
\subsection{任意のブロック表示に対する生成元}
$q_s=\iota_{df_s}(\partial_{x_{s1}}\wedge\partial_{x_{s2}}\wedge\partial_{x_{s3}})$ とし、
\[
A_{s,i}=\frac{\partial f_s}{\partial x_{si}},\qquad
p_s=\sum_i x_{si}A_{s,i}=3f_s-\sum_\alpha u_\alpha\frac{\partial f_s}{\partial u_\alpha}
\]
とおく。$D_{2j-2}(\pi)=Z(\pi^j)$ をスキームとして考える。
\begin{proposition}\label{prop:pg-deg-ideal}
$1\leq j\leq k$ に対し、$D_{2j-2}(\pi)$ のイデアル層は、全ての
$I\subset\{1,\ldots,k\}$、$|I|=j$ について
\begin{equation}
 u_\alpha\prod_{s\in I}A_{s,i_s},\qquad
 p_t\prod_{s\in I\setminus\{t\}}A_{s,i_s}\ (t\in I),\qquad
 x_{t,l}\prod_{s\in I}A_{s,i_s}\ (t\notin I)\label{eq:pg-deg-generators}
\end{equation}
を生成元とする。各添字は可能な範囲を動く。
$X$ 上のイデアルはさらに $F$ を加えて得られる。
各生成元の次数は $2j+1$ である。
\end{proposition}
\begin{proof}
錐の原点以外では $E$ の生成する直線を局所的な直和因子に選べる。
その商束の外積から $E\wedge-$ は分裂単射となるので、
射影化の零点イデアルは $E\wedge q^j$ の係数イデアルに等しい。
$q_s^2=0$ より $q^j=j!\sum_{|I|=j}\bigwedge_{s\in I}q_s$。
$E$ の中心成分は第一の型を与える。
$t\in I$ では $(\sum_lx_{tl}\partial_{x_{tl}})\wedge q_t
=p_t\partial_{x_{t1}}\wedge\partial_{x_{t2}}\wedge\partial_{x_{t3}}$ なので第二の型となる。
$t\notin I$ は第三の型を与える。
ブロック別外積次数は、中心に一、選択したブロックに三、選択しないブロックに一、という異なる型をもつ。
残りの選択ブロックは二であるため、異なる型の間に相殺はない。
最後に $T_X\subset T_{\PP(V)}|_X$ は局所的に分裂するので、外積への包含は切断の零点イデアルを変えない。
\end{proof}

\subsection{有限平坦 critical cover}
本小節以降、特に断らない限り
\begin{equation}
f_{s,0}=f_s(x_s,0)\text{ が滑らかな平面三次曲線 }E_s\text{ を定める}\label{eq:pg-smooth-block}
\end{equation}
と仮定する。$E_s$ は $u=0$、$x_t=0$（$t\neq s$）により $X$ 内に入り、互いに交わらない。
\begin{lemma}\label{lem:pg-critical-cover}
$m>0$ なら
\[
K_s=V(A_{s,1},A_{s,2},A_{s,3})\subset\PP(x_s,u)=\PP^{m+2},
\qquad\Delta=V(A_{s,i}\mid s,i)
\]
は $u=0$ と交わらない。
中心座標への射影は $K_s\to\PP_u^{m-1}$ を有限平坦八次被覆、
$\Delta\to\PP_u^{m-1}$ を有限平坦 $8^k$ 次被覆とし、
$\Delta\simeq K_1\times_{\PP_u^{m-1}}\cdots\times_{\PP_u^{m-1}}K_k$。
$m\geq0$ の場合を含め、$\Gamma_I=V(A_{s,i}\mid s\in I,1\leq i\leq3)$ は
非空なら型 $(2^{3|I|})$ の完全交叉で、
\[
\dim\Gamma_I=N-1-3|I|,\qquad\deg\Gamma_I=8^{|I|}.
\]
特に $m>0$ では
$\dim\Delta=m-1$、$\deg\Delta=8^k$、$\omega_\Delta\simeq\OO_\Delta(3k-m)$。
$m=0$ では $\Delta=\varnothing$ とする。
\end{lemma}
\begin{proof}
$R=\CC[u_1,\ldots,u_m]$、$Q_s=R[x_s]/(A_{s,1},A_{s,2},A_{s,3})$ とする。
$Q_s/(u)$ は滑らかな三次式の Jacobian 環で、Hilbert 級数 $(1+t)^3$ をもつ。
その八個の斉次基底を持ち上げれば、次数による帰納法から $Q_s$ は有限 $R$-加群となる。
三方程式による次元下界と有限性より $\dim Q_s=m$。
従って方程式は正則列で $Q_s$ は Cohen--Macaulay である。
$u_1,\ldots,u_m$ は斉次パラメータ系なので $Q_s$-正則列である。
持上げからの全射
\[
R\oplus R(-1)^{\oplus3}\oplus R(-2)^{\oplus3}\oplus R(-3)\longrightarrow Q_s
\]
は同型である。実際、$u$ の Koszul 完全性により
$\operatorname{Tor}_1^R(Q_s,R/(u))=0$ であり、全射の核を $(u)$ で割ると零となる。
次数付き Nakayama が核を零にする。
複数ブロックの環は $R$ 上のテンソル積なので階数 $8^{|I|}$ の有限自由加群である。
残りの active 変数を加えれば、$\Gamma_I$ の高さと完全交叉性が従う。
$m=0$ かつ $I$ が全体のときだけ射影スキームは空である。

$u=0$ では各 gradient の共通零点は原点だけなので、$K_s$ と $\Delta$ は射影的にその部分と交わらない。
$u_\alpha\neq0$ で全ブロックを同時に $u_\alpha=1$ と正規化すれば、有限自由性が被覆の有限平坦性と fiber product 表示を与える。
重なりでも同時正規化するため貼り合わさる。次数と双対化層は Hilbert 級数および随伴公式から従う。
\end{proof}

\subsection{退化集合と零点スキーム}
空でない $L\subset\{1,\ldots,k\}$ に対し
\[
J_L=V(u,x_t\ (t\notin L),f_{s,0}\ (s\in L))
\]
を $E_s$（$s\in L$）の射影的 join とする。
独立な錐の積の射影化なので、被約既約で次元 $2|L|-1$ である。
\begin{proposition}\label{prop:pg-deg-support}
$0\leq r<k$ について
\[
|D_{2r}(\pi)|=\bigcup_{|I|=k-r}|\Gamma_I|\ \cup\ \bigcup_{|L|=r+1}J_L.
\]
$X$ 上では右辺を $X$ と交わらせる。
零点についてはスキームとして
\[
D_0(\pi)=\left(\coprod_sE_s\right)\sqcup\Delta,
\qquad D_0(\pi_X)=\left(\coprod_sE_s\right)\sqcup Z,
\qquad Z=\Delta\cap X.
\]
\end{proposition}
\begin{proof}
非零 active gradient をもつブロック数を $\ell$ とするとアフィン階数は $2\ell$。
$u\neq0$ では Euler 方向が像に入らず、射影化で階数は変わらない。
$u=0$ では\eqref{eq:pg-smooth-block}により gradient が零であることと $x_s=0$ が同値である。
非零ブロックの像は $\ker df_{s,0}$ なので、Euler 方向が像に入る条件は各 $3f_{s,0}=0$。
この場合に限り射影階数は $2\ell-2$ となる。これが台の表示を与える。

スキーム構造には命題\ref{prop:pg-deg-ideal}を用いる。
$u_\alpha\neq0$ では $u_\alpha A_{s,i}$ が全 active 偏微分を生成するので零点は $\Delta$ である。
$E_s$ の点ではある $A_{s,i}$ が単元となり、生成元は局所的に $u,x_t$（$t\neq s$）および $p_s$ を生成する。
これらを法として $p_s=3f_{s,0}$ で、他の生成元は消える。
従って零点イデアルは厳密に $(u,x_t\ (t\neq s),f_{s,0})$。
各 $E_s$ は被約であり、$\Delta$ と交わらないので所要の分解となる。
\end{proof}
\begin{remark}
台の表示を高次の退化スキームの被約分解と取り違えてはならない。
例えば $m=0$ の $\Gamma_I$ は座標線形空間を台とするが、Jacobian イデアルは一般に非被約な厚みを定める。
\end{remark}

\subsection{完全交叉性と一般の滑らかさ}
\begin{proposition}\label{prop:pg-zero-ci}
$m=1$ なら $Z=\varnothing$。
$m\geq2$ なら一般性を仮定せず $Z$ は型 $(2^{3k},3)$ の完全交叉で、
\[
\dim Z=m-2,\qquad \deg Z=3\cdot8^k,\qquad
\omega_Z\simeq\OO_Z(3k+3-m).
\]
滑らかな $f_{s,0}$ を固定して他の係数を一般に選ぶと、$K_s,\Delta$ および $m\geq2$ の場合の $Z$ は滑らかである。
これらが正次元ならば、滑らかな場合には既約である。
\end{proposition}
\begin{proof}
$\Delta$ の斉次環を $Q$ とする。
$Q$ は $R=\CC[u]$ 上有限自由な Cohen--Macaulay 環で、各最小素イデアルは $R$ の零イデアルに縮約する。
実際、各成分は純次元性により次元 $m$ であり、有限性から像も次元 $m$ である。
ある成分上 $F=0$ とし、その関数体を $L$ とする。
$L/\CC(u)$ は有限分離拡大なので、$du_1,\ldots,du_m$ は $\Omega^1_{L/\CC}$ の基底である。
全 active 偏微分が零なので
$0=dF=\sum_\alpha F_{u_\alpha}du_\alpha$、従って全中心偏微分も零となる。
この成分は $u$-空間を支配し原点以外の点をもつため、$X$ の滑らかさに矛盾する。
よって $F$ はどの最小素点にも入らず、Cohen--Macaulay 性により非零因子である。
これが $m\geq2$ の完全交叉性と数値を与える。
$m=1$ では交点があれば $u_1\neq0$ と $3F=u_1F_{u_1}$ により全偏微分が零となり矛盾する。

一般性を確認する。$u$ を含む係数を動かすパラメータ空間では、
$u_\alpha\neq0$ の点で項 $x_{si}u_\alpha^2$ の係数が $A_{s,i}$ の値を独立に変える。
従って全 critical incidence は滑らかである。
$u=0$ には解がなく、incidence はパラメータ空間と射影空間の積の閉部分なので、その射影は proper である。
全 fiber は非空な有限被覆なので射影は支配的であり、標数零の generic smoothness から一般の $K_s,\Delta$ は滑らかとなる。
そのような $\Delta$ を固定し $F$ に中心三次式 $c(u)$ を加えても active 偏微分は変わらない。
切断 $c(u)|_\Delta$ は基点をもたず、Bertini から一般の $Z$ は滑らかとなる。
$X$ の滑らかさも非空開条件である。
実際 $\sum_sf_{s,0}+\sum_\alpha u_\alpha^3$ が滑らかな例を与える。
従ってこれらの開条件を同時に課せる。
正次元の射影完全交叉の Koszul resolution と線束の中間コホモロジー消滅から $H^0(\OO)=\CC$。
よって連結であり、滑らかならば既約である。
\end{proof}

\section{六次元および七次元の幾何}\label{sec:pg-six-seven}
\subsection{可能な階数}
定理\ref{thm:pg-normal}により次の表を得る。
\begin{center}
\begin{tabular}{ccccc}\toprule
$n$ & $k$ & 一般階数 & $m$ & Poisson 超平面\\\midrule
6&1&2&5&$\PP^4$\\
6&2&4&2&$\PP^1$\\
7&1&2&6&$\PP^5$\\
7&2&4&3&$\PP^2$\\
7&3&6&0&なし\\\bottomrule
\end{tabular}
\end{center}
全双ベクトル場の空間の次元は $\binom83=56$、$\binom93=84$。
階数二の構造の射影化は $\operatorname{Gr}(3,8)$、$\operatorname{Gr}(3,9)$ で、次元は十五、十八である。

\subsection{階数二の零点：一般型および Calabi--Yau 完全交叉}
$k=1$ とし、active 座標を $x,y,z$ と書く。$m=n-1$ である。
$C=X\cap\{u=0\}$、$\Gamma=V(F_x,F_y,F_z)$、$Y=\Gamma\cap X$ とする。
\begin{proposition}\label{prop:pg-rank2}
固定した任意の滑らかな $X$ に対し、一般の三次元部分空間を active ブロックに選ぶと、
$C,\Gamma,Y$ は全て滑らかで、スキームとして
$D_0(\pi)=C\sqcup\Gamma$、$D_0(\pi_X)=C\sqcup Y$ である。
さらに
\[
\Gamma:(2,2,2),\quad\deg\Gamma=8,\quad K_\Gamma=(4-n)H,
\qquad
Y:(2,2,2,3),\quad\deg Y=24,\quad K_Y=(7-n)H.
\]
\end{proposition}
\begin{proof}
全偏微分の張る二次式の線形系は射影空間上基点をもたない。
Bertini を三回適用すると一般の三線形結合の零点 $\Gamma$ は滑らかとなる。
同じ線形系を滑らかな $X$ に制限すれば $Y$ も滑らかとなる。
一般の平面切断 $C$ の滑らかさも Bertini による。
これらは active 部分空間の Grassmann 多様体上の非空開条件なので同時に満たせる。
命題\ref{prop:pg-deg-support}、\ref{prop:pg-zero-ci}と随伴公式を適用する。
\end{proof}
$n=6$ では $\Gamma$ は $-K_\Gamma=2H$ の次数八の Fano 四次元多様体、
$Y$ は $K_Y=H$、$K_Y^3=24$ の一般型三次元多様体となる。
$n=7$ では $\Gamma$ は $-K_\Gamma=3H$ の Fano 五次元多様体、
$Y$ は次数二十四の Calabi--Yau 四次元多様体である。
ここで Calabi--Yau とは $K_Y\simeq\OO_Y$ および $H^i(Y,\OO_Y)=0$（$0<i<4$）を意味し、
Lefschetz 定理から単連結性も従う。
$h^{p,q}=\dim H^q(\Omega^p)$、$p_g=h^0(\omega)$、$e$ を位相的 Euler 数と書くと
\begin{align*}
n=6:&\quad p_g(Y)=8,\quad h^{1,1}(Y)=1,\quad h^{2,1}(Y)=246,\quad e(Y)=-504,\\
n=7:&\quad h^{1,1}(Y)=1,\quad h^{2,1}(Y)=0,\quad h^{3,1}(Y)=183,\quad h^{2,2}(Y)=780,\quad e(Y)=1152.
\end{align*}
計算は\S\ref{sec:pg-numerical}で与える。

\subsection{階数四と中心変数}
$F=f(a,u)+g(b,u)$、$a,b\in\CC^3$ とし、
$A_i=f_{a_i}$、$B_j=g_{b_j}$、$\rho_f=\sum_i a_iA_i$、$\rho_g=\sum_jb_jB_j$ とおく。
命題\ref{prop:pg-deg-ideal}は
\begin{align*}
I_{D_2(\pi)}&=(u_\alpha A_iB_j,\rho_gA_i,\rho_fB_j),\\
I_{D_0(\pi)}&=(u_\alpha A_i,u_\alpha B_j,\rho_f,\rho_g,b_jA_i,a_iB_j)
\end{align*}
を与える。
$f_0,g_0$ が滑らかな平面三次曲線 $E_f,E_g$ を定めるとし、
$\Gamma_f=V(A_1,A_2,A_3)$、$\Gamma_g=V(B_1,B_2,B_3)$、
$J=V(u,f_0,g_0)$ とおく。
$J$ は二曲線の三次元 join であり、
\[
|D_2(\pi)|=|\Gamma_f|\cup|\Gamma_g|\cup J,
\qquad |D_2(\pi_X)|=|\Sigma_f|\cup|\Sigma_g|\cup J,
\quad\Sigma_f=\Gamma_f\cap X,
\]
および同様の $\Sigma_g$ を得る。
\begin{proposition}\label{prop:pg-cubic-surface-family}
$\Sigma_f,\Sigma_g$ は型 $(2,2,2,3)$、次元 $m+1$、次数二十四の完全交叉であり、
一般のデータでは既約である。
$K_f=V(A_1,A_2,A_3)\subset\PP(a,u)$ および同様の $K_g$ を底として、
それぞれの双有理モデルは平坦な三次曲面族をなす。
相対 Jacobian イデアルで定まる fiber の特異スキームは $Z=\Delta\cap X$ である。
一般 fiber は滑らかだが全 fiber が滑らかではない。
\end{proposition}
\begin{proof}
$\Gamma_f$ は $K_f$ と反対側の $\PP_b^2$ の join であり、その斉次環は
$\CC[a,u]/(A)$ に変数 $b$ を加えた環である。
$F$ の最高 $b$ 次数部分 $g_0(b)$ は非零で係数は $\CC$ に属する。
従って随伴次数環での $g_0$ の乗法は単射であり、元の環で $F$ は非零因子となる。
補題\ref{lem:pg-critical-cover}から完全交叉性と数値を得る。

一般のデータを取り $K_f$ を滑らかとする。
join の頂点の解消として、直線をパラメータ化する
$\mathbb P_f=\PP_{K_f}(\OO_{K_f}(-1)\oplus\OO_{K_f}^{\oplus3})\to\Gamma_f$ を取る。
底の代表 $(a,u)$ を使う表示では $[t:b]\mapsto[ta:tu:b]$ であり、strict transform は
\begin{equation}
t^3f(a,u)+g(b,tu)=0.\label{eq:pg-surface-fiber}
\end{equation}
代表変更に伴い $t$ が逆に変わるので式は貼り合わさる。
$t=0$ に全成分は含まれず、余分な例外成分はない。
各 fiber の式は $t=0$ で非零の $g_0(b)$ となり、常に三次曲面を定める。
相対超曲面の完全列は fiber への制限後も左端が単射なので、族は平坦である。
$t=0$ では $g_0=0$ が滑らかなので相対特異点はない。
$t=1$ では相対特異イデアルは $(F,B_1,B_2,B_3)$ で、底の $(A_1,A_2,A_3)$ を合わせると正確に $Z$ となる。
$\Delta\to K_f$ は有限で $\dim Z=m-2$。
$K_f$ は一般に滑らかな連結完全交叉、従って既約で次元 $m-1$ なので $Z$ は底を支配しない。
一般 fiber は滑らかな三次曲面で幾何学的に既約。
平坦性により垂直な既約成分はなく、全空間とその像 $\Sigma_f$ は既約である。
$g$ 側も同様であり、$\deg Z=192>0$ から特異 fiber は存在する。
\end{proof}
一般のデータについての数値は次の通りである。
\begin{center}
\begin{tabular}{lll}\toprule
 & $n=6,m=2$ & $n=7,m=3$\\\midrule
$K_f,K_g$ & 次数8、種数5の曲線 & 次数8の K3 曲面\\
$\dim\Sigma_f$ & 3 & 4\\
$\Delta$ & 次数64、種数129の曲線 & 次数64、$K_\Delta=3H$ の曲面\\
$Z$ & 192個の被約点 & 次数192、種数577の滑らかな曲線\\
$D_0(\pi_X)$ & $E_f\sqcup E_g\sqcup Z$ & $E_f\sqcup E_g\sqcup Z$\\\bottomrule
\end{tabular}
\end{center}
$n=7$ の $\Delta$ では
$K_\Delta^2=576$、$c_2(\Delta)=768$、$\chi(\OO_\Delta)=112$、$p_g(\Delta)=111$、$h^1(\OO_\Delta)=0$。
両次元の一般葉の閉包は $\PP^6$ 内の型 $(3,3)$ の四次元完全交叉であり、
$X$ 上の底はそれぞれ $\PP(1,1,3)$、$\PP(1,1,1,3)$。
対数表示の反標準因子の次数は $(1,1,3,3)$、$(1,1,1,3,3)$ である。

\subsection{七次元の階数六}
ここでは $m=0$、$F=f(a)+g(b)+h(c)$、$a,b,c\in\CC^3$ である。
$X$ の滑らかさは各 $f,g,h$ が滑らかな平面三次曲線を定めることと同値である。
一方が特異なら対応する座標平面に $X$ の特異点が生じ、
逆に各 gradient の共通零点が原点だけなら全 gradient についても同じだからである。
曲線を $E_f,E_g,E_h$、二つずつの join を $J_{fg},J_{fh},J_{gh}$、
三曲線の join を $J_3=V(f,g,h)$ とする。
それぞれの次元は三、五である。
一次 Casimir と Poisson 超平面は存在しないが、$f,g,h$ の線形結合は Casimir である。
最高外積は
\begin{align*}
\pi^3&\longleftrightarrow\iota_E(df\wedge dg\wedge dh)\\
&=3(f\,dg\wedge dh-g\,df\wedge dh+h\,df\wedge dg)\\
&=3fgh\,d\log(g/f)\wedge d\log(h/f)\in H^0(\PP^8,\Omega^2(9)).
\end{align*}
$[f:g:h]:\PP^8\dashrightarrow\PP^2$ は $X$ 上で直線 $f+g+h=0$ に値を取る。
$f/g$ は有理 Casimir で、一般六次元葉の閉包は $\PP^8$ 内の型 $(3,3)$ の完全交叉。
$X$ 上の最高外積に対応する一形式は
$f\,dg-g\,df\in H^0(X,\Omega_X^1(6))$ である。
退化の台は
\begin{align*}
|D_4(\pi)|&=J_3\cup V(a)\cup V(b)\cup V(c),\\
|D_2(\pi)|&=J_{fg}\cup J_{fh}\cup J_{gh}\cup\PP_a^2\cup\PP_b^2\cup\PP_c^2,\\
D_0(\pi)&=E_f\sqcup E_g\sqcup E_h.
\end{align*}
$X$ 上では第一行の三つの $\PP^5$ が滑らかな三次四次元多様体となり、
第二行の三つの平面はそれぞれ曲線に制限され、既に join に含まれる。
従って $|D_2(\pi_X)|=J_{fg}\cup J_{fh}\cup J_{gh}$。

\subsection{非被約構造と primary 分解}
$A_i=f_{a_i}$、$B_j=g_{b_j}$、$C_l=h_{c_l}$ とし、
$\mathfrak a=(A_1,A_2,A_3)$、$\mathfrak b=(B_1,B_2,B_3)$、$\mathfrak c=(C_1,C_2,C_3)$ と書く。
\begin{proposition}\label{prop:pg-D4-primary}
周囲では $I_{D_4(\pi)}=(fB_jC_l,gA_iC_l,hA_iB_j)$。
$X$ 上ではイデアル層として
\begin{equation}
I_{D_4(\pi_X)}=I_{J_3}\cap\mathfrak a\cap\mathfrak b\cap\mathfrak c.\label{eq:pg-D4-primary}
\end{equation}
$J_3$ は被約成分で、他の三成分は三次四次元多様体を台とする primary な厚化である。
その一般点での横断長は八である。
\end{proposition}
\begin{proof}
生成元は命題\ref{prop:pg-deg-ideal}の $j=3$ である。
punctured affine cone の $A_i\neq0$ の開集合で $f=-g-h$ とおけば
$I_{D_4}=(g\mathfrak c,h\mathfrak b)$。
この開集合から $b,c$ 空間への写像は $a_i$ 方向の偏微分が単元なので滑らか、従って平坦である。
独立変数環 $\CC[b,c]$ では
\[
(g\mathfrak c,h\mathfrak b)=(g,h)\cap\mathfrak b\cap\mathfrak c.
\]
Euler 等式から $g\in\mathfrak b$、$h\in\mathfrak c$ なので一方の包含は明らか。
逆に $gr+hs$ が右辺に属せば、$\mathfrak b$ を法として非零多項式 $h(c)$ の乗法は単射だから $s\in\mathfrak b$。
同様に $r\in\mathfrak c$ となる。
平坦基底変換は有限交叉と可換し、全点ではいずれかの gradient 成分が単元となるので、これらを貼り合わせて\eqref{eq:pg-D4-primary}を得る。

例えば $\CC[b]/\mathfrak b$ は Hilbert 級数 $(1+t)^3$、長さ八の局所 Artin 環である。
$\mathfrak b$ による錐の商は
$\CC[a,c]/(f+h)\otimes_\CC\CC[b]/\mathfrak b$。
第一因子は滑らかな射影三次四次元多様体の錐の整域である。
第二因子の冪零根基の冪で filter すると、被約商で非零の元は非零因子であると分かる。
従って primary で、一般点の横断長は八。
$J_3$ は独立な三曲線の錐の積の射影化なので被約既約である。
\end{proof}
\begin{proposition}\label{prop:pg-D2-primary}
周囲では
\[
I_{D_2(\pi)}=(fB_j,gA_i,fC_l,hA_i,gC_l,hB_j,
 c_lA_iB_j,b_jA_iC_l,a_iB_jC_l).
\]
$X$ 上では
\begin{equation}
I_{D_2(\pi_X)}=I_{J_{fg}}\cap I_{J_{fh}}\cap I_{J_{gh}}
\cap(\mathfrak a+\mathfrak b)\cap(\mathfrak a+\mathfrak c)\cap(\mathfrak b+\mathfrak c).
\label{eq:pg-D2-primary}
\end{equation}
最後の三項はそれぞれ $E_h,E_g,E_f$ を台とする埋め込まれた primary 成分であり、各厚化の横断長は六十四である。
\end{proposition}
\begin{proof}
生成元は $j=2$ の公式による。
$A_i\neq0$ 上で $F=0$ を使うと
$I_{D_2}=(g,h,c_lB_j,b_jC_l)$。
$\mathfrak m_b=(b_1,b_2,b_3)$、$\mathfrak m_c=(c_1,c_2,c_3)$ とおくと、独立変数環で
\begin{align}
(\mathfrak m_c,g)\cap(\mathfrak m_b,h)&=(g,h,\mathfrak m_b\mathfrak m_c),\label{eq:pg-D2-reduced}\\
(g,h,\mathfrak m_b\mathfrak m_c)\cap(\mathfrak b+\mathfrak c)
&=(g,h,\mathfrak m_c\mathfrak b+\mathfrak m_b\mathfrak c).\label{eq:pg-D2-intersection}
\end{align}
第一式は $(g,h)$ で割って各因子を定数部分と正次数部分に分解すれば分かる。
第二式は $b,c$ 二重次数で確認する。
両次数が正なら $\mathfrak b+\mathfrak c$ のその成分は
$\mathfrak m_c\mathfrak b+\mathfrak m_b\mathfrak c$ に等しく、一方の次数が零なら第一式のイデアルは $(g)$ または $(h)$ であり、Euler 等式により $\mathfrak b+\mathfrak c$ に入る。
平坦基底変換により局所的に
$I_{D_2}=(\mathfrak m_c,g)\cap(\mathfrak m_b,h)\cap(\mathfrak b+\mathfrak c)$ となる。
初めの二項はそこでの $J_{fg},J_{fh}$ のイデアルである。
他の gradient 開集合でも同様に計算し、錐から $X$ へ降ろせば\eqref{eq:pg-D2-primary}を得る。

$\mathfrak b+\mathfrak c$ の商は
$\CC[a]/(f)\otimes\CC[b]/\mathfrak b\otimes\CC[c]/\mathfrak c$。
前命題と同じ filtration により primary で、横断長は $8\cdot8=64$ となる。
この成分は冗長ではない。
$b_jc_l$ は\eqref{eq:pg-D2-reduced}の右辺に属すが二重次数 $(1,1)$ をもち、$\mathfrak b+\mathfrak c$ に入らない。
台 $E_f$ は二つの join に含まれるので埋め込み成分である。他も同様。
\end{proof}
\begin{remark}\label{rem:pg-length49}
六十四は明示した primary 厚化の横断長であり、スキームと被約化との差の長さではない。
$E_f$ の一般点では
$\ker(\OO_{D_2(\pi_X)}\to\OO_{D_2(\pi_X)_{\mathrm{red}}})$ の横断長は四十九である。
実際、$Q_g=\CC[b]/\mathfrak b$、$Q_h=\CC[c]/\mathfrak c$ の極大イデアルを
$\mathfrak n_g,\mathfrak n_h$ とすると、\eqref{eq:pg-D2-reduced}、\eqref{eq:pg-D2-intersection}により核は横断方向で $\mathfrak n_g\otimes\mathfrak n_h$ に等しい。
各 Jacobian 環は長さ八で定数部分が一次元なので、長さは $7^2=49$ となる。
\end{remark}

\section{代数的な次数付き量子化}\label{sec:pg-quantization}
本節は形式的完備化だけでなく、原点を含む代数的な開集合上の量子化を構成する。
中心元の形式的存在には formality と Calabi--Yau 双対性を用い、有限次数の線形代数で代数化する。

\subsection{定義と主張}
\begin{definition}
連結次数付き可換代数 $C$ に次数零の Poisson 括弧があるとする。
$0\in U\subset\mathbb A^1_\hbar$、$K=\OO(U)$ 上の次数付き量子化とは、
$\deg\hbar=0$、$\mathcal C_0=K$、各 $\mathcal C_d$ が有限自由な
次数付き代数 $\mathcal C$ と $\mathcal C/(\hbar)\simeq C$ であって、
\[
\frac{[\widetilde a,\widetilde b]}\hbar\bmod\hbar=\{a,b\}
\]
を満たすものをいう。特殊 fiber の可換性と平坦性により持上げの選択に依存しない。
\end{definition}
連結次数付き代数 $A$ が Koszul であるとは、自明加群 $\CC=A/A_{\geq1}$ が線形自由分解をもつことをいう。
左右の injective dimension が $d$ で、左右それぞれ
\[
\Ext_A^i(\CC,A)=
\begin{cases}\CC(\ell)&i=d,\\0&i\neq d\end{cases}
\]
ならば、次元 $d$、AS parameter $\ell$ の AS-Gorenstein 代数という。
さらに大域次元 $d$ と有限 GK 次元をもてば AS-regular という。
次数移動は $M(\ell)_j=M_{j+\ell}$ とする。
以下の Ext は全内部次数を含み、injective/global dimension は通常の左右加群圏で取る。

\begin{theorem}[代数的量子化]\label{thm:pg-quantization}\label{pthm:quantization}
任意の Poisson 構造 $\pi_X$ に対し、$0\in U\subset\mathbb A^1_\hbar$、
$K=\OO(U)$ 上の二次関係式で定まる次数付き代数 $A$、および三次の中心元 $\Omega$ が存在して次を満たす。
\begin{enumerate}
\item $A$ は $S$ 上の $q$ の量子化であり、全ての閉点 $\lambda\in U$ で
$A_\lambda$ は Koszul AS-regular、
$\operatorname{gldim}A_\lambda=N$、$\operatorname{Hilb}_{A_\lambda}(t)=(1-t)^{-N}$、AS parameter $N$ である。
\item $\Omega\bmod\hbar=F$ で、各 $\Omega_\lambda$ は左右正則元である。
\item $B=A/(\Omega)$ は $S/(F)$ の量子化であり、各 $B_\lambda$ は AS-Gorenstein で、
\[
\operatorname{Hilb}_{B_\lambda}(t)=\frac{1-t^3}{(1-t)^N},\qquad
\operatorname{injdim}B_\lambda=N-1,\qquad
\text{AS parameter}=N-3=n-1.
\]
\end{enumerate}
\end{theorem}

\subsection{ブロック代数と PBW 性}
$R=\CC[u_1,\ldots,u_m]$ とし、$\sym_R$ は中心係数を外へ出して active 変数の語を平均する対称化である。
一ブロック $f(x_1,x_2,x_3,u)$ に対し
\begin{equation}
A_f=R[\hbar]\langle x_1,x_2,x_3\rangle/
\bigl([x_2,x_3]-\hbar\sym_R(f_{x_1}),\ [x_3,x_1]-\hbar\sym_R(f_{x_2}),\
[x_1,x_2]-\hbar\sym_R(f_{x_3})\bigr).\label{eq:pg-block-algebra}
\end{equation}
これは相対 cyclic potential
\begin{equation}
W=x_1x_2x_3-x_2x_1x_3-\hbar\sym_R(f)\label{eq:pg-potential}
\end{equation}
の active cyclic derivatives で定まる。
cyclic word は自由代数を交換子の線形包で割った類であり、
$\partial_i^{\mathrm{cyc}}(a_1\cdots a_r)
=\sum_{a_j=x_i}a_{j+1}\cdots a_ra_1\cdots a_{j-1}$ とする。

\begin{proposition}\label{prop:pg-PBW}
有限個のブロックに共通な $0\in U_{\mathrm{reg}}\subset\mathbb A^1$ を選び、以下を満たせる。
$g_s(x_s)=f_s(x_s,0)$ とおくと、全 $\lambda\in U_{\mathrm{reg}}$ で
\[
S_{s,\lambda}=\CC\langle x_{s1},x_{s2},x_{s3}\rangle/
([x_{s2},x_{s3}]-\lambda\sym((g_s)_{x_{s1}}),\text{循環置換})
\]
は三次元 Koszul AS-regular domain であり、3-Calabi--Yau である。
その Hilbert 級数は $(1-t)^{-3}$。
active 次数で filter し、$u_\alpha$ の filtration degree を零にすると
\[
\operatorname{gr}_xA_{f_s,\lambda}\simeq S_{s,\lambda}\otimes_\CC R.
\]
族全体にも対応する同型があり、$A_{f_s}|_{U_{\mathrm{reg}}}$ は
$R_{\mathrm{reg}}=\OO(U_{\mathrm{reg}})[u]$ 上自由である。
\end{proposition}
\begin{proof}
一つの $g$ について関係式の列は $M_\lambda x$ で表される。ただし
\[
M_\lambda=
\begin{pmatrix}
-\frac\lambda2g_{11}&-x_3-\frac\lambda2g_{12}&x_2-\frac\lambda2g_{13}\\
x_3-\frac\lambda2g_{21}&-\frac\lambda2g_{22}&-x_1-\frac\lambda2g_{23}\\
-x_2-\frac\lambda2g_{31}&x_1-\frac\lambda2g_{32}&-\frac\lambda2g_{33}
\end{pmatrix}.
\]
反対称成分が三つの交換子を与えるので関係式は独立である。
potential 由来の表示は\cite[Corollary 2.7]{PMoriSmith}の意味で standard であり、
同論文 Theorem 1.4 によって AS-regular 性は全 $\PP^2$ 上で $\rank M_\lambda\geq2$ と同値である。
階数一以下の incidence は $\mathbb A^1\times\PP^2$ の閉集合で、その底への像も閉じている。
$M_0$ は全 $\PP^2$ 上で階数二の交代行列なので、像を除けば原点を含む開集合を得る。
有限個のブロックについて交わりを取る。
3-CY 性と domain 性は\cite[Corollary 2.9]{PMoriSmith}による。
三生成元三二次関係式の三次元 AS-regular 代数では、最小分解を双対化すると最後の shift は三となり、
\[
0\longrightarrow S_{s,\lambda}(-3)\longrightarrow S_{s,\lambda}(-2)^3
\longrightarrow S_{s,\lambda}(-1)^3\longrightarrow S_{s,\lambda}
\longrightarrow\CC\longrightarrow0.
\]
これが Koszul 性と Hilbert 級数を与える。

$\lambda$ を固定する。$u=a$ への特殊化では、\eqref{eq:pg-potential}は
$g$ の三次 potential に active 次数二以下の項を加えたものである。
\cite[Theorem 1.1(i)]{PBergerTaillefer}により、その代数は $S_{s,\lambda}$ の PBW 変形である。
$S_{s,\lambda}$ の斉次基底を語の線形結合として持ち上げる。
最高 active 次数の関係式から次数帰納法で、これらは $A_{f_s,\lambda}$ を $R$ 上生成する。
線形関係があれば全 $u=a$ で PBW 基底間の関係となり、係数多項式が全点で零だから元の関係も零となる。
従って所要の PBW 性と自由性を得る。

族の斉次 $S_{s,U_{\mathrm{reg}}}$ の各次数部分は底上有限生成で、全 fiber の次元が一定である。
底環は $\CC[\hbar]$ の局所化で PID なので、各部分は有限自由である。
torsion があればその台で fiber 次元が増えるからである。
その各次数の基底を自由代数へ持ち上げ、同じ帰納法と全 $(\lambda,a)$ への特殊化を使えば、族上の自由性も従う。
\end{proof}
\begin{remark}\label{rem:pg-regular-all}
$g_s$ が滑らかな平面三次曲線を定めるとき、
\cite[Theorem 4.2(1)]{PMoriSmith}により $S_{s,\lambda}$ の AS-regular 性のために除く $\lambda$ はない。
$g_s=0$ の場合も常に多項式環である。
中心元の代数化・正則性のためには、以下でさらに開集合を縮める。
ambient の AS-regular 性は、その中心超曲面商の圏論的滑らかさとは別条件である。
\end{remark}
以下
$A_{\mathrm{reg}}=A_{f_1}\otimes_{R_{\mathrm{reg}}}\cdots\otimes_{R_{\mathrm{reg}}}A_{f_k}$
とする。異なるブロックは可換で、各全次数部分は底上有限自由であり、
$A_{\mathrm{reg}}/(\hbar)\simeq S$。関係式の一次項は\eqref{eq:pg-bracket}である。

\subsection{形式的な中心元}
この小節では $\widehat A_f$、$\widehat R=R[[\hbar]]$ という $\hbar$ 進完備化を用い、
$\widehat\otimes$、双加群分解および $\RHom$ も完備化の意味で解釈する。
内部次数を忘れた $d$-Calabi--Yau 性とは、包絡代数上 perfect で、包絡代数への導来双対が代数自身の $[-d]$ であることをいう。
以下では内部 shift も記録する。
\begin{lemma}\label{lem:pg-formal-center}
任意のブロックに対して、三次の中心元
$\widehat\Omega_f\in Z(\widehat A_f)$ で
$\widehat\Omega_f\bmod\hbar=f$ を満たすものが存在する。
$f(x,0)$ の滑らかさや孤立特異点は仮定しない。
\end{lemma}
\begin{proof}
\cite[Proposition 8.1.1(i)の証明]{EG}の方法を相対係数へ適用するため、
まず絶対 CY 性を確認する。
$A=\widehat A_f$、$r_i=\partial_i^{\mathrm{cyc}}W$ と書く。
相対完備自由双加群複体
\begin{equation}
0\longrightarrow P_3\xrightarrow{d_3}P_2\xrightarrow{d_2}P_1
\xrightarrow{d_1}P_0\xrightarrow\mu A\longrightarrow0\label{eq:pg-potential-resolution}
\end{equation}
を、
\[
P_0=A\widehat\otimes_{\widehat R}A,\quad
P_1=\bigoplus_{i=1}^3Ae_iA,\quad
P_2=\bigoplus_{i=1}^3Av_iA,\quad P_3=AwA
\]
で定める。$\deg e_i=1$、$\deg v_i=2$、$\deg w=3$ であり、
\[
d_1(e_i)=x_i\otimes1-1\otimes x_i,\quad
 d_3(w)=\sum_i(x_iv_i-v_ix_i),\quad
 d_2(v_i)=\sum_j\sum_{r_i\ni ax_jb}ae_jb.
\]
最後の和は語の各出現を係数付きで数える。
cyclic derivative の恒等式 $\sum_i[x_i,r_i]=0$ を非可換微分すると $d_2d_3=0$。
また $d_1d_2(v_i)=r_i\otimes1-1\otimes r_i=0$。
$\hbar=0$ では対角の相対 Koszul 分解である。
cycle を法 $\hbar$ で boundary として持ち上げ、差を $\hbar$ で割る操作を反復すると、
平坦性と完備性により完全性が従う。左端の核も全 $\hbar^rP_3$ に入り零である。

$\partial_j^{\mathrm{nc}}r_i=\sum_{r_i\ni ax_jb}a\otimes b$ とし
$\tau(a\otimes b)=b\otimes a$ と書く。
cyclic word の二切断位置を交換すれば
$\tau(\partial_j^{\mathrm{nc}}r_i)=\partial_i^{\mathrm{nc}}r_j$。
これが中央微分の自己双対性を与え、両端も符号を除き双対である。
逆順の双対微分は $(-d_1,d_2,-d_3)$ となり、各項に $(1,-1,-1,1)$ を掛けると元の複体に一致する。
従って $D_{\mathrm{rel}}=A\widehat\otimes_{\widehat R}A^{\mathrm{op}}$ に対し
\[
\Ext_{D_{\mathrm{rel}}}^q(A,D_{\mathrm{rel}})=
\begin{cases}A(3)&q=3,\\0&q\neq3.\end{cases}
\]
automorphism による twist はない。
絶対包絡代数 $D_{\mathrm{abs}}=A\widehat\otimes_{\CC[[\hbar]]}A^{\mathrm{op}}$ は、PBW 基底により
中心係数 $\CC[u^L,u^R][[\hbar]]$ 上平坦である。
従って $u_\alpha^L-u_\alpha^R$ は中心正則列で、その商は $D_{\mathrm{rel}}$。
その Koszul 分解と change of rings により
\[
\Ext_{D_{\mathrm{abs}}}^q(A,D_{\mathrm{abs}})=
\begin{cases}A(m+3)&q=m+3,\\0&q\neq m+3.\end{cases}
\]
相対分解の各項は $D_{\mathrm{abs}}$ 上 perfect なので、絶対的にも $(m+3)$-CY である。

対称化は完備化後に次数を保つ $\widehat R$-加群同型
$R[x_1,x_2,x_3][[\hbar]]\simeq A$ を与える。
積を移せば相対 Hochschild cochain による形式変形となる。
相対 polydifferential cochain の包含は、対角の Koszul 分解による HKR 計算から quasi-isomorphism である。
Kontsevich の affine formality のグラフ公式\cite[\S4.6, \S6]{PKontsevich}で active 変数の微分だけを用いれば、$R$ 上の formality を得る。
全変数の同時スケーリングの重み零の部分に制限する。
bar--Koszul 比較を次数付きで取れるので HKR 同型と包含はこの部分でも quasi-isomorphism であり、
Kontsevich 写像の各 Taylor 成分も重みを保つ。
$\hbar$ 倍して完備化した pronilpotent dg Lie 代数の Maurer--Cartan 同値を使えば、
形式 Poisson 構造と gauge transformation をともに重み零で選べる。
従って
\[
\Pi=\hbar\pi_f+O(\hbar^2),\qquad
\pi_f=\sum_{\mathrm{cyc}}f_{x_1}\partial_{x_2}\wedge\partial_{x_3}
\]
という相対形式 Poisson 構造の star product が $A$ と次数付き $R$-線形同型をもつ。
$\Pi$ は $\partial_{u_\alpha}$ 成分をもたないため、絶対グラフ公式で現れる微分も active 変数だけであり、
同じ $\Pi$ が絶対 Kontsevich class の代表となる。

\cite[Theorem 3, Proposition 1]{PDolgushev}を適用する。
台は滑らかな affine space で標準束は自明であり、双対化双加群が自明であることは上で確認した。
よって標準体積形式
$\omega_0=dx_1\wedge dx_2\wedge dx_3\wedge du_1\wedge\cdots\wedge du_m$ に対し、
単元 $\rho\in R[x][[\hbar]]^\times$ で $L_\Pi(\rho\omega_0)=0$ となるものが存在する。
全次数 $m+3$ の成分を取る。
$\Pi$ の重みは零であり、$\rho$ の次数零成分は $c(\hbar)\in\CC[[\hbar]]^\times$ だから
$L_\Pi\omega_0=0$ となる。
ここで単元性は $\rho\bmod\hbar$ が多項式環の単元、従って非零定数であることによる。

$\nu=dx_1\wedge dx_2\wedge dx_3$ とすると相対一形式
$\alpha=\iota_\Pi\nu$ は $d_R\alpha=0$ を満たす。
相対 polynomial de Rham 複体の active Euler homotopy によって
$\alpha=d_R\psi$、$\psi=\hbar f+O(\hbar^2)$、$\deg\psi=3$ を得る。
積分定数を $R[[\hbar]]$ で調整して最初の係数を $f$ にする。
$\Pi$ は $\psi$ の相対 Jacobian 場なので $\psi/\hbar$ は Casimir である。
twisted formality の第一 Taylor 成分はこの零次 Poisson cocycle を零次 Hochschild cocycle、すなわち中心元に送る。
gauge 同型で $A$ に戻したものが $\widehat\Omega_f$。
全写像が次数を保ち、法 $\hbar$ で関数上の HKR 写像は恒等なので、その次数は三、residue は $f$ である。
\end{proof}
\begin{remark}
\cite[Proposition 8.1.1]{EG}の孤立特異点の追加仮定は中心全体・cohomologyを扱う (ii),(iii) に課される。
上の証明は中心元一つの存在を扱い、中心全体が一つの potential で生成されるとは主張しない。
相対 formality および絶対 CY 双対性という外部入力を明示しており、計算実験だけから中心元を推測する議論ではない。
\end{remark}

\subsection{中心元の代数化と正則中心拡大}
\begin{lemma}\label{lem:pg-algebraize-center}
$0\in U_0\subset\mathbb A^1_\hbar$ 上の次数付き代数 $D$ の各次数部分が有限自由であり、
$D$ が次数一の元 $y_1,\ldots,y_r$ と中心元で生成されるとする。
完備化に次数 $d$ の中心元 $\widehat\eta$ があれば、
ある $0\in U\subset U_0$ 上で中心元 $\eta\in D_d$ を
$\eta\bmod\hbar=\widehat\eta\bmod\hbar$ となるように選べる。
\end{lemma}
\begin{proof}
$O=\CC[\hbar]_{(\hbar)}$ として有限自由加群間の写像
$(D_O)_d\to(D_O)_{d+1}^{\oplus r}$、$a\mapsto([a,y_i])_i$ の核を $L$ とする。
完備化の平坦性から完備化後の核は $L\otimes_O\CC[[\hbar]]$。
$L$ の有限個の生成元 $l_j$ を用いて
$\widehat\eta=\sum_ja_j(\hbar)l_j$ と書けば
$\eta_O=\sum_ja_j(0)l_j$ は同じ residue をもつ。
係数の分母は原点で消えないので、それらの零点を除けばよい。
\end{proof}
\begin{lemma}\label{lem:pg-central-linear}
$D$ を連結次数付き代数、$z\in D_1$ を中心正則元とする。
$C=D/(z)$ が次元 $d$、AS parameter $d$ の Koszul AS-regular なら、
$D$ は次元 $d+1$、AS parameter $d+1$ の Koszul AS-regular である。
\end{lemma}
\begin{proof}
$0\to D(-1)\xrightarrow{z}D\to C\to0$ より、
$\operatorname{Tor}^D_q(C,\CC)$ は $q=0,1$ でそれぞれ $\CC,\CC(-1)$、他は零である。
change-of-rings spectral sequence
\[
E^2_{p,q}=\operatorname{Tor}^C_p(\operatorname{Tor}^D_q(C,\CC),\CC)
\Longrightarrow\operatorname{Tor}^D_{p+q}(\CC,\CC)
\]
の各項は内部次数 $p+q$ に集中する。
微分は内部次数を保ち全ホモロジー次数を一減らすため全て零であり、
$D$ は Koszul、$j>d+1$ で $\operatorname{Tor}^D_j(\CC,\CC)=0$、$j=d+1$ では非零である。
$D^!$ を二次双対とすると、Koszul 双加群複体の項は $D\otimes(D_i^!)^*\otimes D$。
全内部次数を固定し右の $D$ 次数で filter すれば有限 filtration となり、
随伴複体の正ホモロジーは左 Koszul 分解で消える。
よって増大双加群複体は長さ $d+1$ の完全列であり、右加群として分裂完全なので
任意の左加群とテンソルしても完全である。従って通常の大域次元も $d+1$。
さらに $\RHom_D(C,D)\simeq C(1)[-1]$ と導来随伴から
\[
\Ext_D^{i+1}(\CC,D)\simeq\Ext_C^i(\CC,C)(1).
\]
右加群でも同様で、左右 AS 条件と parameter を得る。
$\operatorname{Hilb}_D=\operatorname{Hilb}_C/(1-t)$ なので GK 次元も有限である。
\end{proof}
\begin{lemma}\label{lem:pg-central-cubic}
$A$ を次元 $N$、AS parameter $N$ の AS-regular 代数、$\eta\in A_3$ を中心正則元とする。
$B=A/(\eta)$ は左右 injective dimension $N-1$、AS parameter $N-3$ の AS-Gorenstein 代数で、
$\operatorname{Hilb}_B=(1-t^3)\operatorname{Hilb}_A$ である。
\end{lemma}
\begin{proof}
$0\to A(-3)\xrightarrow{\eta}A\to B\to0$ が Hilbert 級数を与える。
導来随伴により任意の左 $B$-加群 $M$ について
$\Ext_B^i(M,B)\simeq\Ext_A^{i+1}(M,A)(-3)$。
右辺は $i\geq N$ で零、$M=\CC$、$i=N-1$ で $\CC(N-3)$ となる。
よって injective dimension はちょうど $N-1$ で AS 条件も成立する。右側も同様。
\end{proof}

\subsection{代数的量子化定理の証明}
\begin{proof}[定理\ref{thm:pg-quantization}の証明]
零構造は $A=S[\hbar]$、$\Omega=F$ でよい。
非零構造に対して $A_{\mathrm{reg}}$ を構成する。
PBW 基底により各 fiber で $u_1,\ldots,u_m$ は中心正則列で、
\[
A_\lambda/(u_1,\ldots,u_m)\simeq\bigotimes_{s=1}^kS_{s,\lambda}.
\]
右辺の自明加群の分解、双加群分解およびその双対は各因子のもののテンソル積で得られる。
従って次元 $3k$、AS parameter $3k$ の Koszul AS-regular 代数である。
補題\ref{lem:pg-central-linear}を $m$ 回適用して $A_\lambda$ の次元 $N=3k+m$ と AS parameter $N$ を得る。
Hilbert 級数は PBW 基底から $(1-t)^{-N}$ である。

補題\ref{lem:pg-formal-center}と\ref{lem:pg-algebraize-center}を各ブロックに適用し、開集合を共通に縮めると
$\Omega_s\in Z(A_{f_s}|_U)_3$、$\Omega_s\bmod\hbar=f_s$ を得る。
異なるブロックは可換なので $\Omega=\sum_s\Omega_s$ は $A$ の中心元で、residue は $F$。

少なくとも一つの $f_s$ は active 変数に依存する。
そうでなければ $F$ は中心変数だけに依存し、active 座標点に特異点をもつ。
その $s$ を固定し、PBW 基底での $\Omega_s$ の active 成分のうち原点で非零の係数が消えないようさらに縮める。
各 $\lambda$ で $\Omega_{s,\lambda}$ の active 次数は正となる。
最高次数が $\lambda$ によらず一定である必要はない。
$\lambda$ を固定し $x_s$ 次数だけで filter すると
\[
\operatorname{gr}_sA_\lambda\simeq
(S_{s,\lambda}\otimes_\CC R)\otimes_R A_{\mathrm{rest},\lambda},
\qquad A_{\mathrm{rest},\lambda}=\bigotimes_{t\neq s,R}A_{f_t,\lambda}.
\]
他の $\Omega_t$ は filtration degree 零なので、$\Omega_\lambda$ の最高記号は
$\Omega_{s,\lambda}$ の非零最高記号 $a_s\in S_{s,\lambda}\otimes R$ である。
$S_{s,\lambda}$ は domain だから多項式拡大も domain で、$a_s$ の左右乗法は単射。
$R$-自由な $A_{\mathrm{rest},\lambda}$ とテンソルしても単射なので、最高次数を比較して
$\Omega_\lambda$ の左右正則性が従う。非可換 domain のテンソル積が domain であるとは仮定していない。

従って各 fiber に
$0\to A_\lambda(-3)\xrightarrow{\Omega_\lambda}A_\lambda\to B_\lambda\to0$
があり、所要の Hilbert 級数を得る。
各次数の余核は PID $\OO(U)$ 上で全 fiber 次元が一定なので有限自由である。
法 $\hbar$ での商は $S\to S/(F)$ で、交換子の一次項もその商の括弧と一致する。
補題\ref{lem:pg-central-cubic}が残りの AS-Gorenstein の主張を与える。
\end{proof}
二次関係式は各生成元対に一つで、反対称成分により独立なので個数は $\binom N2$ である。
六次元では生成元八、関係式二十八、大域次元八、商の injective dimension 七、AS parameter 五。
七次元ではそれぞれ九、三十六、九、八、六となる。
\begin{remark}\label{rem:pg-quantization-scope}
定理\ref{thm:pg-quantization}は、個々の Poisson 構造に対する量子化の存在定理である。
代数の同型分類、中心全体の記述、$\qgr B_\lambda$ の滑らかさは含まない。
とくに本定理だけを根拠に全パラメータの Kuznetsov 成分の smooth proper 性や可換復元を結論することはできない。
\end{remark}

\section{特殊なブロック制限と Fermat 型}\label{sec:pg-special}
\subsection{零の三変数制限}
\begin{proposition}\label{prop:pg-zero-restriction}
$X$ が滑らかで、あるブロックの制限 $f_s(x_s,0)$ が恒等的に零ならば $m\geq3$ である。
$m=3$ ではこのような滑らかな例が存在する。
\end{proposition}
\begin{proof}
対応する座標平面 $\PP_s^2$ 上では全 active 偏微分が零であり、
残る偏微分は高々 $m$ 個の二次式 $F_{u_\alpha}(x_s,0)$ である。
$m\leq2$ なら三変数多項式環でこれらのイデアルの高さは高々二なので、非零の共通零点をもつ。
そこで $F$ の全偏微分が消え、滑らかさに反する。

存在について、九変数 $a_1,a_2,a_3,b_1,b_2,b_3,u_1,u_2,u_3$ の三次式
\begin{equation}\label{eq:pg-zero-restriction-example}
F=\sum_{i=1}^3(u_i a_i^2+u_i^3)+b_1^3+b_2^3+b_3^3
\end{equation}
を考える。全偏微分が零なら $2u_ia_i=0$ および $a_i^2+3u_i^2=0$ から
$a_i=u_i=0$、さらに $3b_i^2=0$ から $b_i=0$ となる。従って射影超曲面は滑らかである。
これは $k=2,m=3$ の正規形で、$a$ ブロックの $u=0$ 制限は零である。
\end{proof}
この例にも正規形定理\ref{thm:pg-normal}、全退化イデアルの命題\ref{prop:pg-deg-ideal}、
量子化定理\ref{thm:pg-quantization}を適用できる。
一方、平面楕円曲線を用いた零点の非交和表示の仮定は満たさないため、
退化スキームは一般式\eqref{eq:pg-deg-generators}から計算する必要がある。

\subsection{Fermat 型の中心元}
\begin{lemma}\label{lem:pg-fermat-center}
結合代数において
\[
[x,y]=\hbar z^2,\qquad [y,z]=\hbar x^2,\qquad [z,x]=\hbar y^2
\]
が成り立つとする。$S=x^3+y^3+z^3$、$T=\sym(xyz)$ とおくと $S-2\hbar^2T$ は中心元である。
ここで $\sym(xyz)$ は六個の語の平均である。
\end{lemma}
\begin{proof}
関係式と積の交換子を展開すると
\[
[S,x]=\hbar(zy^2z-yz^2y)=\hbar[zy,yz]=\hbar^2[zy,x^2]
\]
を得る。一方、$6T$ の六個の語を関係式で並べ替えると
$2T=xzy+zyx+\hbar x^3$ となる。従って $2[T,x]=[zy,x^2]$ であり、
$[S-2\hbar^2T,x]=0$ を得る。循環置換により $y,z$ との交換子も零である。
\end{proof}
特に
\begin{equation}\label{eq:pg-fermat}
F=\frac13\left(\sum_{\alpha=1}^m u_\alpha^3+
\sum_{s=1}^k\sum_{i=1}^3x_{si}^3\right)
\end{equation}
に対する量子化では、各ブロックで $[x_{s1},x_{s2}]=\hbar x_{s3}^2$ とその循環置換を課し、
異なるブロックと中心変数は可換とする。このとき
\begin{equation}\label{eq:pg-fermat-center}
\Omega=\frac13\left[\sum_{\alpha=1}^m u_\alpha^3+
\sum_{s=1}^k\left\{\sum_{i=1}^3x_{si}^3-2\hbar^2\sym(x_{s1}x_{s2}x_{s3})\right\}\right]
\end{equation}
は全パラメータで中心元であり、$\Omega|_{\hbar=0}=F$ を満たす。
定理\ref{thm:pg-quantization}の適切な開集合上で中心正則超曲面の族を与える。

\subsection{Fermat 型の非被約な零点}
\begin{proposition}\label{prop:pg-fermat-thickening}
\eqref{eq:pg-fermat}で $m>0$ ならば $\Delta=V(x_{si}^2\mid s,i)$ は
中心線形空間 $\PP_u^{m-1}$ を台とする横断長 $8^k$ の厚化である。
$m\geq2$ ならば $Z=\Delta\cap X$ は
$V(\sum_\alpha u_\alpha^3)\subset\PP_u^{m-1}$ の同じ横断厚化となる。
$m=0$ の場合には $\Delta=\varnothing$ である。
\end{proposition}
\begin{proof}
active 偏微分は $A_{s,i}=x_{si}^2$ なので、その商環は
\[
\CC[u]\otimes_\CC\bigotimes_{s=1}^k
\CC[x_{s1},x_{s2},x_{s3}]/(x_{s1}^2,x_{s2}^2,x_{s3}^2)
\]
である。各 Artin 因子の長さは八であり、この商では全 $x_{si}^3$ が零なので、
$F$ の像は $\frac13\sum_\alpha u_\alpha^3$ となる。これが厚化の記述を与える。
$m=0$ では偏微分の共通零点はアフィン原点だけであり、射影化は空である。
\end{proof}
六次元・階数四では、$Z$ は中心 $\PP^1$ 上の三点の各々に長さ六十四の厚みをもつ。
七次元・階数四では、中心 $\PP^2$ 上の Fermat 楕円曲線に横断長六十四の厚みをもつ。
次数はともに百九十二で、後者の算術種数は五百七十七である。
従って一般の場合との違いは次数や算術種数には現れず、被約性と滑らかさに現れる。

\section{完全交叉の数値不変量}\label{sec:pg-numerical}
本節は\S\ref{sec:pg-six-seven}で用いた数値を導出する。
滑らかな完全交叉 $W=(d_1,\ldots,d_c)\subset\PP^M$ に対し、
$H=c_1(\OO_W(1))$ とおくと
\begin{equation}\label{eq:pg-ci-formulas}
\omega_W\simeq\OO_W\left(\sum_i d_i-M-1\right),\qquad
c(T_W)=\frac{(1+H)^{M+1}}{\prod_i(1+d_iH)},\qquad
\int_W H^{M-c}=\prod_i d_i.
\end{equation}
これは順に随伴公式、接束の完全列と Euler 列、完全交叉の次数公式である。
また Koszul 分解から
$H^i(W,\OO_W(t))=0$（$0<i<\dim W$、$t\in\ZZ$）を得る。
以下の滑らかさを要する Hodge 数の主張は、一般のデータ、または滑らかさを個別に確認したデータに対して適用する。

\subsection{階数二の零点成分}
$n=6$ で $Y=(2,2,2,3)\subset\PP^7$ の場合、
\[
K_Y=H,\qquad c(T_Y)=\frac{(1+H)^8}{(1+2H)^3(1+3H)}
=1-H+7H^2-21H^3.
\]
従って $K_Y^3=24$、$e(Y)=-21\cdot24=-504$ である。
全方程式が二次以上なので $p_g(Y)=h^0(\OO_Y(1))=8$。
Lefschetz 定理と Hodge 分解により
$h^{1,1}=1$、$h^{1,0}=h^{2,0}=0$ であり、
\[
-504=e(Y)=4-2p_g(Y)-2h^{2,1}(Y)
\]
から $h^{2,1}(Y)=246$ を得る。

$n=7$ で $Y=(2,2,2,3)\subset\PP^8$ の場合は
\[
K_Y=0,\qquad c(T_Y)=\frac{(1+H)^9}{(1+2H)^3(1+3H)}
=1+6H^2-14H^3+48H^4,
\]
従って $e(Y)=48\cdot24=1152$ である。
Koszul 分解は $H^i(\OO_Y)=0$（$0<i<4$）を、
Lefschetz 定理は $h^{1,1}=1$、$h^{2,1}=0$ を与える。
$h^{3,1}$ を求めるため normal sequence を使う。
次数二の方程式は三個、次数三の方程式はそれらの一次式倍と独立な三次式一個だから
\[
h^0(\OO_Y(2))=\binom{10}{2}-3=42,\qquad
h^0(\OO_Y(3))=\binom{11}{3}-3\cdot9-1=137.
\]
Euler 列を $Y$ に制限し上記の消滅を使うと
$h^0(T_{\PP^8}|_Y)=9\cdot9-1=80$、$H^1(T_{\PP^8}|_Y)=0$ である。
$K_Y\simeq\OO_Y$ による $T_Y\simeq\Omega_Y^3$、$H^3(\OO_Y)=0$ と Hodge 対称性から
$H^0(T_Y)=0$ である。従って
\[
0\longrightarrow T_Y\longrightarrow T_{\PP^8}|_Y
\longrightarrow\OO_Y(2)^{\oplus3}\oplus\OO_Y(3)\longrightarrow0
\]
は $h^{3,1}(Y)=h^1(T_Y)=3\cdot42+137-80=183$ を与える。
最後に Hodge 対称性と Poincar\'e 双対性により
$1152=e(Y)=6+2h^{3,1}(Y)+h^{2,2}(Y)$、従って $h^{2,2}(Y)=780$ である。

\subsection{階数四の critical covers と零点}
$m=2$ の $K_f,K_g$ は型 $(2,2,2)$ の $\PP^4$ 内の曲線である。
従って次数八、$K_{K_f}=H$、$2g(K_f)-2=8$ で、種数五を得る。
$m=3$ では $\PP^5$ 内の同じ型の曲面であり、標準束が自明で $H^1(\OO)=0$ なので K3 曲面である。

$n=6,k=2,m=2$ の場合、$\Delta=(2^6)\subset\PP^7$ は次数六十四の曲線で $K_\Delta=4H$。
従って $2g(\Delta)-2=4\cdot64=256$、$g(\Delta)=129$ である。
$Z=(2^6,3)$ は長さ $3\cdot64=192$ の完全交叉で、一般のデータでは滑らか、すなわち百九十二個の被約点となる。

$n=7,k=2,m=3$ の場合、$\Delta=(2^6)\subset\PP^8$ は次数六十四の曲面で、
\[
c(T_\Delta)=\frac{(1+H)^9}{(1+2H)^6}=1-3H+12H^2.
\]
よって
\[
K_\Delta^2=9\cdot64=576,\qquad c_2(\Delta)=12\cdot64=768,\qquad
\chi(\OO_\Delta)=\frac{576+768}{12}=112.
\]
$H^1(\OO_\Delta)=0$ なので $p_g(\Delta)=111$ であり、これは
$h^0(\OO_\Delta(3))=\binom{11}{3}-6\cdot9=111$ からも確認できる。
最後に $Z=(2^6,3)\subset\PP^8$ は次数百九十二の曲線で $K_Z=6H$ だから
\[
2g(Z)-2=6\cdot192=1152,\qquad g(Z)=577.
\]
この式は特殊な非被約完全交叉でも算術種数の公式として保たれる。

\section{Hesse型Poisson量子化の明示的な三次中心元}
\label{sec:explicit-hesse-center}
本節では量子化の定理が与える三次中心元を、滑らかなHesse型
三変数ブロックについて閉公式で指定する。
この公式はSklyanin代数の既知の中心三次式
\cite[Theorem 1およびSection 2]{DeLaetLeBruyn}
（原式は\cite[p.~35, Remark (1)]{SmithTate}）
のPoisson量子化パラメータへの代入である。
従って中心元の存在という既知の結果自体を新規な主張とはしない。
本研究で有用なのは、$0$を含む明示的なパラメータ開集合と
$\hbar=0$での留数を正規化した公式を同時に得ることである。

\subsection{公式と適用開集合}
$t\in\CC$、$t^3\ne-1$とし、
\[
 f_t(x,y,z)=\frac{x^3+y^3+z^3}{3}+txyz
\]
と置く。$\hbar$を次数零のパラメータとして、
\[
 [y,z]=\hbar\left(x^2+\frac{t}{2}(yz+zy)\right),
 \quad
 [z,x]=\hbar\left(y^2+\frac{t}{2}(zx+xz)\right),
 \quad
 [x,y]=\hbar\left(z^2+\frac{t}{2}(xy+yx)\right)
 \tag{H}\label{eq:hesse-commutators}
\]
による次数付き代数を$S_{\hbar,t}$と書く。
ここで
\[
 S=x^3+y^3+z^3,\quad
 C=xyz+yzx+zxy,\quad
 D=yxz+xzy+zyx,\quad T=(C+D)/6.
\]
次の$\hbar$の多項式を定める：
\[
 U(\hbar,t)=1+\frac{\hbar^2t^2}{2}
  +\frac{\hbar^4t}{6}-\frac{\hbar^4t^4}{48},
 \qquad
 W(\hbar,t)=
 \frac{(4-\hbar^2t^2)
 (12t+\hbar^2(t^3-8))}{48}.
 \tag{HU}\label{eq:hesse-UW}
\]
また
\[
 \Delta_{\hbar,t}=
 \left(t+\frac{\hbar^2}{3}
       \left(1+\frac{t^3}{4}\right)\right)^3
  +\left(1-\frac{\hbar^2t^2}{4}\right)^3.
 \tag{HD}\label{eq:hesse-discriminant}
\]
と置く。$\Delta_{0,t}=1+t^3\ne0$である。

\begin{theorem}[Hesse中心三次式の閉公式]
\label{thm:explicit-hesse-center}
$t^3\ne-1$に対し、
\[
 U_t=\{\hbar\in\mathbb A^1_{\CC}:
   U(\hbar,t)(1-\hbar^2t^2/4)\Delta_{\hbar,t}\ne0\}
\]
とする。$0\in U_t$であり、
各$\hbar\in U_t$に対して
$S_{\hbar,t}$は三生成元のKoszul AS-regular domainで、
Hilbert級数$(1-z)^{-3}$を持つ。
元
\[
 \boxed{\quad
 \Omega_{\hbar,t}=
   \frac{U(\hbar,t)}3(x^3+y^3+z^3)
   +W(\hbar,t)\operatorname{sym}(xyz)
 \quad}
 \tag{HC}\label{eq:hesse-central-element}
\]
は$S_{\hbar,t}$の次数三の中心正則元であり、
$\Omega_{0,t}=f_t$である。
特に$t=0$では
$\Delta_{\hbar,0}=1+\hbar^6/27$で、
\[
 \Omega_{\hbar,0}
 =\frac{x^3+y^3+z^3}{3}
   -\frac{2\hbar^2}{3}\operatorname{sym}(xyz)
\]
を得る。
\end{theorem}
\begin{proof}
$q=\hbar t/2$、
$a=1-q$、$b=-(1+q)$、$c=-\hbar$と置く。
\eqref{eq:hesse-commutators}は
\[
 ayz+bzy+cx^2=0,\quad
 azx+bxz+cy^2=0,\quad
 axy+byx+cz^2=0
\]
となる。$\hbar\ne0$かつ
$(1-\hbar^2t^2/4)\Delta_{\hbar,t}\ne0$なら、
$abc\ne0$で、標準Sklyaninの
\[
 (a^3+b^3+c^3)^3\ne(3abc)^3
\]
が成り立つ。
実際、左辺の差は
$-27\hbar^3\Delta_{\hbar,t}$である。
従ってこの開部分では点曲線が滑らかな楕円曲線であり、
$S_{\hbar,t}$は三次元Sklyanin AS-regular domainである
\cite{SmithTate,DeLaetLeBruyn}。
$\hbar=0$の代数は可換多項式環である。

既知の中心元の公式は
\[
 G=a(b^3-c^3)C+
   b(c^3-a^3)D+
   c(a^3-b^3)S
 \tag{G}\label{eq:hesse-sklyanin-G}
\]
である\cite[Section 2]{DeLaetLeBruyn}。
一方、三つの定義関係式にそれぞれ$x,y,z$を左から掛けて足すと
\[
 C-D=\hbar S+\frac{\hbar t}{2}(C+D)
 \tag{R}\label{eq:hesse-CD}
\]
を得る。この等式で$G$を$S,C+D$で表し、
$a,b,c$を代入して係数を整理すると
\[
 G=-9\hbar\left(
       \frac{U(\hbar,t)}3S+
       \frac{W(\hbar,t)}6(C+D)\right)
   =-9\hbar\Omega_{\hbar,t}.
 \tag{I}\label{eq:hesse-G-identity}
\]
これは$\hbar$で割る操作をせずに成立する多項式恒等式である。
$\hbar\ne0$の滑らかな楕円Sklyanin開部分では
$G$が中心なので$\Omega_{\hbar,t}$も中心である。
各次数成分が$\hbar$について有限自由であることは
三次元SklyaninのHilbert級数と零ファイバーにおける
同じ級数、および元論文のPBW構成から従う。
従って各交換子$[\Omega_{\hbar,t},x]$等は
有限自由加群の切断として稠密開集合上で零であり、
$U_t$の全体で零となる。
$\Omega_{0,t}=f_t$は$T_0=xyz$と
$U(0,t)=1$、$W(0,t)=t$から分かる。
既知のSklyanin中心元$G$は楕円開部分で非零であり、
代数はdomainだから正則である。
零ファイバーでも$f_t$は非零因子である。
\end{proof}


\begin{corollary}[任意の滑らかな三変数三次式]
\label{cor:arbitrary-smooth-ternary-cubic}
$g(x_1,x_2,x_3)$が滑らかな平面三次曲線を定める
斉次三次式であるとする。$g$のJacobian Poisson括弧
\[
 [x_2,x_3]=\hbar\operatorname{sym}(\partial_{x_1}g),
 \quad [x_3,x_1]=\hbar\operatorname{sym}(\partial_{x_2}g),
 \quad [x_1,x_2]=\hbar\operatorname{sym}(\partial_{x_3}g)
\]
から得られる量子化は、$0$を含むZariski開集合上で
三次中心元の閉公式を持つ。
実際、Hesse標準形\cite[Lemma 1]{ArtebaniDolgachev}により
ある$M\in\mathrm{GL}_3(\CC)$、$\lambda\in\CC^\times$、
$t^3\ne-1$が存在して
\[
 g(x)=\lambda f_t(Mx)
\]
となる。$y=Mx$、$\hbar'=\lambda\det(M)\hbar$とおけば、
\[
 \boxed{\ \Omega^g_\hbar(x)=
  \lambda\Omega_{\hbar',t}(y)\big|_{y=Mx}\ }
 \tag{HN}\label{eq:arbitrary-ternary-central}
\]
は三次中心元で、$\Omega^g_0=g$である。
\end{corollary}
\begin{proof}
Poissonの定数三重ベクトル
$\partial_{x_1}\wedge\partial_{x_2}\wedge\partial_{x_3}$
は座標変換$y=Mx$の下で
$\det(M)\,
 \partial_{y_1}\wedge\partial_{y_2}\wedge\partial_{y_3}$
に移る。従って$y$の交換子は
\[
 [y_2,y_3]
 =\lambda\det(M)\hbar\,
   \operatorname{sym}(\partial_{y_1}f_t),
\]
およびその巡回置換となる。
対称化は線形座標変換と可換なので、量子化代数は
$S_{\hbar',t}$への次数付き同型を持つ。
その同型によって定理\ref{thm:explicit-hesse-center}の
中心元を戻し、留数が$g$になるよう$\lambda$倍する。
\end{proof}

\begin{corollary}[中心変数と独立ブロックの閉公式]
\label{cor:hesse-block-center}
相異なる三変数ブロック$x_s=(x_s,y_s,z_s)$と
中心変数$u$について
\[
 F=\sum_{s=1}^k
 \left(\frac{x_s^3+y_s^3+z_s^3}{3}
       +t_sx_sy_sz_s\right)+P(u),
 \qquad t_s^3\ne-1,
\]
$P$は任意の斉次三次式とする。
各ブロックに\eqref{eq:hesse-commutators}を課し、
異なるブロック間と$u$を可換に取る。
各$U_{t_s}$の共通開部分で
\[
 \Omega_\hbar
   =\sum_s\Omega_{\hbar,t_s}+P(u)
 \tag{HB}\label{eq:hesse-multiblock}
\]
は次数三の中心元で、$\Omega_0=F$である。
各$u$の一次形式$\ell_{si}(u)$を選んで
$x_{si}$を$x_{si}+\ell_{si}(u)$に置き換えた
非自明な混合項を含むPoissonブロックでも、
同じ置換による明示公式が成立する。
\end{corollary}
\begin{proof}
各$S_{\hbar,t_s}$の中心元は他のブロックと可換であり、
$P(u)$も中心なので和は中心である。
中心線形シフトは生成元の次数付き同型を定め、
同型によって中心元が中心元へ移る。
\end{proof}


\subsection{可換復元の候補となる点曲線}
\begin{proposition}[楕円点曲線のHesseパラメータ]
\label{prop:hesse-point-cubic}
定理\ref{thm:explicit-hesse-center}の条件で$\hbar\ne0$とする。
$S_{\hbar,t}$のpoint schemeを定める滑らかな平面三次曲線は
\[
 E_{\hbar,t}:\quad
 x^3+y^3+z^3+3\,\vartheta(\hbar,t)xyz=0,
\quad
 \boxed{\quad
 \vartheta(\hbar,t)=
 \frac{t+\frac{\hbar^2}{3}(1+t^3/4)}
 {1-\hbar^2t^2/4}
 \quad}
 \tag{EP}\label{eq:hesse-point-parameter}
\]
である。$\vartheta(0,t)=t$と形式的に延長され、
\[
 \Delta_{\hbar,t}
 =\left(1-\frac{\hbar^2t^2}{4}\right)^3
  \left(1+\vartheta(\hbar,t)^3\right).
 \tag{ED}\label{eq:hesse-point-delta}
\]
特に$\hbar\in U_t$なら、右辺のHesse三次曲線は滑らかである。
\end{proposition}
\begin{proof}
滑らかな楕円型三次元Sklyanin代数のpoint schemeは
\[
 abc(x^3+y^3+z^3)-(a^3+b^3+c^3)xyz=0
\]
である\cite{SmithTate,DeLaetLeBruyn}。
定理の証明で与えた
$a=1-\hbar t/2$、$b=-(1+\hbar t/2)$、$c=-\hbar$
について
\[
 abc=\hbar(1-\hbar^2t^2/4),\qquad
 a^3+b^3+c^3
 =-3\hbar t-\hbar^3(1+t^3/4)
\]
を代入し、共通係数$\hbar$で割ると表示を得る。
最後の恒等式は\eqref{eq:hesse-discriminant}の定義から従う。
\end{proof}
\begin{corollary}[一楕円ブロックのrotation付き可換復元]
\label{cor:hesse-reconstruction-established}
$m\ge2$、$V(\partial P)=\{0\}$、$h\in U_t$とする。
$S_{h,t}[u]/(\Omega_{h,t}+P)$のrotation付き残余圏は
$Y=V(f_{\vartheta(h,t)}+P)$のrotation付きKuznetsov成分と同値である。
平行移動の有限位数性は仮定しない。
\end{corollary}
\begin{proof}
$h\ne0$では$\Omega_{h,t}$は標準Sklyanin中心元の非零スカラー倍なので、
その商は滑らかな点曲線$E_{h,t}$のTHCRであり、自己同型は平行移動である。
$h=0$でも中心商は同じ$E_{0,t}$の通常の斉次座標環である。
従って系\ref{cor:elliptic-central-cubic}の全仮定を満たす。
正規化、rotationの共役、周期性、中心変数の追加と生成性の詳しい証明は
独立論文\texttt{elliptic\_cubic\_reconstruction.tex}第2--4節にも記す。
\end{proof}
\begin{remark}[ファイバーごとの復元と相対復元]
上の結論は各複素ファイバーでの同値であり、底上の一つの相対同値は未構成である。
$h=0$の周囲の代数のpoint schemeは$\PP^2$であり、中心商の曲線$E_{0,t}$と区別する。
本質的な中心混合項と二つ以上の無限位数ブロックの全生成は未解決である。
\end{remark}

\begin{remark}[適用範囲と新規性]
\label{rem:hesse-scope}
一般の$f_s(x_s,u)$には$u$とactive変数の本質的な
混合項があり、必ずしも上の中心シフト形にはならない。
その場合の中心元は依然として形式的存在と
有限自由加群間の線形方程式による代数化で定まる。
また、\eqref{eq:hesse-central-element}は
周囲の量子化と中心元の明示式であり、
残余dg圏の相対saturationや
rotation付き可換復元を結論するものではない。
中心元の基本公式は古典的Sklyanin理論に由来し、
本節のHesseパラメータへの翻訳の新規性は未確認である。
\end{remark}


\part{rotation付き残余圏と可換モデル}
\section{可換cubic模型の存在を妨げるrotationの不変量}
\label{sec:rotation-obstruction}

本節では、次数付き代数の条件と、可換cubicのrotationを持つという条件を区別する。
三次中心正則元を持つAS-regular代数の商というだけでは、可換模型の存在は従わない。
さらに、商の射影圏が滑らかな可換cubicのcoherent圏と同値であり、周囲の代数が
Calabi--Yauであっても、与えられたrotationを保つ可換模型が存在しない例がある。
以下の二つの反例は有限位数のZhang twistで構成する。
これらが本稿で構成する特定のPoisson量子化族のファイバーであることは主張しない。

\subsection{Serre関手とrotationの必要条件}

\begin{proposition}[Serre--rotation障害]
\label{prop:serre-rotation-obstruction}
 \(\cC\) をSerre関手 \(S_{\cC}\) を持つproperなdg圏、
 \(\rho\) をその自己同値とする。\(N=n+2\ge5\) とし、
 滑らかなcubic \(Y^n\subset\PP^{N-1}\) とのenhanced pair同値
 \[
 (\cC,\rho)\simeq(\Ku(Y),\rho_Y)
 \]
 が存在するなら、
 \begin{equation}\label{eq:serre-rotation-necessary}
 S_{\cC}\simeq\rho^{\,3-N}[N-2],\qquad
 \Theta_N(\cC,\rho):=S_{\cC}\rho^{N-3}[2-N]\simeq\Id
 \end{equation}
 である。特に \(\rho^3\simeq[2]\) だけでは十分でない。
\end{proposition}
\begin{proof}
 可換斉次座標環 \(R_Y=\CC[z_1,\ldots,z_N]/(F_Y)\) のKrull次元は
 \(N-1\)、標準加群は \(R_Y(3-N)\) である。
 滑らかさによって原点以外の錐は正則であり、次数付きAuslander--Reiten双対性は
 \(\Dsg^{\ZZ}(R_Y)\) のSerre関手を
 \((3-N)[N-2]\) と与える
 \cite[Theorem 7, p.~161]{UeyamaIsolated}。
 Orlovのwindowで内部次数シフト \((1)\) は \(\rho_Y\) に対応するので、
 \(S_{\Ku(Y)}\simeq\rho_Y^{3-N}[N-2]\) を得る。
 Serre関手は同値によって保存されるから、pair同値を通してこの式を移せばよい。
\end{proof}

\subsection{有限位数Zhang twistと次数シフト}

以後 \(\zeta\) は \(\zeta^3=1\)、\(\zeta\ne1\) を満たす。
\(S=\CC[x_0,\ldots,x_{N-1}]\)、\(F=\sum_i x_i^3\)、
\(X=V(F)\subset\PP^{N-1}\) とする。
\(q_i\in\{1,\zeta,\zeta^2\}\) を選び、
\(\tau(x_i)=q_i x_i\) と置く。\(\tau^3=\Id\)、\(\tau(F)=F\) である。
\(f:X\to X\) を \(f^*=\tau\) に対応する射影自己同型とする。
右Zhang twist \(A=S^{\tau}\) の積を、斉次元に対して
\begin{equation}\label{eq:right-zhang-product}
 a*b=a\,\tau^{\deg a}(b)
\end{equation}
で定義する。右辺の積は元の可換多項式環の積である。

\begin{lemma}\label{lem:zhang-cubic-structure}
 \(A\) は左右noetherianなKoszul AS-regular domainで、
 global dimensionは \(N\)、Hilbert級数は \((1-t)^{-N}\) である。
 同じ次数付きベクトル空間内の元 \(\Omega=F\in A_3\) は中心正則元であり、
 \[
 B=A/(\Omega)\simeq (S/(F))^{\tau},\qquad
 H_B(t)=\frac{1-t^3}{(1-t)^N}.
 \]
 さらに
 \begin{equation}\label{eq:zhang-qgr-pair}
 \qgr B\simeq\operatorname{coh}X,\qquad
 (\mathcal K_B,\rho_B)\simeq
 (\Ku(X),f_*\rho_X)
 \end{equation}
 である。従って \(\mathcal K_B\) はsmooth properで、
 \(\rho_B^3\simeq[2]\) を満たす。
\end{lemma}
\begin{proof}
 積\eqref{eq:right-zhang-product}は結合的であり、生成元の関係式は
 \[
 x_i*x_j=(q_j/q_i)x_j*x_i\qquad(i<j)
 \]
 となる。従って \(A\) はscalarによる可逆なOre拡大の反復である。
 順序付き単項式 \(x_0^{a_0}*\cdots*x_{N-1}^{a_{N-1}}\) は
 元の可換単項式の非零scalar倍なので基底をなし、domain性とHilbert級数が従う。
 Ore拡大のHilbert基底定理から左右noetherian性を得る。
 多項式環のKoszul分解をtwistすると、\(k\) の長さ \(N\) の線形自由分解になる。
 その双対もtwistされたKoszul分解であるから
 \[
 \operatorname{Ext}^j_A(k,A)=
 \begin{cases}k(N)&j=N,\\0&j\ne N\end{cases}
 \]
 となり、左側にも同じ式が成り立つ。
 同じtwistを双加群Koszul分解に施した有限分解、またはOre拡大のglobal dimensionの公式により、
 global dimensionは \(N\) である。これはZhang twistがAS-regular性を保つことの
 この場合の具体的な確認である \cite{ZhangTwist,RRZ}。

 任意の斉次 \(a\in A_d\) に対して
 \[
 a*F=a\tau^d(F)=aF,
 \qquad F*a=F\tau^3(a)=Fa
 \]
 であるから、\(F\) は中心的である。
 domain性から正則であり、\((F)\) の次数付きベクトル空間は元の多項式環の
 \((F)\) と一致する。このため商も右twistであり、表示したHilbert級数を得る。
 なお \(x_i*x_i*x_i=q_i^3x_i^3=x_i^3\) なので、
 \(\Omega\) はtwist後の生成元の三乗和としても同じ元である。

 次数付き右 \(R=S/(F)\)-加群 \(M\) に、
 \(m*a=m\tau^{\deg m}(a)\) という作用を入れた加群を \(\mathsf T(M)\) と書く。
 \(\mathsf T\) は逆twistを逆関手に持つ完全同値で、有限生成性とtorsionを保つ。
 従って \(\operatorname{gr}R\simeq\operatorname{gr}B\) および
 \(\operatorname{coh}X\simeq\qgr B\) を与える。

 次数シフトの規約を \(M(1)_d=M_{d+1}\) とする。
 \(\mathsf T(M)(1)\) をuntwistすると、次数 \(d\) の元 \(m\in M_{d+1}\) に対して
 \[
 m\cdot a=m*\tau^{-d}(a)=m\tau(a)
 \]
 となる。従って、untwist後の \((1)\) は \(M(1)_{\tau}\) であり、
 coherent圏では \(f_*(-\otimes\OO_X(1))\) に対応する。
 特に \(\OO_B(i)\) は \(\OO_X(i)\) に対応する。
 \(f_*\OO_X\simeq\OO_X\) なので、\(f_*\) は \(\OO_X\) によるleft mutationと可換である。
 よって残余圏とrotationについて\eqref{eq:zhang-qgr-pair}が従う。
 \(f_*\) は \(\rho_X\) と可換で、\(f_*^3\simeq\Id\) であるため、
 \((f_*\rho_X)^3\simeq\rho_X^3\simeq[2]\) となる。
\end{proof}

\begin{remark}[右twistにおけるNakayamaの規約]
\label{rem:zhang-nakayama}
 上の右twistのNakayama自己同型は
 \begin{equation}\label{eq:right-twist-nakayama}
 \mu_A=\tau^{-N}\xi_{\det\tau},\qquad
 \mu_A(x_i)=\frac{\prod_jq_j}{q_i^N}x_i,
 \end{equation}
 と取れる。ここで \(\xi_c\) は次数 \(d\) 上で \(c^d\) 倍する自己同型である。
 \cite[Theorem 5.4]{RRZ}は左twist
 \(a\circ b=\sigma^{\deg b}(a)b\) に対して
 \(\mu=\sigma^N\xi_{\det\sigma}^{-1}\) を与える。
 \(a\mapsto\tau^{-\deg a}(a)\) は右twistから \(\sigma=\tau^{-1}\) の左twistへの同型であり、
 この公式を移すと\eqref{eq:right-twist-nakayama}になる。
 このように左右の規約を固定してから用いる。
\end{remark}

\subsection{三次元の滑らかな反例}

\begin{theorem}[Serre--rotation条件を破るsmoothな非可換cubic]
\label{thm:zhang-threefold-counterexample}
 \(N=5\)、\(\tau=\operatorname{diag}(\zeta,1,1,1,1)\) とし、
 補題\ref{lem:zhang-cubic-structure}の \(B\) を取る。
 この \(B\) はnoetherianなKoszul AS-regular五次元代数の三次中心正則商であり、
 \(\qgr B\simeq\operatorname{coh}X\) は滑らかなFermat cubic threefoldのcoherent圏と同値である。
 しかし、いかなる滑らかな可換cubic threefold \(Y\) に対しても
 \[
 (\mathcal K_B,\rho_B)\not\simeq(\Ku(Y),\rho_Y).
 \]
\end{theorem}
\begin{proof}
 代数の諸性質とsmooth proper性は補題\ref{lem:zhang-cubic-structure}で示した。
 \(T=f_*|_{\Ku(X)}\) と書くと、同補題により
 \(\rho_B=T\rho_X\)、\(T^3\simeq\Id\)、\(T\rho_X\simeq\rho_XT\) である。
 underlying圏は \(\Ku(X)\) なので、そのSerre関手は
 \(S_X\simeq\rho_X^{-2}[3]\) である。
 もし可換pair模型が存在すれば、命題\ref{prop:serre-rotation-obstruction}により
 \[
 \rho_X^{-2}[3]\simeq (T\rho_X)^{-2}[3]
       \simeq T^{-2}\rho_X^{-2}[3]
 \]
 となり、\(T^2\simeq\Id\) を得る。\(T^3\simeq\Id\) と合わせると \(T\simeq\Id\) となる。

 これはHochschild homologyへの作用に反する。
 例外対象のHochschild homologyは次数零だけにあるため、加法性とHKR同型から
 \[
 \HH_1(\Ku(X))\simeq\HH_1(X)\simeq H^{2,1}(X)
 \]
 である。Griffithsのresidue同型\cite{Griffiths}は
 \[
 J(F)_1\xrightarrow{\sim}H^{2,1}(X),\qquad
 a\longmapsto\operatorname{Res}_X\frac{a\Omega_{\mathrm{proj}}}{F^2},
 \quad
 \Omega_{\mathrm{proj}}=
 \sum_{i=0}^4(-1)^i x_i\,dx_0\wedge\cdots\wedge\widehat{dx_i}\wedge\cdots\wedge dx_4
 \]
 と書ける。\(\tau(F)=F\)、\(\tau^*\Omega_{\mathrm{proj}}=(\det\tau)\Omega_{\mathrm{proj}}
 =\zeta\Omega_{\mathrm{proj}}\) である。
 \(J(F)=\CC[x_0,\ldots,x_4]/(x_0^2,\ldots,x_4^2)\) だから、
 \(f^*\) の \(H^{2,1}(X)\) 上の固有値は
 \[
 \zeta^2\ \text{が1個},\qquad \zeta\ \text{が4個}
 \]
 となる。\(f_*\) の作用はその逆であり、いずれにしても恒等ではない。
 従って \(T\not\simeq\Id\) であり、矛盾する。
\end{proof}

\begin{remark}
 この例では \(\rho_B^3\simeq[2]\) とsmooth proper性の両方が成立する。
 一方、\eqref{eq:right-twist-nakayama}のNakayama自己同型は非自明であり、
 周囲の \(A\) はCalabi--Yauではない。
 次の例は、この最後の条件を加えても一般的な存在命題が成立しないことを示す。
\end{remark}

\subsection{Calabi--Yauな周囲を持つ四次元の反例}

\begin{theorem}[Calabi--Yauな非可換cubicのrotation障害]
\label{thm:zhang-cy-fourfold-counterexample}
 \(N=6\)、\(\tau=\operatorname{diag}(\zeta,\zeta^2,1,1,1,1)\) とし、
 補題\ref{lem:zhang-cubic-structure}の \(A\)、\(B\) を取る。
 このとき \(A\) は六次元Calabi--Yauなnoetherian Koszul AS-regular代数であり、
 \(\Omega\) は三次中心正則元である。
 \(\qgr B\) は滑らかなFermat cubic fourfold \(X\) のcoherent圏と同値であり、
 \(\mathcal K_B\) はsmooth properな2-Calabi--Yau圏である。
 \(\rho_B^3\simeq[2]\) および\eqref{eq:serre-rotation-necessary}も満たす。
 しかし、いかなる滑らかな可換cubic fourfold \(Y\) に対しても
 \[
 (\mathcal K_B,\rho_B)\not\simeq(\Ku(Y),\rho_Y).
 \]
 より具体的に、Hochschild \emph{homology}の次数零部分について
 \begin{equation}\label{eq:cy-fourfold-hh0-obstruction}
 \dim\HH_0(\mathcal K_B)^{\rho_B}=8,
 \qquad
 \dim\HH_0(\Ku(Y))^{\rho_Y}=20
 \end{equation}
 である。
\end{theorem}
\begin{proof}
 \(\det\tau=1\)、\(\tau^6=\Id\) だから、
 \eqref{eq:right-twist-nakayama}によって \(\mu_A=\Id\) である。
 AS-regular性とNakayamaの自明性はCalabi--Yau性を与える
 \cite[Lemma 1.2, Theorem 5.4]{RRZ}。
 残りの代数的性質と \(\mathcal K_B\simeq\Ku(X)\) は
 補題\ref{lem:zhang-cubic-structure}による。
 cubic fourfoldのKuznetsov成分は2-Calabi--Yauなので
 \(S_{\mathcal K_B}\simeq[2]\) であり、
 \(\rho_B^{-3}[4]\simeq[2]\) からSerre--rotation条件も成立する。

 Hochschild homologyの作用を計算する。
 Mukai--HKR同定と残余成分の射影によって
 \begin{equation}\label{eq:cubic-fourfold-hh0-split}
 \HH_0(\Ku(X))\simeq
 H^{2,2}(X)_{\mathrm{prim}}\oplus (A_2\otimes_{\ZZ}\CC)
 \end{equation}
 と書ける。ここで前者は20次元、後者は2次元である。
 \(A_2\) はcubic fourfoldのKuznetsov成分に標準的に含まれる二つのMukai classの張る格子である。
 \cite[Proposition 3.12]{Huy}により、\(\rho_X\) はprimitive部分上で恒等、
 \(A_2\) 上で三つの根を巡回置換するので、後者の固有値は \(\zeta,\zeta^2\) である。
 この記述は一般の滑らかなcubic fourfoldにも成り立つため、
 \(\dim\HH_0(\Ku(Y))^{\rho_Y}=20\) を得る。

 射影自己同型 \(f\) の作用は分解\eqref{eq:cubic-fourfold-hh0-split}を保ち、
 \(A_2\) 上で恒等である。
 実際、その二つの標準classは直線 \(\ell\subset X\) 上の
 \(\OO_\ell(1)\)、\(\OO_\ell(2)\) の残余射影から作られる。
 \(f\) は直線を直線に送り、全ての直線のcohomology classは同じである。
 残余射影のcohomology上の作用は \(\OO_X,\OO_X(1),\OO_X(2)\) から定まり、
 これらを \(f\) が固定するので、二つのMukai classも固定される。
 Mukai--HKR同定に入るTodd classは射影自己同型で不変であるから、
 これは \(\HH_0\) 上の作用についても成立する。

 primitive部分ではresidue表示
 \[
 J(F)_3\xrightarrow{\sim}H^{2,2}(X)_{\mathrm{prim}},\qquad
 a\longmapsto\operatorname{Res}_X\frac{a\Omega_{\mathrm{proj}}}{F^3}
 \]
 を用いる。ここでresidue表示はpole filtrationの商を取った
 \(J(F)_3\simeq\operatorname{Gr}_F^2H^4(X)_{\mathrm{prim}}\) と、
 Hodge分解による \(\operatorname{Gr}_F^2\simeq H^{2,2}_{\mathrm{prim}}\) の合成を意味する。
 射影自己同型は両方の構造を保つ。
 \(\det\tau=1\)、\(\tau(F)=F\) なので、ここで \(f^*\) の作用は
 \(J(F)_3\) 上の通常の置換 \(a\mapsto\tau(a)\) に一致する。
 FermatのJacobian環は \(J(F)=\CC[x_0,\ldots,x_5]/(x_0^2,\ldots,x_5^2)\) であり、
 次数三の基底は平方自由な三変数単項式20個である。
 その固有値を分類すると
 \[
 \begin{array}{c|c|c}
 \text{単項式に含まれる変数}&\text{個数}&\tau\text{の固有値}\\\hline
 x_0,x_1\text{の両方}&4&1\\
 x_0,x_1\text{のどちらも含まない}&4&1\\
 x_0\text{だけ}&6&\zeta\\
 x_1\text{だけ}&6&\zeta^2
 \end{array}
 \]
 となる。従って \(f_*\) のprimitive部分の固定空間も8次元である。

 補題\ref{lem:zhang-cubic-structure}から \(\rho_B=f_*\rho_X\) である。
 primitive部分で \(\rho_X\) は恒等なので、\(\rho_B\) の固定空間は8次元である。
 \(A_2\) 上では \(f_*\) が恒等であり、\(\rho_X\) に固有値1はない。
 よって \eqref{eq:cy-fourfold-hh0-obstruction}が成立する。
 enhanced pair同値は \(\rho\) の \(\HH_0\) 上の作用を共役するので、
 固定空間の次元の不一致が同値の不存在を示す。
\end{proof}

\subsection{smooth properな連結族で保存される不変量}

前の反例と可換cubicのpairとの違いは、連結な圏の族を通じても検出できる。
Hochschild homologyの各次数の次元が個別に一定であることを仮定せず、
perfect複体のEuler標数を用いる。

\begin{proposition}[rotation不変Euler標数の局所一定性]
\label{prop:rotation-euler-locally-constant}
 \(S\) を複素スキームとし、\(\mathcal C_S\) をsmooth properな代数的dg圏の族とする。
 \(\rho_S\) は可逆perfect bimoduleで与えられる自己同値の族とし、
 整数 \(d>0\)、\(a\in\ZZ\) と自然同型 \(\rho_S^d\simeq[2a]\) を仮定する。
 幾何学的点 \(s\) に対して
 \[
 \chi_{\rho}(\mathcal C_s)
   =\sum_{i\in\ZZ}(-1)^i
        \dim_{\kappa(s)}\HH_i(\mathcal C_s)^{\rho_s}
 \]
 と定義する。この和は有限であり、\(s\mapsto\chi_{\rho}(\mathcal C_s)\) は局所一定である。
 特に \(S\) が連結なら一定である。
\end{proposition}
\begin{proof}
 相対Hochschild複体 \(H=\HH(\mathcal C_S/S)\) は \(S\)-perfectであり、
 導来底変更と両立する \cite[Lemmas 3.4, 3.6]{PerryIHC}。
 実際、étale局所的に \(\mathcal C_S=\Perf(E)\) と書けば、
 \(H=E\otimes^{\mathbf L}_{E^{\mathrm{op}}\otimes^{\mathbf L}_S E}E\) である。
 smoothnessによって対角は有限cell双加群の直和因子として表され、propernessによって
 各cellをtraceした複体は \(S\)-perfectとなる。perfect性は直和因子で保たれる。
 これらの構成から底変更も従う。

 \(\rho_S\) が誘導する \(H\) の自己同型を \(r\) と書く。
 Hochschild homologyの加法性により \([1]\) の誘導作用は \(-1\)、従って
 \([2a]\) の作用は恒等である \cite[Section 4, Example 5.6]{CisinskiTabuada}。
 よって \(r^d=1\) であり、
 \[
 e=\frac1d\sum_{j=0}^{d-1}r^j\in\operatorname{End}_{D(S)}(H)
 \]
 は冪等元である。\(\Perf(S)\) は冪等完備なので、その像 \(H^{\mathrm{inv}}\) はperfectである。
 幾何ファイバーで \(e\) の作用は有限巡回群の不変部分への射影であるから、
 \[
 H^{-i}(H^{\mathrm{inv}}\otimes_S^{\mathbf L}\kappa(s))
     =\HH_i(\mathcal C_s)^{\rho_s}.
 \]
 perfect複体のファイバーEuler標数は、局所的な有限階数ベクトル束複体の
 階数の交代和であるため局所一定である。これが主張である。
\end{proof}

\begin{corollary}[Calabi--Yau反例と可換pairの連結成分の分離]
\label{cor:zhang-cy-component-obstruction}
 定理\ref{thm:zhang-cy-fourfold-counterexample}のpairは、
 \(\rho^3\simeq[2]\) を相対的に備えた連結smooth proper代数的dg圏の族によって、
 滑らかな可換cubic fourfoldの標準pairに結ぶことができない。
 実際、両者の不変Euler標数はそれぞれ
 \[
 \chi_{\rho_B}(\mathcal K_B)=10,
 \qquad
 \chi_{\rho_Y}(\Ku(Y))=22
 \]
 である。
\end{corollary}
\begin{proof}
 cubic fourfoldのKuznetsov成分のHochschild homologyは次数 \(-2,0,2\) にのみあり、
 次数 \(-2,2\) は各一次元である。
 標準rotationはsymplecticなので、この二つの部分に恒等として作用する
 \cite[Proposition 3.12]{Huy}。
 定理\ref{thm:zhang-cy-fourfold-counterexample}の \(f\) も、
 \(\det\tau=1\)、\(\tau(F)=F\) により \(H^{3,1}(X)\) のresidueを固定し、
 その複素共役 \(H^{1,3}(X)\) も固定する。
 従って \(f_*\rho_X\) についても \(\HH_{\pm2}\) の不変部分は各一次元である。
 次数零の不変部分は\eqref{eq:cy-fourfold-hh0-obstruction}で計算したので、
 Euler標数はそれぞれ \(1+8+1=10\)、\(1+20+1=22\) となる。
 命題\ref{prop:rotation-euler-locally-constant}を適用する。
\end{proof}

\begin{remark}[量子化族との論理的関係]
\label{rem:counterexample-quantization-scope}
 定理\ref{thm:zhang-cy-fourfold-counterexample}は、
 「多項式環と同じHilbert級数を持つCalabi--Yau AS-regular代数の三次中心正則商であり、
 射影圏がsmooth properである」という条件による一般的な可換pair復元を反証する。
 またSerre--rotation条件を追加するだけでも十分でない。
 underlyingなKuznetsov成分には既に可換模型 \(X\) があるが、指定されたrotationが異なるのである。

 一方、この代数が本稿の特定のPoisson構造から作る、零パラメータで可換な
 中心三次量子化族のファイバーであることは示していない。
 \(\tau\) の有限位数性は中心元の構成に本質的であり、
 \(\zeta\) を任意の複素パラメータに置き換えても三次中心元が残るわけではない。
 従って本例から、その特定の量子化族に対する存在問題の否定は導かない。
 研究目標では、扱う量子化族と残余pairの相対構成を指定し、
 その上での存在問題を述べる必要がある。
\end{remark}

\section{完全忠実比較を同値へ強める条件}
\label{sec:serre-generation}

多ブロック族では、既に構成した完全忠実関手の像が全体を生成するかが問題になる。
外部積による直接的な生成証明とは別に、properness、Serre関手および連結性から
本質的全射性を判定できる。ここでの \(\HH^0\) はcohomologyであり、
前節の反例に用いたhomologyの \(\HH_0\) とは異なる。

\begin{proposition}[Serreと連結性による全生成判定]
\label{prop:serre-connected-generation}
 \(\mathcal A\) を非零のsmooth properな、冪等完備pretriangulated dg圏、
 \(\mathcal D\) をproperな、冪等完備pretriangulated dg圏とする。
 \(\mathcal D\) はSerre関手 \(S_{\mathcal D}\) を持つとし、
 dg完全忠実関手 \(\Phi:\mathcal A\to\mathcal D\) が
 \[
 \Phi S_{\mathcal A}\simeq S_{\mathcal D}\Phi
 \]
 を満たすと仮定する。
 さらに環 \(\HH^0(\mathcal D)\) が \(0,1\) 以外の冪等元を持たないなら、
 \(\Phi\) はdg Morita同値である。
 特に \(\HH^0(\mathcal D)=\CC\) なら、この連結性条件は満たされる。
 \(\mathcal D\) のsmoothnessは仮定しなくてよい。
\end{proposition}
\begin{proof}
 まず \(\Phi\) の像がadmissibleであることを示す。
 \(\mathcal A\) はsmooth properなので、compact generator \(G\) と
 smooth proper dg代数 \(E=\REnd_{\mathcal A}(G)\) を用いて
 \(\mathcal A\simeq\Perf(E)\) と書ける。
 \(D\in\mathcal D\) に対し、\(\RHom_{\mathcal D}(\Phi G,D)\) は
 \(\mathcal D\) のpropernessにより \(\CC\)-perfectである。
 smoothなdg代数上では、このようなpseudo-perfect加群はperfectである。
 実際、対角 \(E\) を有限cell双加群の直和因子として表し、この加群をtensorすると、
 \(\CC\)-perfectな係数を持つ有限cell \(E\)-加群の直和因子として表される。
 したがって \(\mathcal A\) の対象で表現され、\(\Phi\) の右随伴を与える。
 opposite圏で同じ議論を行えば左随伴も得る。
 この構成はdg加群の表現可能性として行われるので、随伴もenhancedである。

 \(\mathcal B=\Phi(\mathcal A)^\perp\) とする。
 \(D\in\mathcal B\)、\(A\in\mathcal A\) に対し、Serre双対性から
 \[
 \Hom_{\mathcal D}(D,\Phi A[i])
 \simeq
 \Hom_{\mathcal D}(S_{\mathcal D}^{-1}\Phi A[i],D)^*
 \simeq
 \Hom_{\mathcal D}(\Phi S_{\mathcal A}^{-1}A[i],D)^*=0
 \]
 となる。従って右直交と左直交は一致する。
 admissibilityによる各対象の分解三角形は、二方向のHom消滅によって分裂する。
 よって \(\mathcal D\) は \(\Phi(\mathcal A)\) と \(\mathcal B\) の直交直和である。
 全てのcohomological shiftのHomが消えるため、これはdg Moritaの意味でも直和分解となる。

 この直和分解は対角双加群の分解を与え、第一成分への射影が
 \(\HH^0(\mathcal D)=\operatorname{End}(\Delta_{\mathcal D})\) の冪等元になる。
 \(\mathcal A\ne0\) なのでこの冪等元は零でない。
 もし \(\mathcal B\ne0\) なら恒等でもないため、仮定に反する。
 従って \(\mathcal B=0\) であり、\(\Phi\) は本質的全射である。
\end{proof}

\begin{corollary}[任意位数の独立ブロックに対する十分条件]
\label{cor:arbitrary-order-conditional-model}
 \(N=3r+m\ge5\) とし、全ての \(h_j^6\ne-27\) とする。
 \(P(u)\) は原点のみを臨界点に持つ中心三次式とし、\(m=0\) では省く。
 独立ブロックの代数
 \[
 B=\left(\bigotimes_{j=1}^rS_{h_j}\right)[u]
     \Big/\left(\sum_jG_{h_j}+P(u)\right)
 \]
 の通常次数付きMF圏を \(\mathcal K_B\)、次数シフトを \(\rho_B\) とする。
 次の三条件を仮定する。
 \begin{enumerate}
 \item \(\mathcal K_B\) はproperである。
 \item Serre関手が存在し、\(S_{\mathcal K_B}\simeq\rho_B^{3-N}[N-2]\) である。
 \item \(\HH^0(\mathcal K_B)\) は非自明な冪等元を持たない。
 \end{enumerate}
 このとき、平行移動の有限位数性を仮定せずに、
 \[
 (\mathcal K_B,\rho_B)\simeq
 \left(\Ku\!\left(V\!\left(\sum_j
 (x_j^3+y_j^3+z_j^3+h_j^2x_jy_jz_j)+P(u)\right)\right),\rho_Y\right)
 \]
 が成り立つ。
\end{corollary}
\begin{proof}
 各楕円ブロックのdegree shiftを保つ同値、外部積の完全忠実性、
 および有限群による通常次数への降下から、
 表示した滑らかな可換cubic \(Y\) について
 \(\Phi:\Ku(Y)\to\mathcal K_B\) が得られ、\(\Phi\rho_Y\simeq\rho_B\Phi\) である。
 中心三次式 \(P\) は一つの可換因子として扱えばよく、
 可換側のequivariant Thom--Sebastianiと非可換側の外部積の完全忠実性は
 Fermat型への分解を必要としない。
 \(Y\) の滑らかさは各変数群の勾配の孤立性から従う。
 \(\Ku(Y)\) は非零のsmooth proper圏であり、
 \eqref{eq:serre-rotation-necessary}と第2条件はSerre関手との互換性を与える。
 命題\ref{prop:serre-connected-generation}を適用すればよい。
\end{proof}

\begin{remark}[この十分条件の検証状態]
 この系は、二つ以上の無限位数ブロックの全生成性を無条件に証明したものではない。
 その場合に残る課題を、三つの明示した圏論的条件へ帰着する結果である。
 与えられた量子化族の全ファイバーについて、properness、Serre公式、
 および \(\HH^0\) の連結性を確認することは依然必要である。
 前節の反例が示すように、一般のAS-regular中心商に対しては、
 比較する完全忠実関手の存在も自動的には得られない。
\end{remark}

\section{可換復元の定式化と主結果}\label{sec:intro}

滑らかなcubicのKuznetsov成分を
\[
 D^b(X)=\langle\Ku(X),\OO_X,\OO_X(1),\ldots,\OO_X(n-2)\rangle
\]
の規約で取る。rotationは
\[
 \rho_X=\mathbf L_{\OO_X}\circ(-\otimes\OO_X(1))
\]
である。すべての圏はdg enhancementと冪等完備化を保って扱う。
可換cubicについては、\((\Ku(X),\rho_X)\) が \(X\) の射影同型類を決定する
\cite[Theorem 1.7]{LZ}。このTorelli定理は可換モデルの一意性を与えるが、
任意の非可換な圏の組がその像に属することを述べない。
幾何側のenhanced同値は、包含・射影と合成してFourier--Mukai核で表されるため、
同定理のFourier--Mukai条件を満たす。
本稿は、この存在問題の形式近傍、飽和な代数的族のétale近傍、
および明示的な楕円ブロックの族を扱う。

\subsection{形式的な存在と一意性}

\begin{theorem}[次数付き変形の形式的可換モデル]\label{thm:main}
 \(X=V(F)\subset\PP^{n+1}_{\CC}\) を滑らかなcubic、\(n\ge3\) とする。
 \(R\) を剰余体が \(\CC\) の有限次元可換局所Artin代数とし、
 \(B_R=\bigoplus_{d\ge0}(B_R)_d\) を中心的な \(R\)-代数で、
 \((B_R)_0=R\)、各次数部分が有限自由、かつ指定された次数付き同型
 \[
 B_R\otimes_R\CC\simeq\CC[x_1,\ldots,x_{n+2}]/(F)
 \]
 を持つものとする。次を定義する。
 \[
 \cP_R=D(\QGr B_R)^c,\qquad
 \cK_R=\langle\OO_R,\ldots,\OO_R(n-2)\rangle^\perp\subset\cP_R,\qquad
 \rho_R=\mathbf L_{\OO_R}\circ(1).
 \]
 このとき \((\cK_R,\rho_R)\) はsmooth properな \(R\)-linear dg圏と自己同値の組である。
 あるmarkedな滑らかなcubic変形
 \(X_R\subset\PP^{n+1}_R\) と、markingを保つ同値
 \[
 (\cK_R,\rho_R)\simeq(\Ku(X_R),\rho_{X_R})
 \]
 が存在する。\(X_R\) のmarked射影同型類は一意である。
 これらはArtin底の変更と両立する。
\end{theorem}

ここで \(\QGr\) は全次数付き加群をtorsion加群で割ったGrothendieck圏であり、
上付き \(c\) はcompact対象を表す。Artin底上では
\(\cP_R\) を \(D^b(\qgr B_R)\) 全体に置き換えない。
例えば可換な定数族でも、底 \(R\) が非正則なら両者は異なる。
後述する体上の有限パラメータでは、残余圏を通常のgraded matrix factorization圏、
すなわちOrlov同値で \(D^b(\qgr B)\) に埋め込まれる成分として扱う。
特異なファイバーで、この成分と \(D(\QGr B)^c\) の直交成分を混同しない。

\begin{corollary}[形式族]\label{cor:formal}
 \(B_\hbar\) が次数ごとに \(\CC[[\hbar]]\)-flatな、定理の斉次座標環の形式的変形なら、
 \[
 F_\hbar\in\CC[[\hbar]][x_1,\ldots,x_{n+2}]_3,\qquad F_\hbar\bmod\hbar=F
 \]
 が存在し、各 \(\CC[\hbar]/(\hbar^{r+1})\) 上で残余圏の組を
 \(V(F_\hbar)\) の組と同一視できる。これらの同値は互換的に選べる。
 形式的cubicのmarked同型類は一意である。
\end{corollary}

これは形式的な主張である。係数の収束や、指定された代数的パラメータ空間への
降下、任意の非零パラメータでの可換モデルの存在を含まない。
また \(F_\hbar=F\) を意味しない。

\subsection{代数的族における局所的存在}

定理\ref{thm:etale}は、有限型複素底 \(S\) 上のsmooth properな代数的dg圏の族
\(\cK_S\) と、可逆perfect bimoduleで表された自己同値 \(\rho_S\) を入力とする。
ある点 \(s\) で \((\cK_s,\rho_s)\simeq(\Ku(X),\rho_X)\) なら、
\(s\) のétale近傍 \(T\) 上に滑らかなcubic \(Y/T\) と
\[
 (\cK_T,\rho_T)\simeq(\Ku(Y/T),\rho_{Y/T})
\]
が存在する。従って、可換cubicで実現される複素点はZariski開部分をなす。
これはgraded flatな代数の族から、相対smooth properなdg模型を作る課題を
省略するものではない。

また \(n\ne4\) では \(\HH^2(\Ku(X))=J(F)_3\)、\(\HH^1=0\) であり、
rotationを忘れた形式的復元も成立する（定理\ref{thm:unpolarized-formal}）。
同じ条件下では、代数的dg圏の族についてもrotationを最初から指定せずに
étale局所モデルを作れる（系\ref{cor:etale-unpolarized}）。
これらは局所的な実現の結果であり、rotationなしの大域的な一意性とは分けて扱う。

\subsection{有限パラメータを含む明示的族}

\[
 S_h=\CC\langle x,y,z\rangle/
 ([x,y]-hz^2,\ [y,z]-hx^2,\ [z,x]-hy^2)
\]
と置き、\(T=\sym(xyz)\) を六つの順列の平均、
\(G_h=x^3+y^3+z^3-2h^2T\) とする。
中心変数 \(u_1,\ldots,u_m\) と斉次三次式 \(P(u)\) に対して
\[
 B_{h,P}=S_h[u_1,\ldots,u_m]/(G_h+P(u))
\]
を考える。\(P=\sum_i u_i^3\) の場合を \(B_{h,m}\) と略記する。

\begin{theorem}[一ブロックと任意の中心三次式の可換モデル]\label{thm:global}
 \(m\ge2\)、\(n=m+1\)、\(h\in\CC\)、\(h^6\ne-27\) とする。
 \(P\in\CC[u_1,\ldots,u_m]_3\) のアフィン臨界点が原点だけであると仮定する。
 \(B_{h,P}\) の残余dg圏とrotationは、滑らかなcubic
 \[
 Y_{h,P}=V\!\left(x^3+y^3+z^3+h^2xyz+P(u)\right)
 \subset\PP^{m+2}
 \]
 の \((\Ku(Y_{h,P}),\rho_{Y_{h,P}})\) と同値である。
 この \(Y_{h,P}\) は、同じrotation付き残余圏を持つ滑らかな可換cubicとして
 射影同型を除いて一意である。
\end{theorem}

この定理はsixfoldでは \(m=5\)、sevenfoldでは \(m=6\) を含む。
中心変数どうしの任意の三次混合項を許すが、一つのSklyaninブロックと中心部分が
分離されているという仮定を保持する。\(P=\sum u_i^3\) はその特殊例である。
複数の非可換ブロックについても、対応する可換cubicの残余圏からの完全忠実埋込みを得る
（命題\ref{prop:multiblock-embedding}）。次の部分的な有限位数条件の下では、その像が全体を生成する。

\begin{theorem}[一つの任意位数ブロックを許す多ブロック可換モデル]\label{thm:multiblock-model}
 \(r\ge2\)、\(m\ge0\)、\(n=3r+m-2\) とする。
 各 \(j\) について \(h_j^6\ne-27\) とし、楕円曲線
 \[
 E_{h_j}=V(x_j^3+y_j^3+z_j^3+h_j^2x_jy_jz_j),\qquad
 o_j=[1:-1:0],\quad p_j=[1:-1:-h_j]
 \]
 上の点 \(p_j\) は、高々一つを除いて有限位数であると仮定する。
 すなわち因子を並べ替えた後、\(j\ge2\) にのみ有限位数を仮定し、\(p_1\) の位数には条件を課さない。
 \(m>0\) の場合、\(P\in\CC[u_1,\ldots,u_m]_3\) は原点だけに臨界点を持つとし、
 \(m=0\) の場合は中心部分を省く。
 異なるSklyaninブロックの元を可換とした代数
 \[
 B_{\mathbf h,P}=
 (S_{h_1}\otimes\cdots\otimes S_{h_r})[u_1,\ldots,u_m]
 \big/\left(\sum_jG_{h_j}+P\right)
 \]
 のrotation付き残余dg圏は、滑らかなcubic
 \[
 Y_{\mathbf h,P}=V\!\left(
 \sum_{j=1}^r(x_j^3+y_j^3+z_j^3+h_j^2x_jy_jz_j)+P(u)
 \right)\subset\PP^{3r+m-1}
 \]
 の \((\Ku(Y_{\mathbf h,P}),\rho_{Y_{\mathbf h,P}})\) と同値である。
 この可換モデルは射影同型を除いて一意である。
\end{theorem}

ここで「有限パラメータ」は形式変数でなく実際の複素数であることを意味し、
楕円曲線上の点が「有限位数」であることとは異なる。
定理は一つの無限位数ブロックを許す。
具体的に \(h_1=2\) は無限位数（補題\ref{lem:order-infinite}）、
\(h_j^6=5\ (j\ge2)\) は位数五（補題\ref{lem:order-five}）となる。
従って二ブロックのsixfold、三ブロックのsevenfoldにも、真に無限位数の因子を含む例が得られる。
二つ以上の無限位数ブロックを任意に許す主張は、本定理には含まれない。
証明は第\ref{sec:finite-center}節に置く。

\begin{proposition}[例外パラメータのhomological singularity]\label{prop:bad-intro}
 \(m\ge0\) とする。Fermat中心部分の族 \(B_{h,m}\) で \(h^6=-27\) とすると、\(S_h[u]\) はAS-regularで
 \(G_h+\sum u_i^3\) は中心正則元であるが、そのgraded matrix factorization圏は
 properではなく、さらに
 \[
 \operatorname{hdim}(\qgr B_{h,m})=\infty
 \]
 である。ここでhomological dimensionは、非零のExtが存在する次数の上限である。
 従って、滑らかな射影cubicのKuznetsov成分とは同値にならない。
\end{proposition}

これはAS-regularな周囲、三次の中心正則元、期待されるHilbert series、
および \(\rho^3\simeq[2]\) だけでは、大域的な可換モデルの存在を保証しないことを示す。
本稿の可換点からのPoisson量子化族に限定した、良い有限パラメータ全体での存在問題は別に残る。
一般のAS-regular中心三次商については、smooth properかつCalabi--Yauな例にも
定理\ref{thm:zhang-cy-fourfold-counterexample}の反例がある。

\subsection{先行研究と本稿の位置}

Orlov \cite{Orlov} とCassidy--Conner--Kirkman--Moore \cite{CCKM} は、
中心正則元による商、特異圏、matrix factorization、残余成分の比較を与える。
体上のwindowのdg模型にはCăldăraru--Tu \cite[Theorem 5.8]{CT} を用いる。
本稿の形式的な構成では、これらの体上の同値をArtin底上へ無条件に延長せず、
De Deken--Lowen \cite{DDL} のabelian deformationとcompact generationを用いる。
smoothnessとpropernessの持ち上げはnilpotent descent
\cite{Mathew}による。

残余圏のHochschild cohomologyはBallard--Favero--Katzarkov \cite{BFK}、
幾何学的Kodaira--Spencer写像はHuybrechts--Rennemo \cite{HR} と、
三次元の場合にはJacovskis--Lin--Liu--Zhang \cite{JLLZ} を用いる。
四次元の余分な二方向に対するrotationの作用は
Huybrechts \cite[Proposition 3.12]{Huy} から確認する。
actual dg変形をHochschild形式変形と同一視する際のcurvature問題
\cite{KL}は、Lurie \cite{Lurie} とBlanc--Katzarkov--Pandit
\cite{BKP} の比較写像の単射性で処理する。

明示的族の楕円ブロックはEtingof--Ginzburg \cite{EG} と
Sklyaninの点曲線の公式 \cite{Walton} に依拠する。
本稿では次数シフトの共役、中心変数を加える際の生成性、
有限群による通常次数への降下を証明する。
可換な係数環についてのThom--Sebastiani定理を、
非可換な係数環に対して仮定を確認せず引用することはしない。

有限中心代数の外部積生成については、Elagin--Lunts--Schn\"urer
\cite[Theorem 4.18]{ELS} が、perfect field上のnoetherian J-2 schemesに載る
coherentな代数に対し、ungradedな有界導来圏の生成定理を証明している。
本稿では、\(\ZZ^r\) の任意の粗化次数を保ち、さらに一方の係数代数には
有限中心性を仮定しない形を証明する（定理\ref{thm:one-sided-boxes}）。
\cite[Theorem 4.18]{ELS} は双方の因子をcoherentな構造代数とするため、
任意のnoetherian係数代数を許すここでの主張は、その定理の直接の適用ではない。
新たな入力はBeilinsonの対角分解 \cite{Beilinson} を非可換係数付きで用いる
多項式拡大の補題である。中心のNoether正規化と分離冪等元によるretractを組み合わせ、
斉次中心局所化の核を次元について帰納する。
対象の細分次数へのliftや、係数代数の有限global dimensionは仮定しない。
有限位数THCRの有限中心性 \cite[Corollary 2.3]{SmithTate} をこの片側のみに適用し、
一つの任意位数ブロックを含む多ブロックのMF比較を得る。
He--Ma--Ye \cite{HMY} のtensor productに対する一般化Kn\"orrer周期性は、
中心元の次数が二である非可換quadricとClifford変形を扱う。
本稿の三次中心元への直接の適用ではない。
周囲のnoetherian性については、Chan \cite{ChanMulti} の多重Serre同値と
面・帯による議論を任意個の楕円因子へ展開し、有限位数の仮定なしで証明する。
Keeler \cite{KeelerMulti} のTheorem 2.7は任意個の可換な可逆bimoduleの
NC-amplenessを判定するが、同論文のTheorem 3.3とCorollary 3.5のnoetherian性は
二因子についての主張である。本稿では任意個の積に必要な帰納法を記す。
基礎となる対角分解と有限中心のungraded生成定理は既知である。
一方有限中心の次数付き定式化および本適用の先行研究上の新規性は確認継続中である。

étale局所復元では、Toën--Vaquié \cite{TV} のperfect objectのmoduliと
smoothnessの無限小判定 \cite{Stacks} を用いる。
新たに示すのは、同値bimoduleとrotation intertwinerを同時にパラメータ化したスタックが、
形式的復元定理によって底上smoothになることである。
Pirozhkov \cite{Pirozhkov} およびLin--Zhang \cite{LZ} のTorelliは
二つの可換モデルの一意性に関する結果であり、この局所的な存在の証明とは役割が異なる。
Huybrechts \cite[Section 9]{HuyUpdate} の四次元での一般的な非可換変形の問題提起と、
本稿のrotation付きの変形問題との接写像の詳細な比較は研究ノートに残す。

\paragraph{主張の状態.}
本稿の定理には以下で証明を与える。これらを一つの一般次元の復元定理として
整理したことの先行研究上の新規性は未確認であり、
既知の基礎定理自体を新しい成果として扱わない。
本稿で構成する可換点からのPoisson量子化族について、
良い有限パラメータ全体での存在は未解決として残す。

\section{次数付き変形から相対残余圏を作る}\label{sec:relative}

\subsection{nilpotentな底変更によるcompactnessの検出}

\begin{lemma}\label{lem:descent}
 \(R\) を上の局所Artin代数、\(k=R/\fm\) とする。
 \(A\) を \(R\)-linear dg代数、\(M\) を \(A\)-加群とし、
 \(A_0=A\otL_R k\)、\(M_0=M\otL_R k\) と書く。
 \(M_0\) が \(A_0\)-compactなら \(M\) は \(A\)-compactである。
 また導来底変更 \(-\otL_R k\) は保守的である。
\end{lemma}
\begin{proof}
 \(\fm\) の冪による有限濾過から
 \(R\in\thick_{D(R)}(k)\) である。
 \(A\)-加群の任意の族 \((N_\alpha)\) に対し、
 \[
 \bigoplus_\alpha\RHom_A(M,P\otL_RN_\alpha)
 \longrightarrow
 \RHom_A\!\left(M,P\otL_R\bigoplus_\alpha N_\alpha\right)
 \]
 が同値となる \(P\in D(R)\) の全体はthickである。
 \(P=k\) では導来随伴と \(M_0\) のcompactnessによって同値となるので、
 \(P=R\) でも同値となる。保守性も、
 \(P\otL_RM=0\) となる \(P\) のthick部分圏を使って同様に従う。
 これは \cite[Propositions 3.28, 3.35; Remark 3.29]{Mathew}
 の特別な場合である。
\end{proof}

\begin{remark}
 この補題は \(k\) が \(R\)-perfectであるとは仮定していない。
 また普通の商 \(A/\fm A\) だけを調べる命題ではなく、
 必ず導来ファイバー \(A\otL_Rk\) を使う。
\end{remark}

\subsection{abelian deformationと明示的な生成子}

\begin{proposition}\label{prop:qgr}
 定理\ref{thm:main}の \(B_R\) について、
 \(\QGr B_R\) は \(\operatorname{Qcoh}(X)\) のflat abelian deformationである。
 \[
 E_R=\bigoplus_{i=0}^{n+1}\OO_R(i)
 \]
 は \(D(\QGr B_R)\) のcompact generatorであり、
 \[
 \Lambda_R=\REnd(E_R),\qquad
 \Lambda_R\otL_Rk\simeq\REnd_X(E_0)
 \tag{\(\ast\)}\label{eq:derived-basechange}
 \]
 がdg代数として成り立つ。
\end{proposition}
\begin{proof}
 まず \(B_R\) は左右noetherianである。
 \(\fm^aB_R/\fm^{a+1}B_R\) は、
 flatnessにより \(B_0\) の有限個の直和である。
 有限濾過と \(B_0\) のnoetherian性を使えばよい。
 次数付き生成元を持ち上げる同じ議論から、
 \(B_R\) は有限生成で、付随する \(\ZZ\)-代数のtailsも通常のtorsion商で記述できる。

 \(X\) の射影正規性により \(B_0\) の \(\ZZ\)-代数は
 \((\OO_X(i))_{i\in\ZZ}\) のdirected Hom-categoryである。
 この列は \(\ZZ\)-generatingであり、\(i\le j\) に対して
 \[
 \Ext_X^{1,2}(\OO_X(i),\OO_X(j))
 =H^{1,2}(X,\OO_X(j-i))=0.
 \]
 hypersurfaceの完全列と射影空間のcohomologyからこの消滅が従う。
 よって \cite[Theorem 4.5]{DDL} を適用できる。
 \(\OO_R(i)\) は明示的なflat liftである。
 実際 \(B_R(i)\) との有限 \(R\)-加群のテンソル積は完全で、
 \(\operatorname{Gr}B_R\to\QGr B_R\) は完全である。
 これらの対象はabelian圏も生成するので、必要なflat resolutionを取れる。

 閉埋め込み \(j:X\hookrightarrow\PP^{n+1}\) を考える。
 \(E_0=j^*(\bigoplus_{i=0}^{n+1}\OO_{\PP^{n+1}}(i))\) はcompactである。
 \(\RHom_X(E_0,M)=0\) なら、随伴とBeilinsonの生成定理
 \cite{Beilinson}から \(Rj_*M=0\)、従って \(M=0\) となる。
 したがって \(E_0\) はcompact generatorである。

 \(E_R\) は有界flat liftなので、
 \cite[Propositions 5.8, 5.9]{DDL} によりcompact generatorとなる。
 同論文Proposition 5.5と導来随伴は
 \[
 k\otL_R\RHom(E_R,E_R)
 \simeq\RHom_X(k\otL_RE_R,k\otL_RE_R)
 \]
 を与える。同論文Theorem 5.13の証明でも同じ底変更が使われている。
 dg enhancement上の導来還元関手からこの比較を作るため、
 比較は合成と単位を保つ。必要ならhomotopically \(R\)-flatなdg模型を選ぶ。
 これで \eqref{eq:derived-basechange} は複体だけでなくdg代数として成立する。
\end{proof}

\begin{corollary}\label{cor:smooth-proper}
 \(\cP_R=D(\QGr B_R)^c\) はsmooth properであり、
 \(\cP_R\otL_Rk\simeq\Perf(X)\) である。
\end{corollary}
\begin{proof}
 導来Morita理論 \cite{Keller} によって
 \(\cP_R\simeq\Perf(\Lambda_R)\) である。
 \(X\) のpropernessから \(\Lambda_0\) は \(k\)-perfectなので、
 補題\ref{lem:descent}で \(\Lambda_R\) は \(R\)-perfectとなる。
 次に
 \[
 A=\Lambda_R^{\mathrm{op}}\otL_R\Lambda_R,\qquad M=\Lambda_R
 \]
 とする。この導来ファイバーは \(\Lambda_0\) のenveloping algebraと対角加群である。
 \(X\) のsmoothnessにより対角加群はその上でperfectであり、
 同じ補題から \(\Lambda_R\) の対角もperfectとなる。
\end{proof}

\subsection{例外列、mutation、rotation}

\begin{proposition}\label{prop:relative-pair}
 \(\OO_R,\ldots,\OO_R(n-2)\) は \(\cP_R\) 内のrelative exceptional collectionである。
 その右直交 \(\cK_R\) はsmooth properで、導来特殊ファイバーは \(\Ku(X)\) である。
 \(\rho_R=\mathbf L_{\OO_R}\circ(1)\) は \(\cK_R\) の自己同値となる。
 これらの構成はArtin底変更と両立する。
\end{proposition}
\begin{proof}
 Hom複体の導来底変更と補題\ref{lem:descent}から
 \[
 \RHom(\OO_R(i),\OO_R(i))\simeq R,\qquad
 \RHom(\OO_R(j),\OO_R(i))=0\quad(j>i,\ 0\le i,j\le n-2)
 \]
 である。自己Homについては単位写像のconeを還元する。
 propernessにより評価・余評価はperfect対象を保つため、各例外成分はadmissibleである。
 評価のconeによるmutationを順に取って
 \[
 p_R=\mathbf L_{\OO_R}\cdots\mathbf L_{\OO_R(n-2)}
 \]
 と置くと、\(G_R=p_R(E_R)\) は \(\cK_R\) を生成する。
 mutationの導来底変更は \cite[Lemma 5.14]{DDL} による。
 \(G_R\) の導来特殊ファイバーは \(\Ku(X)\) の生成子なので、
 その自己準同型dg代数へ補題\ref{lem:descent}を適用すればsmooth proper性も従う。

 \(T=(1)\) は次数付き射影圏の自己同値で、\(T\OO_R(i)=\OO_R(i+1)\) である。
 \(M\in\cK_R\)、\(1\le i\le n-2\) に対して
 \[
 \RHom(\OO_R(i),TM)=\RHom(\OO_R(i-1),M)=0.
 \]
 \(\OO_R\) によるmutationはこの消滅を保ち、\(i=0\) の直交性も与える。
 よって \(\rho_R\) は \(\cK_R\) を保つ。
 その特殊ファイバーは通常の \(\rho_X\) である。
 Homの比較写像は還元後に同値なので、nilpotentな保守性から
 \(\rho_R\) は完全忠実である。

 \(\rho_R(G_R)\) の特殊ファイバーは生成子である。
 これに直交するInd対象を導来還元すると零になり、
 同じ保守性で元の対象も零になる。従って \(\rho_R(G_R)\) もcompact generatorである。
 冪等完備な圏からの完全忠実完全関手の像はthickなので、\(\rho_R\) は本質的全射でもある。

 底変更については、次数付き導来テンソル関手による生成子の自己Homの自然な比較を使う。
 剰余体上でこの比較は恒等写像となるため、変更後のArtin底上で同値となる。
 評価写像とそのconeも底変更と両立するので、残余圏とrotationまで両立する。
\end{proof}

\section{rotation付きの形式的復元}\label{sec:formal}

\subsection{正確な変形問題}

\begin{definition}
 \((\Ku(X),\rho_X)\) のmarked pair deformationとは、
 smooth properな \(R\)-linear dg圏 \(\cK_R\)、invertible bimoduleで表される
 自己同値 \(\rho_R\)、および特殊ファイバーでのenhanced pairの指定された同定からなる。
 同値にはこのmarkingとの両立を要求する。
 \(\rho_R^3\simeq[2]\) の指定されたcoherenceは定義に含めない。
\end{definition}

\begin{theorem}[pairの形式的復元]\label{thm:pair}
 \(X\) を滑らかなcubic \(n\)-fold、\(n\ge3\) とする。
 marked cubic deformationからmarked pair deformationへの写像は、
 各古典的局所Artin \(\CC\)-代数上で同型類の全単射を与える。
 互換的な形式族についても同型類の一意性と互換的な実現が成り立つ。
\end{theorem}

この主張はfull higher deformation groupoidの同値ではない。
pairには自然変換のscalar automorphismがあり、実現する圏同値も一意とは限らない。

\subsection{actual変形とHochschild形式変形}

\(\cC=\Ind(\Ku(X))\) とする。
actual compactly generated deformationの空間と、その形式moduli problemの間には
\[
 \theta_{\cC}:\operatorname{CatDef}_{\cC}^{c}
 \longrightarrow\operatorname{CatDef}_{\cC}^{\wedge}
\]
がある。右辺の接複体は \(CC^\bullet(\Ku(X))[2]\) である。
actual Morita変形が任意のHochschild classを実現するとは一般にはいえない
\cite{KL}。本稿ではその全射性を使わない。

\begin{lemma}\label{lem:theta}
 上の \(\theta_{\cC}\) は \((-1)\)-truncatedである。
 従って、右辺で同値となった二つのactual変形の同値はactual側へ戻せる。
 この性質はrotationによる \(\ZZ\)-homotopy fixed pointを取った後にも成り立つ。
\end{lemma}
\begin{proof}
 \(\Ku(X)\) はproperでcompact generatorを持つためtameである。
 \cite[Proposition 5.3.21, Theorem 5.1.9]{Lurie} を適用する。
 \cite[Definition 3.3, Section 3.3]{BKP} にも、
 \(\pi_0\) で単射、\(\pi_i\), \(i\ge1\), で同型という形でこの帰結が明記されている。
 比較写像は圏同値に自然であり、markingを移すrotationの作用と両立する。
 空間のmonomorphismはhomotopy limitで保たれる。
\end{proof}

一つのenhanced autoequivalenceを持ち上げる問題は、
markingに作用する \(\ZZ\) のhomotopy fixed pointである。
\(\ZZ\) は一つの可逆射で自由生成されるので、追加の有限群関係を課していない。
\[
 P=\left(\operatorname{CatDef}_{\cC}^{\wedge}\right)^{h\ZZ}
\]
と置くと、これはformal moduli problemであり、
\[
 T_P=R\Gamma(\ZZ,CC^\bullet(\Ku(X))[2]).
\]
\(\ZZ\) の自明加群の長さ一の自由分解から、次の完全列を得る。
\[
 0\longrightarrow
 \operatorname{coker}(1-\rho:\HH^1\to\HH^1)
 \longrightarrow H^0(T_P)
 \longrightarrow
 \ker(1-\rho:\HH^2\to\HH^2)
 \longrightarrow0.                                      \tag{1}\label{eq:tangent}
\]

\subsection{Hochschild計算とKodaira--Spencer写像}

\(N=n+2\) および
\[
 J(F)=\CC[x_1,\ldots,x_N]/(\partial_1F,\ldots,\partial_NF)
\]
を用いる。

\begin{lemma}\label{lem:HH}
 \(n\ge3\) のとき
 \[
 \HH^{<0}(\Ku(X))=0,\quad \HH^0(\Ku(X))=\CC,\quad
 \HH^1(\Ku(X))=0,
 \]
 \[
 \HH^2(\Ku(X))=
 \begin{cases}
 J(F)_3& n\ne4,\\
 J(F)_3\oplus\CC^2& n=4.
 \end{cases}
 \]
 幾何学的Kodaira--Spencer写像は同型
 \[
 H^1(X,T_X)\xrightarrow{\sim}
 \HH^2(\Ku(X))^{\rho_X}=J(F)_3
 \]
 を与える。その次元は \(\binom{n+2}{3}\) である。
\end{lemma}
\begin{proof}
 \cite[Theorem 5.39]{BFK} のKoszul cohomologyによる公式を
 三次のequivariant matrix factorizationへ適用する。
 恒等sectorは偶数次数 \(2r\) に \(J(F)_{3r}\) を与える。
 非自明な \(g\in\mu_3\) では固定点集合は原点だけであり、
 conormalのrankは \(N\)、determinant characterは \(-N\) である。
 cohomological degree \(q\) で残るweightは
 \[
 N+3(q-N)/2=(3q-N)/2.
 \]
 従って二つの非自明sectorは、\(3\mid N\) のときに限り、
 \(q=N/3\) に各一次元を与える。これで列記した群が得られる。

 BFKの恒等sectorのKoszul模型では、\(G\) の乗法が \([G]\) を表す。
 \(F+\varepsilon G\) の実際の変形類との一致は次の比較による。
 \(n\ge4\) では、periodicityから作るkernel morphismと
 Atiyah--Kodaira--Spencer morphismの一致を
 \cite[Lemma 2.19]{HR} とその前後の構成で確認できる。
 ここで三次に必要な次数条件 \(3\le(n+2)/2\) は満たされる。
 \(n=3\) では \cite[Theorem 4.4(2)]{JLLZ} が必要な接写像の同型を与える。
 normal sequenceとEuler sequenceは \(H^1(T_X)=J(F)_3\) を与える。

 幾何学的変形はrotationを持つので、この部分空間は不変である。
 \(n\ne4\) ではそれが \(\HH^2\) 全体である。
 \(n=4\) では \cite[Proposition 3.12]{Huy} により、
 rotationはMukai格子の \(A_2^\perp\) を固定し、\(A_2\) の三つの根を巡回置換する。
 invariantな2-Calabi--Yau volumeによる
 \(\HH^2\simeq\HH_0\) の下で、余分な二方向の固有値は
 \(\zeta,\zeta^2\) となり、不変部分はない。
 最後に偏微分は二次の正則列であるから
 \(H_{J(F)}(t)=(1+t)^N\) であり、次数三の次元は \(\binom N3\) である。
\end{proof}

\subsection{持ち上げの差と自己同型}

\begin{lemma}[Artin帰納法]\label{lem:induction}
 \(D\) をsmoothかつinfinitesimally homogeneousなmarked変形関手で、
 古典的Artin底上ではset-valuedとし、
 \(P\) をformal moduli problemとする。
 自然な写像 \(f:D\to P\) が接空間で同型を誘導すると仮定する。
 さらに、各small extension \(R'\to R\) と \(x'\in D(R')\) に対し、
 \(f(x'\vert_R)\) の任意の自己同型類が \(f(x')\) 上へ持ち上がると仮定する。
 このとき \(D(R)\to\pi_0P(R)\) は各古典的Artin底上で全単射である。
\end{lemma}
\begin{proof}
 Artin長について帰納法を行う。核 \(I\) のsmall extension上に
 \(y'\in P(R')\) があり、その還元を \(f(x)\) と同一視できたとする。
 \(D\) のsmoothnessにより \(x\) を \(x'\) に持ち上げる。
 固定した還元の同定を持つ \(f(x')\) と \(y'\) は、
 \(P(R')\to P(R)\) の同じhomotopy fibreに属する。
 このfibreが空でなければ、その連結成分は
 \(H^0(T_P)\otimes I\) のtorsorである。
 接写像の全射性によって \(x'\) を修正すれば差を消せる。

 単射性では、\(f(x')\simeq f(z')\) の還元を考える。
 帰納法により幾何学的還元は一致するが、
 圏同値の還元はその像の自己同型となり得る。
 仮定によってこの自己同型を上へ持ち上げ、その逆を合成する。
 すると還元を恒等に固定したfibreで二つのliftを比較できる。
 差の接ベクトルの像は零であり、接写像の単射性で元の差も零となる。
\end{proof}

\begin{lemma}\label{lem:automorphisms}
 cubicから得たpairのmarked deformationにおいて、
 small extensionに沿う上の自己同型持ち上げ条件が成り立つ。
\end{lemma}
\begin{proof}
 smooth properな圏ではHochschild複体はperfectで、その導来底変更が成り立つ。
 補題\ref{lem:HH}と局所Artin環上のminimal perfect complexを使うと
 \[
 \HH^0(\cK_R/R)=R,\qquad \HH^1(\cK_R/R)=0
 \]
 である。実際特殊ファイバーの負次数と次数一が零なので、
 minimal complexにはそれらの次数の項がなく、次数零は単位で生成される。

 恒等関手に還元されるinvertible bimoduleは相対対角と同型である。
 perfect bimoduleのliftの差はsmall extensionの核 \(I\) に対して
 \(\Ext^1(\Delta,\Delta)\otimes I=\HH^1(\Ku(X))\otimes I=0\)
 で測られ、対角自身がliftを与えるためである。
 この同型は各段階で指定した特殊ファイバーのmarkingを保つように選ぶ。
 従ってpairの自己同型の基礎となる圏同値は恒等と同型であり、
 残るのは \(\rho_R\Rightarrow\rho_R\) のscalar unitである。
 markingを固定すればこのscalarは \(1+\fm_R\) に入り、必ず \(R'\) へ持ち上がる。
 scalarな二次自己同型も持ち上がる。
 formal moduli problem側の自己同型をactual bimoduleで扱えることは、
 補題\ref{lem:theta}の \(\pi_1\) の同型による。
\end{proof}

\begin{proof}[定理\ref{thm:pair}の証明]
 cubicのmarked変形関手はsmoothである。方程式の係数を持ち上げればよく、
 \(H^1(\OO_X)=H^2(\OO_X)=0\) によって \(\OO_X(1)\) は一意に延長する。
 infinitesimalな射影自己同型もない。実際、二次の偏微分の正則列には
 一次係数の非自明なsyzygyがなく、Euler scalarを除けば \(H^0(T_X)=0\) である。
 これはmarkedな幾何学的変形をset-valuedとして扱うことを正当化する。

 式\eqref{eq:tangent}と補題\ref{lem:HH}により接写像は同型である。
 補題\ref{lem:automorphisms}と補題\ref{lem:induction}を適用し、
 最後に補題\ref{lem:theta}でactual pairの同値へ戻す。
 帰納法では各段階で指定した還元の同値を保てるので、
 \(\CC[\hbar]/(\hbar^{r+1})\) に沿った実現も互換的に選べる。
\end{proof}

\begin{proof}[定理\ref{thm:main}と系\ref{cor:formal}の証明]
 命題\ref{prop:relative-pair}の組へ定理\ref{thm:pair}を適用する。
 互換的な幾何学的liftの方程式の係数の逆極限が \(F_\hbar\) を与える。
 この最後の操作は形式族を作るものであり、非零複素数で評価する操作ではない。
\end{proof}

\begin{corollary}\label{cor:period}
 上で構成したArtin pairについて、指定を伴わない同型
 \(\rho_R^3\simeq[2]\) が存在する。
\end{corollary}
\begin{proof}
 \(\rho_R^3[-2]\) は恒等に還元されるinvertible bimoduleである。
 補題\ref{lem:automorphisms}で使った \(\HH^1\)-rigidityを適用する。
 または可換モデルからこの同型を移してもよい。
 特定のmatrix factorization表示が与えるperiodicity同型とのcoherenceの一致は、
 この系の主張に含めない。
\end{proof}

\begin{remark}
 sevenfoldでは \(\HH^3(\Ku(X))\simeq\CC^2\) があるが、
 上の帰納法は \(\HH^3=0\) を要求しない。
 比較している二つのliftは既に存在し、その差を接空間で修正する。
 higher operationsを無視したformalityの仮定も使っていない。
\end{remark}

\subsection{四次元以外ではrotationを忘れた形式復元}\label{sec:unpolarized-formal}

四次元以外では、形式近傍に限れば、復元の入力からrotationを省くことができる。
ここでいうmarked dg変形は、特殊ファイバーを \(\Ku(X)\) と同定した
smooth properな \(R\)-linear dg圏であり、同値にもこのmarkingとの両立を要求する。

\begin{theorem}[rotationを忘れた形式復元]\label{thm:unpolarized-formal}
 \(X\subset\PP^{n+1}_{\CC}\) を滑らかなcubicとし、\(n\ge3\)、\(n\ne4\) とする。
 任意の古典的局所Artin \(\CC\)-代数 \(R\) について、
 marked cubic変形から \(\Ku(X)\) のmarked smooth proper dg変形への写像は、
 同型類の全単射を与える。従って、任意のそのような \(\cK_R\) は
 一意なmarked cubic変形 \(X_R\) の \(\Ku(X_R)\) と同値である。

 また、固定した \(\cK_R\) 上で、\(\rho_X\) を持ち上げるinvertible bimodule
 \(\rho_R\) が存在し、そのmarked同型類は一意である。
 従ってrotationを忘れる写像
 \[
  \{\text{marked smooth proper pair変形}\}/\simeq
  \ \longrightarrow\
  \{\text{marked smooth proper dg変形}\}/\simeq
 \]
 も全単射である。これらの主張はArtin底変更と互換的な形式族についても成り立つ。
\end{theorem}
\begin{proof}
 \(Q=\operatorname{CatDef}_{\Ind\Ku(X)}^{\wedge}\) と置く。
 補題\ref{lem:HH}の幾何学的Kodaira--Spencer写像は、\(n\ne4\) では
 \[
  H^1(X,T_X)\xrightarrow{\sim}
  \HH^2(\Ku(X))=H^0(T_Q)
 \]
 である。cubicのmarked変形関手はsmoothかつset-valuedである。
 補題\ref{lem:induction}を適用するため、幾何学的な変形の像の自己同型を調べる。

 \(\Ku(X_R)\) の恒等に還元される自己同値を表すinvertible bimoduleは、
 \(\HH^1(\Ku(X))=0\) により相対対角bimoduleと同型である。
 これは補題\ref{lem:automorphisms}と同じsmall extensionの帰納法で従う。
 特殊ファイバー上の同型を指定しても、そこでのscalarは \(R^\times\) へ持ち上がるので、
 同型をmarkingと両立するように選べる。
 従ってmarkedな自己同値の同型類は恒等だけであり、
 必要な自己同型持ち上げ条件は満たされる。
 \(Q\) 側の自己同型についても、補題\ref{lem:theta}の比較写像が
 \(\pi_1\) 上の同型を誘導するため同じことが成り立つ。

 補題\ref{lem:induction}から、cubic変形と \(\pi_0Q(R)\) の全単射が得られる。
 幾何学的変形の像はactual smooth proper変形であり、
 補題\ref{lem:theta}によりactual変形間の同値を検出できる。
 よって最初の主張が従う。この議論はHochschild形式変形の一般的なactual実現性を
 仮定せず、この場合には幾何学的なliftで全てを実現している。

 得られた同値 \(\Ku(X_R)\simeq\cK_R\) を通じて
 \(\rho_{X_R}\) を移せば、\(\rho_X\) のliftが存在する。
 二つのmarkedなlift \(\rho_R,\rho'_R\) があれば、
 \(\rho'_R\rho_R^{-1}\) は恒等に還元されるinvertible bimoduleである。
 再び \(\HH^1=0\) による対角のrigidityを用いると、
 \(\rho'_R\rho_R^{-1}\simeq\mathrm{id}\) となり、markingを保つ同型も選べる。
 これがliftの同型類の一意性とforgetful mapの単射性を与える。
 各small extensionで還元上の同値を保って構成できるため、
 形式族での互換性も従う。
\end{proof}

\begin{remark}\label{rem:unpolarized-scope}
 この定理の一意性は同型類についてのものである。
 rotationとのintertwinerにはscalar自由度があり、
 rotationを忘れる写像がfull higher deformation groupoidの同値だとは主張しない。
 実際、\(\HH^0=\CC\)、\(\HH^1=0\) とrotationのscalar上の自明作用から、
 pairの接複体には \(H^{-1}(T_P)\simeq\CC\) がある一方、
 \(H^{-1}(T_Q)=0\) である。

 \(n=4\) では幾何学的な20方向に加えて \(\HH^2\) に二方向があり、
 rotationの固定部分を取って初めて接写像が同型となる。
 この次元でrotationを省く結論を上の証明から導くことはできない。
 ここでは、Hochschildの全22方向が各Artin底上でactual変形として
 実現されるとは主張していない。

 \(n=7\) では \(\HH^3\simeq\CC^2\) が残るが、
 幾何学的なliftと既に存在するliftを比較する帰納法により結論は成り立つ。
 障害を収め得るcohomology群が非零であることと、
 この変形問題で実際に障害が現れることは同じではない。
 一般の有限パラメータの圏からrotationを省けるという大域的主張も含めない。
\end{remark}

\section{飽和な代数的族に対するétale局所復元}\label{sec:algebraization}

形式的な復元を有限パラメータへ移す際には、圏の族そのものの代数性が必要である。
本節ではそれを明示的に仮定し、可換モデルと圏同値の両方をétale近傍上で構成する。
各次数がflatであることだけから、この仮定が自動的に従うとは主張しない。

\begin{definition}\label{def:algebraic-family}
 \(S\) を有限型複素スキームとする。
 本節で\emph{smooth properな代数的dg圏の族}とは、
 étale局所的に \(\Perf(A)\) と書け、\(A\) が底環上のdg代数で、
 \(A\) 自身が底上perfect、対角 \(A\) が \(A^{\mathrm{op}}\otL A\) 上perfectである族をいう。
 圏は小さく、pretriangulatedかつ冪等完備とし、すべての底変更を導来的に取る。
 自己同値 \(\rho\) の代数的族は、可逆なperfect bimoduleによって与える。
\end{definition}

\begin{theorem}[étale局所的な可換復元]\label{thm:etale}
 \(S\) を有限型複素スキーム、\(s\in S(\CC)\) とし、
 \((\cK_S,\rho_S)\) を定義\ref{def:algebraic-family}の組とする。
 滑らかなcubic \(X^n\subset\PP^{n+1}\)、\(n\ge3\)、およびenhanced pair同値
 \[
 \phi_s:(\cK_s,\rho_s)\xrightarrow{\sim}(\Ku(X),\rho_X)
 \]
 が与えられたとする。このとき、あるpointed étale近傍
 \((T,t)\to(S,s)\)、滑らかな相対cubic \(Y\subset\PP^{n+1}_T\)、および
 \[
 \Phi_T:(\cK_T,\rho_T)\xrightarrow{\sim}(\Ku(Y/T),\rho_{Y/T})
 \]
 が存在し、\(t\) 上で \(Y_t=X\)、\(\Phi_t\simeq\phi_s\) とできる。
 この同値はperfect bimoduleとrotationを結ぶ同型によって与えられる。

 従って、可換cubicのpairで実現できる複素点は \(S\) のZariski開部分をなす。
 その各複素点で、実現する滑らかなcubicの射影同型類は一意である。
\end{theorem}

この定理は、底 \(S\) 全体上のcubicの一意な降下や、同値の一意性を述べない。
また、\(S\) が連結であることだけでは、上の開部分が \(S\) 全体であるとはいえない。
\(n\ne4\) ではrotationを最初から与えなくても、同じ方法でétale局所的な
可換モデルが存在する（系\ref{cor:etale-unpolarized}）。

\subsection{指定された還元同値を保つ持ち上げ}

\begin{lemma}\label{lem:relative-realization}
 \(R'\twoheadrightarrow R\) を剰余体 \(\CC\) の局所Artin代数のsmall extensionとする。
 \((\cK_{R'},\rho_{R'})\) をsmooth proper pairとし、滑らかなcubic \(Y_R\) と同値
 \[
 \phi_R:(\cK_R,\rho_R)\simeq(\Ku(Y_R),\rho_{Y_R})
 \]
 を固定する。このとき \(Y_R\) を延長する \(Y_{R'}\) と、
 還元後に \(\phi_R\) と同型になるpair同値 \(\phi_{R'}\) が存在する。
 \(Y_R\) の射影空間への埋込みも指定してよい。
\end{lemma}
\begin{proof}
 特殊ファイバーでの \(\phi_R\) をmarkingとして固定する。
 まずcubicの係数を持ち上げ、幾何側に一つのlift \(Y^0_{R'}\) を取る。
 \(P\) を第\ref{sec:formal}節のpairのformal moduli problemとし、
 \(I=\ker(R'\to R)\) と置く。
 与えられた圏とこの幾何学的liftを、\(\phi_R\) によって
 \(P(R')\to P(R)\) の同じhomotopy fibre内で比較する。
 存在するliftの差は \(H^0(T_P)\otimes I\) のtorsorで測られる。
 幾何学的Kodaira--Spencer写像
 \(H^1(T_{Y_0})\simeq H^0(T_P)\) により \(Y^0_{R'}\) を修正し、その差を消す。
 補題\ref{lem:theta}の完全忠実性で、指定した還元同値を保つactual同値へ戻る。
 途中で還元上の自己同型が現れる表示を用いた場合には、
 補題\ref{lem:automorphisms}でそれを持ち上げて補正できる。

 抽象的な偏極cubicのliftに対しても、\(\OO(1)\) の延長と
 \(H^0(\OO(1))\) の自由な基底の持ち上げにより、指定された埋込みを延長できる。
 これは \(\operatorname{PGL}_{n+2}\) のsmoothnessを用いることと同じである。
\end{proof}

\subsection{同値とrotationの同型を表す代数的スタック}

\(U\subset\PP(\CC[x_1,\ldots,x_{n+2}]_3)\) を滑らかなcubicの方程式空間とし、
 \(\mathcal X\to U\) を普遍cubicとする。
 相対的な例外列とmutationにより
 \((\Ku(\mathcal X/U),\rho_U)\) はsmooth properな代数的pairの族となる。
 相対生成子として \(\bigoplus_{i=0}^{n+1}\OO_{\mathcal X}(i)\) の残余射影を使える。
 この生成性は、射影空間の相対Beilinson生成子の引戻しと、
 閉埋込みの押出しが保守的であることから従う。

\begin{proposition}\label{prop:equivalence-stack}
 \(S\)-スキーム \(T\) に対し、
 \[
 (f:T\to U,\ P,\ \alpha:P\circ\rho_T\xrightarrow{\sim}\rho_f\circ P)
 \]
 を対象とするスタック \(\mathcal E\) を考える。
 ここで \(P:\cK_T\simeq\Ku(\mathcal X_f/T)\) はperfect bimoduleで表される同値、
 射は \(\alpha\) と両立するbimoduleの同型である。
 このスタックは \(S\times U\) 上、従って \(S\) 上で局所有限表示のArtin 1-stackである。
\end{proposition}
\begin{proof}
 すべての構成はétale降下と両立するため、底をaffineにして生成子を選んでよい。
 \(\cD=\cK_S^{\mathrm{op}}\otL\Ku(\mathcal X/U)\) は \(S\times U\) 上smooth properである。
 Toën--VaquiéのTheorem 3.6とCorollary 3.15 \cite{TV}により、
 そのperfect bimoduleのmoduliは局所幾何的で局所有限表示である。
 同論文は一般の可換底環上で定式化されているので、定数族だけに限定されない。

 \(P\) が同値を与える条件は開である。
 実際、右随伴を表すperfect bimoduleを \(P^R\) とすると、
 単位と余単位
 \[
 \Delta_{\cK}\longrightarrow P^R P,
 \qquad PP^R\longrightarrow\Delta_{\Ku}
 \]
 のconeがともに零となることが必要十分である。
 生成子でdg代数を選んだ模型で、これらのconeを評価する。
 smooth proper性によりその底上の複体はperfectであり、消滅はそこで検出される。
 従って消滅座は開である。
 これは \cite[Lemma 3.25とCorollary 3.24の証明]{TV} の二つの圏に対する同じ議論である。

 同値 \(P\) の幾何ファイバーで
 \[
 \Ext^i(P,P)=\HH^i(\Ku(X_f)),\qquad i\in\ZZ.
 \]
 補題\ref{lem:HH}から負次数は消える。
 perfectな自己Hom複体のTor振幅の開性により、これは任意の底変更後にも負次数が
 消える開部分として扱える。\cite[Corollary 3.21]{TV}により、
 同値を表す部分は通常のArtin 1-stackとなる。

 最後に \(P\rho_T\) と \(\rho_fP\) の間の同型を付加する。
 射のmoduliの局所有限表示性は \cite[Corollary 3.19]{TV} による。
 その中でconeが零となる部分を取れば同型のmoduliを得る。
 同型が存在する点では両側は可逆bimoduleであり、
 負次数のHomは再び \(\HH^{<0}\) の消滅によって零である。
 従ってこの同型の選択は代数的空間で表され、
 \(\mathcal E\) はArtin 1-stackのままである。
 この構成には \(\rho^3\simeq[2]\) の選択や、そのcoherenceは課していない。
\end{proof}

\subsection{smoothnessとétale局所section}

\begin{proof}[定理\ref{thm:etale}の証明]
 \(\phi_s\) と \(X\) の方程式を選んで \(e\in\mathcal E(\CC)\) を得る。
 \(V\to\mathcal E\) をsmooth atlasとし、\(e\) の上に複素点 \(v\) を選ぶ。
 \(v\) における \(V\to S\) のsmall extensionの持ち上げ問題を考える。
 補題\ref{lem:relative-realization}により、対応する \(\mathcal E\) の対象を
 指定された還元を保って持ち上げられる。
 さらに \(V\to\mathcal E\) のsmoothnessにより、それを \(V\) へ持ち上げる。
 底が局所Artin環で剰余体が \(\CC\) なので、atlasへのétale局所持ち上げはsectionを持ち、
 この持ち上げを同じ底上で取れる。

 \(V\to S\) は局所有限型であるから、Noetherian底上の無限小判定
 \cite[Tag 02HX]{Stacks}により \(v\) でsmoothである。
 \(v\) の近傍を縮め、smooth morphismのétale局所section
 \cite[Tag 055Uとその証明中のLemma 37.38.5]{Stacks}を使うと、
 \(\kappa(v)=\kappa(s)=\CC\) であることから、指定した \(v\) を通る
 \((T,t)\to(S,s)\) と \(T\to V\to\mathcal E\) が得られる。
 これが求める \(Y\)、\(P\)、\(\alpha\) を同時に与える。

 \(T\to S\) の像はZariski開であり、その各複素点の上には複素点を取れる。
 従って可換モデルが存在する複素点の集合は開部分の複素点集合である。
 各点でのモデルの一意性は \cite[Theorem 1.7]{LZ} による。
\end{proof}

\begin{corollary}\label{cor:generic-realization}
 定理\ref{thm:etale}の仮定で \(S\) が既約なら、\(s\) を含む空でないZariski開部分
 \(S^\circ\) の全複素ファイバーに可換cubicモデルが存在する。
 幾何学的generic fibreにも、有限型étale底変更を通じてモデルが存在する。
\end{corollary}
\begin{proof}
 上のétale近傍の像を \(S^\circ\) とする。既約性からこれはdenseである。
 \(S\) を必要なら被約化してgeneric pointを考えると、
 \(T\) の支配的な成分のgeneric residue fieldは \(\CC(S)\) の有限分離拡大である。
 その上のcubicと同値を代数閉包へ底変更すればよい。
\end{proof}

\begin{remark}\label{rem:algebraization-boundary}
 本節では形式的なpower seriesの収束を証明していない。
 指定された代数的圏の族に対し、代数的な同値のmoduliのsmoothnessから
 別途、実際の代数的近傍と同値を作っている。
 その完備化は形式的復元定理のモデルとmarked同型類で一致する。
 しかし、任意に選ばれた形式的座標表示や形式的bimoduleそのものが
 元の底上へ一意に降下する、という結論ではない。
\end{remark}


\begin{corollary}[rotationを指定しない局所的存在]\label{cor:etale-unpolarized}
 定理\ref{thm:etale}で \(n\ne4\) とする。
 \(\rho_S\) は仮定せず、\(\cK_s\simeq\Ku(X)\) というenhanced同値だけを与えても、
 pointed étale近傍 \(T\) 上にcubic \(Y/T\) と、指定された特殊ファイバーの同値を
 延長する \(\cK_T\simeq\Ku(Y/T)\) が存在する。
\end{corollary}
\begin{proof}
 命題\ref{prop:equivalence-stack}から \(\alpha\) の指定を除いた同値のスタックを用いる。
 定理\ref{thm:unpolarized-formal}の証明を、還元同値を指定したhomotopy fibre内で行えば、
 補題\ref{lem:relative-realization}のrotationなしの版が得られる。
 後は定理\ref{thm:etale}と同じsmooth atlasとétale局所sectionの議論を適用する。
 得られたモデルからrotationを移せるが、その \(T\) 全体上での一意性は主張しない。
\end{proof}

\section{Sklyaninブロックと次数シフト}\label{sec:elliptic}

\subsection{中心元と点曲線}

\begin{lemma}\label{lem:center}
 \(G_h=x^3+y^3+z^3-2h^2\sym(xyz)\) は、すべての \(h\) で
 \(S_h\) の非零な中心正則元である。
\end{lemma}
\begin{proof}
 \(D=x^3+y^3+z^3\) とする。関係式から
 \[
 [D,x]=h(zy^2z-yz^2y)=h[zy,yz]=h^2[zy,x^2].
 \]
 また六つの順列を同じ順序へ移すと
 \[
 \sym(xyz)=xyz+\frac h2(-x^3+y^3-z^3)
 \]
 となる。この式を \(x\) と交換させ、再び二次関係式を使えば
 \(2[\sym(xyz),x]=[zy,x^2]\) を得る。
 従って \([G_h,x]=0\)。巡回置換で \(y,z\) についても従う。

 非零性は特殊化で失われない形で確認できる。
 自由代数の次数三部分の線形汎関数 \(\ell\) を
 \[
 \ell(x^3)=\ell(y^3)=\ell(z^3)=1,\quad
 \ell(xyz)=\ell(yzx)=\ell(zxy)=h/2,
 \]
 \[
 \ell(xzy)=\ell(zyx)=\ell(yxz)=-h/2
 \]
 とし、ちょうど二文字が等しい語では零とする。
 三つの二次関係式の左右に一文字を掛けた18個の式を \(\ell\) はすべて殺す。
 例えば \(\ell(z(xy-yx-hz^2))=h/2+h/2-h=0\) である。
 従って \(\ell\) は \((S_h)_3\) に降り、\(\ell(G_h)=3\) となる。
 \((a,b,c)=(1,-1,-h)\) はSklyanin代数の12個の非正則パラメータに属さず、
 \(S_h\) はAS-regular domainである \cite[Introduction]{CGWZ}。
 \(h=0\) は多項式環である。よって \(G_h\) は左右正則である。
\end{proof}

標準Sklyanin記法では \((a,b,c)=(1,-1,-h)\) である。
標準中心元の式 \cite[Equation (2.16)]{Walton}
\[
 g=c(a^3-c^3)x^3+a(b^3-c^3)xyz+b(c^3-a^3)yxz+c(c^3-b^3)y^3
\]
に代入し、\(yxz=xyz-hz^3\) を用いると \(g=-hG_h\) となる。
\(h\ne0\)、\(h^6\ne-27\) では点曲線の公式
\cite[Theorem 2.14]{Walton}から
\[
 E_h=V(x^3+y^3+z^3+h^2xyz)\subset\PP^2
\]
を得る。これは滑らかな楕円曲線であり、
\[
 S_h/(G_h)\simeq B(E_h,L,\sigma_h),\qquad L=\OO_{E_h}(1),
\]
と書ける。ここで \(\sigma_h=t_p\) は平行移動である。
原点をinflection point \(o=[1:-1:0]\) に取ると、
多重線形化した関係式を第一成分 \(o\) で評価した核は
\(p=[1:-1:-h]\) である。従ってこの規約で \(\sigma_h=t_p\) となる。
この中心商と楕円曲線の比較は \cite[Sections 3.5--3.6]{EG} も参照。
\(h=0\) は通常のFermat平面三次曲線の斉次座標環である。
この場合は \(p=o\)、\(\sigma_0=\Id\) とする。

\subsection{rotationに含まれる平行移動を消す}

\begin{lemma}\label{lem:elliptic-conjugacy}
 \(E\) を滑らかな楕円曲線、\(L\) を次数三の直線束、
 \(\sigma=t_p\) とする。
 \(C=B(E,L,\sigma)\) のgraded singularity categoryを \(D^b(E)\) と同一視したとき、
 その次数シフトは可換平面三次曲線のrotation
 \(\rho_0=T_{\OO_E}\circ(-\otimes L)\) と共役である。
\end{lemma}
\begin{proof}
 右加群の規約を
 \[
 \Gamma_*(F)_d=H^0(E,F\otimes L_d),\qquad
 L_d=L\otimes\sigma^*L\otimes\cdots\otimes(\sigma^{d-1})^*L
 \]
 と取る。tailsの次数シフトは \(H(F)=\sigma_*(F\otimes L)\) である。
 Gorenstein parameter零のOrlov windowを比較すると
 \[
 \rho_{\mathrm{nc}}=T_{\OO_E}\circ H=t_{-p}^*\rho_0.
 \tag{2}\label{eq:rho-nc}
 \]
 このwindow公式を確認する。特異圏への商を \(q\) とし
 \(\Psi_i=qR\Gamma_{\ge i}\) と置くと、
 \((1)\Psi_1=\Psi_0H\)、\(q(k)=\Psi_0(\OO_E)[1]\)、
 \[
 \Psi_1F\longrightarrow\Psi_0F
 \longrightarrow\RHom(\OO_E,F)\otimes q(k)
 \]
 がある。楕円曲線のSerre双対から
 \(\Psi_1=\Psi_0T_{\OO_E}^{-1}\) となり、\eqref{eq:rho-nc}を得る。

 原点を \(o\)、\(P_q=\OO_E(q-o)\)、\(A_q=(-\otimes P_q)\) とする。
 spherical twistの核は \(I_\Delta[1]\) なので、
 \[
 K_{T_{P_q}}=(P_q^{-1}\boxtimes P_q)\otimes I_\Delta[1].
 \]
 see-sawを両方向のファイバーへ適用すると
 \[
 (P_q^{-1}\boxtimes P_q)\otimes I_\Delta
 \simeq I_{\{y=x-q\}}.
 \]
 よって \(T_{P_q}=t_q^*T_{\OO_E}\) であり、
 \[
 A_q\rho_0A_q^{-1}=t_q^*\rho_0,\qquad
 \rho_{\mathrm{nc}}=A_{-p}\rho_0A_{-p}^{-1}.
 \tag{3}\label{eq:conjugacy}
 \]
 これらは核とwindowによるdg同値として実現される。
 三次周期性の二つの同型の差は \(\HH^0(D^b(E))^\times=\CC^\times\) のscalarであり、
 共役の自然同型をその三乗根で調整すれば周期性も両立させられる。
 \(\HH^{<0}(D^b(E))=0\) なので、この自然同型の等式に追加の高次ホモトピー障害はない。
 後で用いる有限次数差の作用は、この周期性から構成する。
\end{proof}

\begin{corollary}\label{cor:block}
 \(h^6\ne-27\) のとき、三次周期性と次数シフトを保つdg Morita同値
 \[
 \MF^{\ZZ}(S_h,G_h)
 \simeq
 \MF^{\ZZ}\bigl(\CC[x,y,z],x^3+y^3+z^3+h^2xyz\bigr)
 \]
 が存在する。
\end{corollary}
\begin{proof}
 \cite[Theorem 3.6.2]{EG}、Orlov window、
 \cite[Theorem 5.8]{CCKM}によって両辺を \(D^b(E_h)\) と比較し、
 補題\ref{lem:elliptic-conjugacy}を使う。
 MFと特異圏の比較は、factorizationを商環上の両側無限の完全周期分解へ送る
 dg関手で実現する。中心正則元の商はGorensteinであり、complete resolutionは
 MCM加群の安定圏のdg模型を与える。
 CCKMの完全忠実性をすべてのshiftに適用すると、Hom複体の比較がquasi-isomorphismとなる。
 周囲がAS-regularであることによる高次syzygyのfactorization表示が本質的全射性を与える。
 特異圏からwindowへのdg比較は \cite[Section 5, Theorem 5.8]{CT} による。
 \(h=0\) では恒等の可換比較を用いる。
\end{proof}

\section{中心三次式を加える比較と通常次数への降下}\label{sec:TS}

\subsection{二つの次数ブロックと外部積}

中心変数側の三次式をFermat型に制限する必要はない。
この節では \(T=\CC[u_1,\ldots,u_m]\)、\(m\ge1\)、とし、
\(P\in T_3\) は
\[
 J(P)=\left(\frac{\partial P}{\partial u_1},\ldots,
                 \frac{\partial P}{\partial u_m}\right),\qquad
 \dim_\CC T/J(P)<\infty                                      \tag{4a}\label{eq:central-isolated}
\]
を満たすものとする。これは斉次三次式 \(P\) のアフィン臨界点が
原点だけであるという条件である。
\[
 \Gamma=(\ZZ e_0\oplus\ZZ e_1)/\langle3e_0-3e_1\rangle,
 \qquad w=3e_0=3e_1                                       \tag{4b}\label{eq:central-gamma}
\]
と置く。\(S_d\) の次数を \(de_0\)、すべての \(u_i\) の次数を \(e_1\) とする。
したがって中心変数どうしの混合項を含む \(P\) も次数 \(w\) である。

\begin{lemma}[任意の孤立中心三次式に対する外部積]\label{lem:TS}
 \(S\) を連結非負次数付きnoetherian AS-regular \(\CC\)-代数、
 \(g\in S_3\) を中心正則元とする。
 \eqref{eq:central-isolated}の下で、外部積はpretriangulated hullと
 冪等完備化の後にdg Morita同値
 \[
 \MF^\ZZ(S,g)\otL_\CC\MF^\ZZ(T,P)
 \simeq
 \MF^\Gamma(S\otimes_\CC T,g+P)                             \tag{4c}\label{eq:central-TS}
 \]
 を与える。この同値は二つの因子の次数シフトと、
 両者の三次内部シフトを \([2]\) と同一視する周期性を保つ。
\end{lemma}
\begin{proof}
 \(R=S\otimes_\CC T=S[u_1,\ldots,u_m]\)、\(W=g+P\)、
 \(A=R/(W)\)、\(C=S/(g)\) と置く。
 まず \(W\) は中心正則元である。
 実際、\(u\) に関する単項式順序を固定すると、\(P\ne0\) の最高項の係数は
 非零のscalarであり、任意の非零多項式との積の最高項は消えない。
 以下では中心正則元の商に対するMFとgraded singularity categoryの比較
 \cite[Theorem 5.8]{CCKM}を用いる。
 \(\Gamma\to\ZZ\)、\(e_0,e_1\mapsto1\)、は有限の核を持ち、
 正の整数次数を与えるので、有限次数付き自由分解による同じ比較は
 \(\Gamma\)-次数についても適用できる。

 \paragraph{外部積と完全忠実性.}
 factorizationの規約を
 \[
 Q=(Q_1\xrightarrow{\varphi}Q_0
                  \xrightarrow{\psi}Q_1(w)),\qquad
 \psi\varphi=g\Id,\quad\varphi\psi=g\Id
 \]
 とする。これを \(Q^{i+2}=Q^i(w)\) を満たす曲がった複体に展開すると、
 \(Q^0=Q_0\)、\(Q^1=Q_1(w)\) である。
 二つのfactorizationの外部積は、二つの奇微分をKoszul符号付きで足すことにより定める。
 奇微分の外部和の平方は \(g+P\) である。
 二つの \([2]\) を同じ内部シフトに対応させる関係が
 \(3e_0=3e_1\) であり、得られるfactorizationは有限自由である。

 有限自由対象 \(Q,Q'\) と \(K,K'\) に対して、
 係数を \(S\)-次数と \(T\)-次数に分解すると
 \[
 \Hom_{\MF^\Gamma}(Q\boxtimes K,Q'\boxtimes K')
 \simeq
 \operatorname{Tot}\!\left(
   \Hom_{\MF^\ZZ(S,g)}(Q,Q')\otimes_\CC
   \Hom_{\MF^\ZZ(T,P)}(K,K')\right)                         \tag{4d}\label{eq:central-hom}
 \]
 を得る。具体的には、左辺の次数 \(r\) の斉次写像は、
 \(a+b=r\) を満たす右辺の次数 \((a,b)\) の写像の和に一意に分解する。
 二つの内部次数の差を移す操作は \((3,-3)\) の整数倍だけであり、
 これは二つのcohomological degreeを \((2,-2)\) だけ変える操作に対応する。
 微分と合成の符号も外部和のKoszul符号に一致する。
 各因子のHom複体は、有限個の自由加群の次数シフトと係数代数の非負次数性から
 下に有界である。従って各全cohomological degreeに寄与する項は有限個であり、
 \eqref{eq:central-hom}に完備化は必要ない。
 これが外部積のdg完全忠実性を与える。

 \paragraph{中心変数を零とした商への生成性の帰着.}
 \(\MF^\Gamma(R,W)\) の有限自由対象の奇微分を \(D\) とする。
 \(u_i\) は中心変数なので、係数行列の \(u_i\)-微分は \(R\) 上の導分であり、
 \[
 D\,\frac{\partial D}{\partial u_i}
 +\frac{\partial D}{\partial u_i}\,D
 =\frac{\partial P}{\partial u_i}\Id
 \]
 が成り立つ。内部シフトを付けると、これは
 \(\partial_iP:N\to N(2e_1)\) の零ホモトピーである。
 斉次イデアル \(J(P)\) は \((u_1,\ldots,u_m)\)-primaryなので、
 ある整数 \(d\ge2\) に対してすべての \(u_i^d\) が \(J(P)\) に属する。
 斉次係数でこれらを偏微分の線形結合に書けば、
 \(u_i^d:N\to N(de_1)\) も特異圏で零となる。

 \(N\) を有限生成 \(A\)-加群の有界複体で表し、
 \(u_1^d,\ldots,u_m^d\) に沿うKoszul複体を \(K(u^d;N)\) と書く。
 零写像のconeは元の対象を直和因子に持つため、\(N\) は
 \(K(u^d;N)\) の直和因子である。
 より正確には、一変数ごとのKoszul cone関手は完全関手でretractを保つので、
 各 \(u_i^d\) の零作用から得られるretractを順に合成すればよい。
 このKoszul複体のcohomologyは \((u_1^d,\ldots,u_m^d)\) に殺され、
 \[
 (u_1,\ldots,u_m)^{m(d-1)+1}\subset(u_1^d,\ldots,u_m^d)
 \]
 である。従って各cohomology加群は、有限個の \(C=A/(u_1,\ldots,u_m)\)-加群を
 商とする有限filtrationを持つ。標準切断の三角形も使うと、
 右辺の特異圏は \(C\)-加群の制限とその次数シフトで厚く生成される。

 ここで制限が特異圏へ降りることを確認する。
 \(u_1,\ldots,u_m\) は \(A\) の中心正則列である。
 一つの変数 \(u_i\) について、\(R\) で \(u_i a=Wb\) とする。
 \(u_i=0\) とした商では、\(\overline W=g+P|_{u_i=0}\) は正則である。
 特殊化した \(P\) が非零なら最高項のscalar係数を使い、
 零なら \(g\) の正則性を使えばよい。
 従って \(\overline b=0\)、すなわち \(b=u_i b_1\) であり、
 多項式環 \(R\) で \(u_i\) を消去して \(a=Wb_1\) を得る。
 先行する変数を順に零とした同じ議論が正則列を与える。
 よって \(C\) は有限自由Koszul分解を持つ \(A\)-perfect加群であり、
 制限は \(\Dsg(C)\to\Dsg(A)\) を誘導する。

 \paragraph{制限加群の安定化の明示.}
 Euler公式から
 \[
 q_i=\frac13\frac{\partial P}{\partial u_i},\qquad
 P=\sum_{i=1}^{m}u_iq_i
 \]
 と書ける。\(T\) 上の通常のKoszul分解を
 \[
 K_j=T(-je_1)\otimes_\CC\Lambda^j\CC^m,
 \qquad d_K=\sum_i u_i\iota_i
 \]
 とする。\(\iota_i\) は外積代数の縮約、\(\varepsilon_i\) は対応する基底による左外積である。
 \[
 h_K=\sum_iq_i\varepsilon_i:K_j\longrightarrow K_{j+1}(w)
 \]
 と置くと
 \[
 d_Kh_K+h_Kd_K=P\Id,\qquad h_K^2=0.                         \tag{4e}\label{eq:central-koszul}
 \]
 ここで \(q_i\) の次数は \(2e_1\) であり、
 \(K_{j+1}(w)\) のシフトによりすべての表示写像は内部次数零である。

 有限Koszul複体を偶奇で折り畳み、
 \[
 (K_P)_0=\bigoplus_{j\ge0}K_{2j}(jw),\qquad
 (K_P)_1=\bigoplus_{j\ge0}K_{2j+1}(jw)
 \]
 と置く。\(d_K+h_K\) は
 \[
 (K_P)_1\longrightarrow(K_P)_0
             \longrightarrow(K_P)_1(w)
 \]
 を与え、\eqref{eq:central-koszul}によってこれは \(P\) の有限自由factorizationである。

 \(M\) を \(C\)-加群とし、その特異圏での像を表すfactorizationを
 \(Q=(Q_1\xrightarrow\varphi Q_0\xrightarrow\psi Q_1(w))\) とする。
 有限回のsyzygyとcohomological shiftを取れば
 \(M=\operatorname{coker}\varphi\) としてよい。
 実際、元の加群との差に現れる \(C\)-perfect対象は制限後も \(A\)-perfectである。
 \(\psi\varphi=g\Id\) と \(g\) の正則性から \(\varphi\) は単射であり、
 \[
 0\longrightarrow Q_1\xrightarrow\varphi Q_0\longrightarrow M\longrightarrow0
 \]
 は \(S\)-自由分解である。
 これをhomological degree \(1,0\) に置き、
 \(d_Q=\varphi\)、\(s_Q:Q_0\to Q_1(w)\) を \(\psi\) とする。
 \(s_Q\) は \(Q_1\) 上では零とする。

 \(u_iM=0\) として \(M\) を \(R\)-加群に制限したときの有限自由分解は
 \[
 F_a=\bigoplus_{b+c=a}Q_b\otimes_\CC K_c,
 \qquad d=d_Q\otimes1+(-1)^b\otimes d_K
 \]
 である。\(b\in\{0,1\}\) は \(Q\) 側のhomological degreeである。
 \[
 s=s_Q\otimes1+(-1)^b\otimes h_K:F_a\longrightarrow F_{a+1}(w)
 \]
 と置くと、異なる因子の奇作用の反可換性から
 \[
 ds+sd=W\Id,\qquad s^2=0.                                  \tag{4f}\label{eq:central-shamash}
 \]
 が成り立つ。

 このデータの周期化が実際に \(A\)-自由分解を与えることも確認しておく。
 \(F_a=0\) for \(a<0\) とし、
 \[
 \mathcal F_n=
 \bigoplus_{j\ge0}(F_{n-2j}\otimes_R A)(-jw),\qquad n\ge0
 \]
 と置く。微分は同じ \(j\) 内の \(d\) と、\(j\ge1\) について
 \[
 s:(F_{n-2j}\otimes_R A)(-jw)
 \longrightarrow(F_{n-2j+1}\otimes_R A)(-(j-1)w)
 \]
 の和である。\eqref{eq:central-shamash}により平方は零になる。
 各 \(n\) で和は有限である。

 \(j\) による増大濾過を使うと、第一ページは
 \(H_*(F\otimes_R A)\) の内部・homological shiftの和となる。
 \(W\) が正則で \(WM=0\) なので、
 \[
 H_i(F\otimes_R A)=
 \begin{cases}
 M&i=0,\\
 M(-w)&i=1,\\
 0&i\ne0,1.
 \end{cases}
 \]
 次の微分は、\(j\ge1\) の \(H_0\) を \(j-1\) の \(H_1\) へ送る同型である。
 これを \(ds+sd=W\Id\) から直接見られる。
 \(a\in F_0\) が \(M\) の元を表すなら、\(s(a)\) の \(A\) 上の像はcycleであり、
 \(d s(a)=Wa\) である。
 短完全列
 \[
 0\longrightarrow F(-w)\xrightarrow{W}F
      \longrightarrow F\otimes_R A\longrightarrow0
 \]
 のconnecting mapはそのcycleを元の \(a\) に送る。
 従ってこの微分は、上の \(H_1\simeq M(-w)\) の標準同定の下で恒等である。
 隣接する二項はすべて相殺し、\(j=0\) の \(H_0=M\) だけが残る。
 濾過は各全homological degreeで有限なので収束し、
 \(\mathcal F\) は \(M\) の \(A\)-自由分解となる。

 \(F\) は有限複体なので、この分解の十分先は二周期的である。
 そのfactorizationは
 \[
 \widehat F_0=\bigoplus_{j\ge0}F_{2j}(jw),\qquad
 \widehat F_1=\bigoplus_{j\ge0}F_{2j+1}(jw),
 \qquad d+s
 \]
 による折り畳みであり、Koszul符号と内部シフトを含めて
 \(Q\boxtimes K_P\) に一致する。
 実際、\(\widehat F_0\) は
 \(Q_0\otimes(K_P)_0\oplus Q_1\otimes(K_P)_1(w)\)、
 \(\widehat F_1\) は
 \(Q_1\otimes(K_P)_0\oplus Q_0\otimes(K_P)_1\) である。
 従って制限加群の安定化は
 \[
 \operatorname{stab}_A(M)\simeq Q\boxtimes K_P.               \tag{4g}\label{eq:restriction-mf}
 \]
 ここで全体のshiftにも曖昧さはない。
 十分大きい \(\ell\) について
 \(\mathcal F_{2\ell}=(\widehat F_0\otimes_R A)(-\ell w)\)、
 \(\mathcal F_{2\ell+1}=(\widehat F_1\otimes_R A)(-\ell w)\) であり、
 そのcokernelは \(\Omega_A^{2\ell}M\) である。
 \(\Omega_A^{2\ell}M\simeq M[-2\ell]\) と \((w)\simeq[2]\) によって
 この内部シフトはちょうど打ち消される。

 最後に、\(\Gamma/\ZZ e_0\simeq\ZZ/3\) であるため、
 任意の \(\Gamma\)-次数付き \(C\)-加群は、
 三つの次数類 \(0,e_1,2e_1\) に属する部分に分解する。
 各部分は通常の \(\ZZ\)-次数付き \(C\)-加群の内部シフトであり、
 \eqref{eq:restriction-mf}の右辺の \(T\) 側を \((j)\)、\(j=0,1,2\)、
 だけシフトすればすべて得られる。
 前段のKoszul filtrationと合わせて、外部積の像は右辺全体を厚く生成する。
 これで本質的全射性が従う。
 周期性および二つの次数シフトとの両立は、ここで用いた外部和の構成そのものから従う。
\end{proof}

\begin{remark}[一変数の場合の明示行列]\label{rem:central-one-variable}
 \(m=1\)、\(P=u^3\) では \(K_P=(u,u^2)\) である。
 上の分解とホモトピーは
 \[
 Q_1(-e_1)\xrightarrow{(-u,\varphi)^t}
 Q_1\oplus Q_0(-e_1)\xrightarrow{(\varphi,u)}Q_0,
 \]
 \[
 s_0=(\psi,u^2)^t,\qquad s_1=(-u^2,\psi)
 \]
 となる。折り畳んだfactorizationは
 \[
 \alpha=
 \begin{pmatrix}\varphi&u\\-u^2&\psi\end{pmatrix}:
 Q_1\oplus Q_0(-e_1)\longrightarrow Q_0\oplus Q_1(w-e_1),
 \]
 \[
 \beta=
 \begin{pmatrix}\psi&-u\\u^2&\varphi\end{pmatrix}:
 Q_0\oplus Q_1(w-e_1)\longrightarrow
 Q_1(w)\oplus Q_0(w-e_1).
 \]
 中心性により \(\alpha\beta=\beta\alpha=(g+u^3)\Id\) である。
 これらの式は、一般の \(K_P\) による安定化の符号と次数シフトの具体的な検算になる。
\end{remark}

\subsection{有限群による整数次数の回復}

\(\Gamma\to\ZZ\)、\(e_0,e_1\mapsto1\)、の核は
\(T_\Gamma\simeq\ZZ/3\) であり、\(e_1-e_0\) で生成される。
分裂 \(1\mapsto e_0\) の下で、その双対群
\(H_\Gamma=\Hom(T_\Gamma,\CC^\times)\simeq\mu_3\) は
\(S\) 上恒等、すべての \(u_i\) を同時に三乗根倍する作用である。
\(P\) は斉次三次式なのでこの作用は \(W=g+P\) を保ち、
\[
 \MF^\Gamma(R,W)
 \simeq\bigl(\MF^\ZZ(R,W)\bigr)^{H_\Gamma}.                \tag{5}\label{eq:eq}
\]
有限群の位数は可逆である。
同変な有限次数付き自由加群は、連結な整数次数と有限群表現の半単純性により
同変な斉次基底を持つので、この比較は有限自由factorizationのdg模型で実現する。
さらに双対群 \(T_\Gamma\) で同変化すると、元の通常次数の圏が回復される
\cite[Theorem 4.2, Section 8, Lemma 8.6, Theorem 8.9]{Elagin}。
冪等完備化はこの比較の一部として取る。

外部積 \eqref{eq:central-TS} の左辺で双対作用の生成元は
\[
 \rho_T\rho_S^{-1}
\]
である。ここで \(\rho_S,\rho_T\) はそれぞれの因子の内部次数一のシフトである。
その三乗を恒等と同一視する自然同型は、両因子の
\((3)\simeq[2]\) から得られる。
二つのperiodicityを含めて \(\rho_S\) を保つdg同値はこの作用と両立し、
\eqref{eq:eq}と双対同変化により通常次数へ降りる。
この際 \(e_0\) のシフトも保存されるので、通常次数のdegree shift \((1)\) が保存される。
有限群作用のcoherenceはfactorizationの内部シフトのcoherenceから引き継ぎ、
単に自己同値の三乗が同型になるという存在命題だけから群作用を作るのではない。

\subsection{Sklyaninブロックへの適用}

\begin{proof}[定理\ref{thm:global}の証明]
 補題\ref{lem:TS}を \((S_h,G_h)\) と
 \((\CC[x,y,z],x^3+y^3+z^3+h^2xyz)\) のそれぞれに適用する。
 系\ref{cor:block}は第一因子の次数シフトと三次周期性を保ち、
 \(\MF^\ZZ(T,P)\) 側では恒等を用いる。
 前節の降下によって通常次数のdegree shiftを保つdg Morita同値
 \[
 \MF^\ZZ(S_h[u],G_h+P)
 \simeq
 \MF^\ZZ\!\left(\CC[x,y,z,u],x^3+y^3+z^3+h^2xyz+P\right)
 \]
 を得る。体上のOrlov同値と \cite{CCKM} により、これは残余圏の比較である。

 可換方程式の \(u_i\)-偏微分がすべて零なら、
 \eqref{eq:central-isolated}によって全 \(u_i=0\) である。
 \(h^6\ne-27\) の下では残る平面三次式のアフィン臨界点も原点だけである。
 よって射影超曲面 \(Y_{h,P}\) は滑らかである。
 次元は \(m+1\ge3\) なので、一意性は \cite[Theorem 1.7]{LZ} による。
\end{proof}



\begin{corollary}[平行移動を伴う一般の楕円ブロック]\label{cor:elliptic-central-cubic}
 \(S\) をHilbert series \((1-t)^{-3}\) を持つnoetherianな
 二次三次元AS-regular \(\CC\)-代数、\(g\in S_3\) を中心正則元とする。
 滑らかな平面三次曲線 \(E\)、\(L=\OO_E(1)\)、および
 ある原点に関する平行移動 \(\sigma=t_p\) に対して、次数付き同型
 \[
 S/(g)\simeq B(E,L,\sigma)
 \]
 が存在すると仮定する。
 \(f_E\in\CC[x,y,z]_3\) を \(E\subset\PP^2\) の方程式とし、
 \(P\in\CC[u_1,\ldots,u_m]_3\)、\(m\ge2\)、は
 \eqref{eq:central-isolated}を満たすとする。
 このとき \(S[u]/(g+P)\) のrotation付き残余圏の滑らかな可換モデルは
 \[
 V(f_E(x,y,z)+P(u))\subset\PP^{m+2}
 \]
 であり、射影同型を除いて一意である。
\end{corollary}
\begin{proof}
 Gorenstein parameter零のOrlov同値により
 \(\MF^\ZZ(S,g)\) は \(D^b(E)\) と同値であり、
 次数シフトは補題\ref{lem:elliptic-conjugacy}の
 \(T_{\OO_E}\sigma_*(-\otimes L)\) に対応する。
 同補題の直線束による共役とperiodicityのscalar調整により、
 \(\MF^\ZZ(\CC[x,y,z],f_E)\) との、次数シフトと三次周期性を保つdg同値を得る。
 あとは補題\ref{lem:TS}と有限次数差による降下を適用すればよい。
 滑らかさと一意性は定理\ref{thm:global}と同じ証明による。
 方程式 \(f_E\) の非零scalar倍による違いは、
 分離された \(x,y,z\) の同時三乗根倍によって射影同型に吸収される。
\end{proof}

\begin{remark}[この拡張の範囲]
 定理\ref{thm:global}では中心変数間の任意の三次混合項を許すが、
 非可換ブロックの生成元と中心変数の混合項は含めない。
 別の非可換ブロックについては、その生成元による通常の偏微分が代数の導分とは限らず、
 本節の \(J(P)\)-primaryな支持への帰着は適用できない。
 この点は独立した生成性の問題として残る。
 本節の拡張の先行研究上の新規性は未確認である。
\end{remark}

\section{複数ブロックの比較と生成性の帰着}\label{sec:multiblock}

\subsection{生成性を仮定しない外部積}

\(A_1,\ldots,A_s\) を、\((A_i)_0=\CC\) を満たす局所有限な
非負次数付き代数、\(f_i\in Z(A_i)_3\) を中心的な三次斉次元とする。
この小節の \(\MF\) は、有限自由次数付きmatrix factorizationのdg圏の
pretriangulated hullを取り、さらに冪等完備化したものを表す。
この定義にはtensor product algebraのnoetherian性を仮定しない。
\[
 \Gamma_s=\ZZ^s/\langle3e_i-3e_1:2\le i\le s\rangle,
 \qquad w=3e_1,
 \qquad A=\bigotimes_{i=1}^s A_i,
 \qquad W=\sum_{i=1}^s f_i
\]
とする。\(A_i\) の次数 \(d\) の斉次元には \(de_i\) を割り当てる。

\begin{lemma}[外部積の完全忠実性]\label{lem:multiblock-ff}
 外部積はdg完全忠実関手
 \[
 \Phi:\Perf\!\left(\bigotimes_{i=1}^s\MF^{\ZZ}(A_i,f_i)\right)
 \longrightarrow\MF^{\Gamma_s}(A,W)
 \]
 を定める。\(i\) 番目の因子のdegree shiftは右辺の \((e_i)\) に対応し、
 \((3e_i)=[2]\) の標準同型とも両立する。
\end{lemma}
\begin{proof}
 まず有限自由factorization \(P_i,Q_i\) を取る。
 tensor productの奇微分はKoszul符号付きの外部和であり、
 その二乗は \(W\) となる。
 各因子のHom複体を \(C_i^\bullet\) と書くと、
 自由加群のHomを行列係数ごとに分解することにより
 \[
 \Hom_{\MF^{\Gamma_s}(A,W)}^\bullet
       (\boxtimes_iP_i,\boxtimes_iQ_i)
 \cong
 \operatorname{Tot}\!\left(\bigotimes_i C_i^\bullet\right)
 \tag{\(\dagger\)}\label{eq:multiblock-hom}
 \]
 を得る。ここでは次数群の関係 \(3e_i=3e_1\) が、
 各因子の二つのcohomological shiftを他の因子へ移す操作に正確に対応する。
 従って右辺の次数 \(q\) の和は
 \(q_1+\cdots+q_s=q\) にわたる通常のtensor product複体の和であり、
 追加の完備化や直積を取らない。

 この和の有限性を確認する。各 \(P_i,Q_i\) は有限個の内部次数shiftからなり、
 \(A_i\) は非負次数なので、\(C_i^q=0\) は十分小さい \(q\) で成り立つ。
 したがって固定した全次数 \(q\) に寄与する \((q_1,\ldots,q_s)\) は有限個である。
 同様に、行列係数を元の \(\ZZ^s\)-次数へ分解したときも、
 固定した \(\Gamma_s\)-次数と非負次数の条件の下では有限個の成分しか現れない。
 これにより \eqref{eq:multiblock-hom} は実際のHom複体の同型である。
 differentialとcompositionの一致はtensor productのKoszul符号から従う。
 完全忠実性は有限個のcone、shift、直和因子を取っても保たれるので、主張を得る。
 内部次数shiftとperiodicityとの両立はこの構成自体に含まれる。
\end{proof}

\subsection{任意個の楕円THCRの積と周囲の正則性}

Chanの多重Serre同値 \cite[Theorem 3.4]{ChanMulti} は、代数のnoetherian性を
仮定せず、right ampleな可逆bimoduleの組に適用できる。
同論文のCorollaries 5.8--5.9は二因子のTHCRと三次元AS-regular代数の積を扱う。
以下では、その証明で用いる面と帯の議論をfactor数について帰納し、任意個の楕円因子を扱う。
Keeler \cite[Theorem 2.7, Corollary 2.8]{KeelerMulti} によっても、
この積のbimoduleの組の左右NC-amplenessを確認できる。
各translationはnumerical Picard groupに自明に作用し、全ての座標が正の外部積はampleだからである。
ただし多重amplenessだけではnoetherian性は自動的に従わず、以下で帯の有限性を別に扱う。

\begin{lemma}[楕円THCRの有限テンソル積]
\label{lem:elliptic-thcr-tensor-noeth}
 \(E_1,\ldots,E_r\) を滑らかな複素射影楕円曲線、\(L_i\) を正次数の可逆層、
 \(\sigma_i\) をtranslationとする。このとき
 \[
 C=\bigotimes_{i=1}^r C_i,\qquad C_i=B(E_i,L_i,\sigma_i)
 \]
 は左右noetherianである。
\end{lemma}
\begin{proof}
 \(X=\prod_iE_i\) とし、各factorから引き戻した可逆bimoduleを
 \(\mathcal L_i\) と書く。異なるfactorの作用は可換であり、
 \(\mathcal L^{\mathbf d}=\mathcal L_1^{d_1}\cdots\mathcal L_r^{d_r}\) と置けば
 \[
 C_{\mathbf d}=H^0(X,\mathcal L^{\mathbf d})\qquad(\mathbf d\in\mathbb N^r)
 \]
 は \(C\) のテンソル積多重次数と乗法を与える。
 可逆bimoduleの逆を用いて \(\mathcal L^{\mathbf d}\) は
 \(\mathbf d\in\mathbb Z^r\) に対しても定義する。

 \emph{Ampleness.}
 外部積のvery ample可逆層 \(H=\boxtimes_iH_i\) を固定する。
 任意のcoherent層 \(\mathcal F\) を、有限個の \(H\) の負冪の直和からの全射で
 \(\dim X\) 段まで近似する。
 各 \(\mathcal L_i^{d_i}\) の台となる可逆層の次数は \(d_i\deg L_i\) なので、
 その各近似項を \(\mathcal L^{\mathbf d}\) でtwistした高次cohomologyは、
 \(d_i\) が全て十分大きいと消える。
 ここでは楕円曲線上の正次数の可逆層の \(H^1\) の消滅とK\"unneth公式を用いる。
 最後の核にはGrothendieck消滅を適用し、dimension shiftingによって
 \[
 H^q(X,\mathcal F\mathcal L^{\mathbf d})=0
 \quad(q>0,\ \mathbf d\gg0)
 \]
 を得る。従ってこの組はright ampleである。
 \(\mathcal F\otimes H^{-q}\) にもこの消滅を適用すると、
 Castelnuovo--Mumford regularityにより
 \(\mathcal F\mathcal L^{\mathbf d}\) の大域生成も \(\mathbf d\gg0\) で従う。
 自動同型によるpushforwardはcohomologyと大域生成を変えないので、
 この議論は可逆bimoduleの左右の規約に依存しない。

 \emph{面と有限幅の帯.}
 代数 \(C\) の右noetherian性と、全てのcoherent \(\mathcal F\)、\(q\ge0\) に対する
 \[
 Q^q(\mathcal F)=\bigoplus_{\mathbf d\in\mathbb N^r}
 H^q(X,\mathcal F\mathcal L^{\mathbf d})
 \]
 の有限生成性を、\(r\) について同時に示す。
 \(r=0\) の場合は有限次元ベクトル空間の主張である。
 \(r-1\) 個の因子まで成立すると仮定し、
 \(C_{\widehat i}=\bigotimes_{j\ne i}C_j\)、\(X_{\widehat i}=\prod_{j\ne i}E_j\) と置く。
 座標 \(d_i=a\) を固定した \(Q^q(\mathcal F)\) の面は、\(C_{\widehat i}\)-加群である。
 射影 \(\pi_i:X\to X_{\widehat i}\) と
 \(\mathcal F_i=\mathcal F\mathcal L_i^a\) に対するLerayスペクトル系列の
 全多重次数の和は
 \[
 E_2^{p,b}=Q^p_{X_{\widehat i}}
       (R^b\pi_{i*}\mathcal F_i)
 \ \Longrightarrow\
 \bigoplus_{\mathbf d_{\widehat i}\ge0}
 H^{p+b}(X,\mathcal F_i\mathcal L_{\widehat i}^{\mathbf d_{\widehat i}})
 \]
 となる。射影公式と残るtranslationに関する自然性により、これは
 \(C_{\widehat i}\)-加群のスペクトル系列である。
 帰納法で各 \(E_2\) 項はnoetherianな \(C_{\widehat i}\) 上有限生成である。
 従って、有限個の核・商・拡張で得られる各面もnoetherianである。

 次の帯に関する観察を使う。
 \(\mathbb N^r\)-次数付き \(C\)-加群 \(M\) の各面 \(M_{d_i=a}\) が
 \(C_{\widehat i}\) 上noetherianであり、\(M_{\ge\mathbf b}=0\) とする。
 \[
 M^{(i)}=\bigoplus_{d_1\ge b_1,\ldots,d_i\ge b_i}M_{\mathbf d}
 \quad(0\le i\le r)
 \]
 は \(M=M^{(0)}\supset\cdots\supset M^{(r)}=0\) という \(C\)-部分加群のfiltrationである。
 各 \(M^{(i-1)}/M^{(i)}\) をさらに \(d_i\ge a\) でfiltrationすると、
 有限個の商は固定面の \(C_{\widehat i}\)-部分加群になる。
 その商では \((C_i)_{>0}\) の作用は零なので、これらは \(C\)-加群としてもnoetherianである。
 よって \(M\) はnoetherianである。この観察は \(C\) 自身のnoetherian性を使わない。
 各面がnoetherianな加群の部分商についても、同じ観察を適用できる。

 \emph{切断加群.}
 \(q>0\) のとき、amplenessにより \(Q^q(\mathcal F)_{\ge\mathbf b}=0\) となる
 \(\mathbf b\) がある。従って前段階から
 \(Q^q(\mathcal F)\) は\emph{noetherian}である。
 特に、後で用いる \(Q^1\) の任意の部分加群も有限生成である。

 \(q=0\) について、\(\mathcal F\mathcal L^{\mathbf a}\) を大域生成する
 \(\mathbf a\ge0\) を選び、逆bimoduleをtensorして
 \[
 0\longrightarrow\mathcal K\longrightarrow
       (\mathcal L^{-\mathbf a})^{\oplus t}
       \longrightarrow\mathcal F\longrightarrow0
 \]
 を得る。\(\mathcal K\) はcoherentである。
 この積の場合、\(\mathbf v\) の一つでも成分が負なら
 \(H^0(X,\mathcal L^{\mathbf v})=0\) である。
 実際、外部積のそのfactorの次数が負であり、K\"unneth公式を適用できる。
 従って切断を全非負多重次数で取った列は
 \[
 C(-\mathbf a)^{\oplus t}\longrightarrow Q^0(\mathcal F)
       \longrightarrow Q^1(\mathcal K)
 \]
 と正確である。最初の写像の像は有限生成であり、余核は既にnoetherianと示した
 \(Q^1(\mathcal K)\) の部分加群なので、\(Q^0(\mathcal F)\) は有限生成である。
 また任意の \(\mathbf a\ge0\) について
 \[
 Q^0(\mathcal F)_{\ge\mathbf a}(\mathbf a)
 \cong Q^0(\mathcal F\mathcal L^{\mathbf a})
 \]
 であるから、第一象限によるtruncationも有限生成である。

 \emph{イデアルの有限生成.}
 多重斉次右イデアル \(I\subset C\) を取る。
 \cite[Theorem 3.4]{ChanMulti} により、sheafificationはexactであり、
 \(\widetilde I\subset\widetilde C=\OO_X\) となる。
 \(X\) はnoetherianなので、このquasi-coherent部分層はcoherentである。
 随伴のunit \(\eta_I:I\to Q^0(\widetilde I)\) の自然性と
 \(\eta_C:C\xrightarrow{\sim}Q^0(\OO_X)=C\) により、\(\eta_I\) は単射である。
 その余核は多重tailsに関するtorsionであり、\(Q^0(\widetilde I)\) の商だから有限生成である。
 有限生成torsion加群は一つの第一象限上で零となるため、ある \(\mathbf a\ge0\) に対して
 \[
 I_{\ge\mathbf a}=Q^0(\widetilde I)_{\ge\mathbf a}
 \]
 となる。右辺の有限生成性は前段階で証明した。
 一方、\(I/I_{\ge\mathbf a}\) は \(C/C_{\ge\mathbf a}\) の部分加群である。
 \(C\) の固定面は有限次元空間 \((C_i)_a\) と \(C_{\widehat i}\) のtensor productなので
 \(C_{\widehat i}\)-noetherianである。
 従って帯の観察により \(C/C_{\ge\mathbf a}\) はnoetherianであり、
 \(I/I_{\ge\mathbf a}\) も有限生成となる。これで \(I\) の有限生成性を得た。

 加法と両立する \(\mathbb N^r\) の整列順序、例えば全次数付き辞書式順序で
 任意の右イデアルの先頭多重斉次項を取る。
 その斉次イデアルの有限生成元を元のイデアルへ持ち上げ、
 先頭次数を下げる帰納法を行えば、任意の右イデアルも有限生成である。
 これで通常の右noetherian性と帰納法が完了する。
 opposite algebraには逆translationを伴う正次数の可逆bimoduleが対応するので、
 同じ証明により左noetherian性も得る。
\end{proof}

\begin{lemma}[正次数中心正則元からのnoetherian性の持ち上げ]
 \label{lem:central-noetherian-lift}
 \(D\) を連結局所有限な非負次数付き代数、
 \(g\in Z(D)\) を正次数の斉次正則元とする。
 \(D/(g)\) がnoetherianなら \(D\) もnoetherianである。
\end{lemma}
\begin{proof}
 \(g\)-adic associated gradedは \((D/(g))[t]\) でありnoetherianである。
 任意の斉次右イデアル \(J\) に誘導された \(g\)-adic filtrationの
 associated gradedはこの環の右イデアルなので有限生成である。
 その二重斉次生成元を \(J\) に持ち上げる。
 内部次数 \(d\) の元のinitial termをこれらの持ち上げで消すと、
 \(g\)-adic orderが真に増える。
 しかし \(g^nD\) の内部次数は少なくとも \(n\deg(g)\) なので、
 固定した \(d\) ではこの操作は有限回で終わる。
 よって全ての斉次右イデアルが有限生成である。
 任意の右イデアルについては通常の内部次数の最高次項を取り、
 生成元を持ち上げて次数に関する帰納法を行えば有限生成性を得る。
 左側も同じである。
\end{proof}

\begin{proposition}[周囲のテンソル積]
\label{prop:tensor-sklyanin-ambient}
 \(h_i^6\ne-27\)、\(S_i=S_{h_i}\)、\(g_i=G_{h_i}\) とし、
 \[
 A=(\bigotimes_{i=1}^rS_i)[u_1,\ldots,u_m],\qquad
 N=3r+m,\qquad W=\sum_i g_i+P(u),\qquad R=A/(W)
 \]
 と置く。ここで \(P\) は任意の中心斉次三次式であり、\(m=0\) なら \(P=0\) とする。
 \(A\) は左右noetherianなKoszul AS-regular代数で、
 \[
 \operatorname{gldim}A=N,\qquad H_A(t)=(1-t)^{-N},\qquad
 \operatorname{Ext}^j_A(\CC,A)=
 \begin{cases}\CC(N)&j=N,\\0&j\ne N\end{cases}
 \]
 を満たす。\(W\) は中心正則元であり、\(R\) はinjective dimension \(N-1\)、
 Gorenstein parameter \(N-3\) のAS-Gorenstein代数である。
 この主張に \(\sigma_i\) の有限位数や \(P\) の孤立性は必要ない。
\end{proposition}
\begin{proof}
 \(S_i/(g_i)\) の斉次基底の持ち上げ \(V_i\subset S_i\) を選ぶと、
 \(g_i\) の正則性と次数に関する帰納法により
 \(S_i\cong V_i\otimes\CC[g_i]\) は中心 \(\CC[g_i]\)-加群として自由である。
 \(g_1,\ldots,g_r\) は \(\bigotimes_iS_i\) 上の中心正則列をなし、その商は
 補題\ref{lem:elliptic-thcr-tensor-noeth}の \(C\) である。
 補題\ref{lem:central-noetherian-lift}により、正次数中心正則元について
 商のnoetherian性を元の連結次数付き代数へ持ち上げられる。
 これを反復し、中心変数についてのHilbert基底定理を用いると \(A\) は左右noetherianとなる。

 各 \(S_i\) の長さ三のKoszul双加群分解と、多項式環のKoszul双加群分解をtensorすると、
 \(A\) の対角双加群の長さ \(N\) の有限射影分解を得る。
 よって \(\operatorname{gldim}A\le N\) である。
 各factorの \(\CC\) のKoszul分解とそのdualのK\"unneth公式は、
 表示したExt公式と左側の同じ式を与える。
 従ってglobal dimensionはちょうど \(N\) であり、AS-Gorenstein条件も成立する。
 Koszul性とHilbert seriesはtensor分解から従い、GK dimensionは \(N\) である。

 \(A\) は \(\CC[g_1,\ldots,g_r,u_1,\ldots,u_m]\) 上自由である。
 この中心多項式環で \(W\) は \(g_1\) に関してmonic一次式なので、\(A\) 上でも左右正則である。
 中心正則元によるchange of ringsにより、商のinjective dimensionは一つ減り、
 Gorenstein parameterは三減る。
\end{proof}

\subsection{多ブロックHesse模型からの埋込み}

\(r\ge1\)、\(m\ge0\) とし、独立した変数 \(x_j,y_j,z_j\) に対して
\[
 S_{h_j}=\CC\langle x_j,y_j,z_j\rangle/
 ([x_j,y_j]-h_jz_j^2,[y_j,z_j]-h_jx_j^2,[z_j,x_j]-h_jy_j^2),
\]
\[
 G_j=x_j^3+y_j^3+z_j^3-2h_j^2\sym(x_jy_jz_j),\qquad
 F_j=x_j^3+y_j^3+z_j^3+h_j^2x_jy_jz_j
\]
と置く。異なるブロックの元は可換とし、
\[
 A_{\mathbf h,m}=(S_{h_1}\otimes\cdots\otimes S_{h_r})[u_1,\ldots,u_m],
 \qquad W_{\mathbf h,m}=\sum_{j=1}^rG_j+\sum_{a=1}^m u_a^3,
\]
\[
 F_{\mathbf h,m}=\sum_{j=1}^rF_j+\sum_{a=1}^m u_a^3,
 \qquad Y_{\mathbf h,m}=V(F_{\mathbf h,m})\subset\PP^{3r+m-1}
\]
とする。

\begin{proposition}[多ブロックの完全忠実比較]\label{prop:multiblock-embedding}
 全ての \(j\) について \(h_j^6\ne-27\) と仮定する。
 degree shiftと三次periodicityを保つdg完全忠実関手
 \[
 \MF^{\ZZ}(\CC[x_j,y_j,z_j,u_a],F_{\mathbf h,m})
 \longrightarrow
 \MF^{\ZZ}(A_{\mathbf h,m},W_{\mathbf h,m})
 \]
 が存在する。\(n=3r+m-2\ge3\) なら、左辺は
 \((\Ku(Y_{\mathbf h,m}),\rho_{Y_{\mathbf h,m}})\) である。
 この命題は、右辺の全対象が像に属するとは主張しない。
\end{proposition}
\begin{proof}
 \(s=r+m\) とし、各非可換ブロックと各一変数ブロックに別の \(e_i\) を割り当てる。
 系\ref{cor:block}は各 \(j\) について
 \[
 \MF^{\ZZ}(\CC[x_j,y_j,z_j],F_j)
 \simeq\MF^{\ZZ}(S_{h_j},G_j)
 \]
 をdegree shiftとperiodicityを保って与える。
 これらをtensorし、\(\MF^{\ZZ}(\CC[u_a],u_a^3)\) には恒等を用いる。
 非可換側へは補題\ref{lem:multiblock-ff}の完全忠実関手を用いる。

 可換側の外部積は本質的全射でもある。
 これは可換なsmooth affine ambientに対するequivariant Thom--Sebastiani定理
 \cite[Theorem 5.15, Corollary 5.18]{BFK}を適用したものである。
 二因子の場合は \(G=H=\mathbb G_m\)、potentialのcharacterを双方 \(t\mapsto t^3\)
 と取る。異なる二因子上でのcharacterの差 \((-3,3)\) は非torsionであり、
 \(G\times_{\mathbb G_m}H\) のcharacter groupは
 \(\ZZ^2/\langle3e_2-3e_1\rangle\) となる。
 各potentialは斉次なので、Euler公式により和のcritical locusは
 個々の零点の積に含まれる。従って同定理の支持条件を除くことができる。
 反復して \(\Gamma_s\) を得る。
 ここで引用している定理のambientは全て可換であり、非可換ambientへ拡張して引用してはいない。

 これで \(\Gamma_s\)-次数付きの完全忠実関手を得た。
 \(\Gamma_s\to\ZZ\)、\(e_i\mapsto1\) の核は
 \(T_s\cong(\ZZ/3)^{s-1}\) である。
 細分次数の圏から通常次数の圏へ戻す双対同変化では、
 \(T_s\) の生成元は因子のdegree shiftの差
 \(\rho_i\rho_1^{-1}\) として作用する。
 全ての比較はこれらと三次periodicityを保つので、
 有限群関係とそのcoherenceを保つ同変関手となる。
 同変対象間のHomは元のHomの有限群不変部分であり、標数零では不変部分を取る操作は完全である。
 従って完全忠実性は同変化で保たれる。
 \cite[Theorem 4.2 and Section 8]{Elagin}の双対同変化を用いて通常次数へ降ろせば、
 求める関手を得る。\((e_1)\) との両立は降下後の \((1)\) との両立になる。

 最後に \(Y_{\mathbf h,m}\) のsmoothnessを確認する。
 特異点では全ての \(u_a=0\) であり、各三変数ブロックの偏微分が同時に零となる。
 \(F_j\) の射影零点がsmoothなので、そのaffine critical locusは原点だけである。
 従って全座標が零となり、射影特異点は存在しない。
 Orlov同値により最後の主張を得る。
\end{proof}

\begin{remark}
 命題\ref{prop:tensor-sklyanin-ambient}により、この独立ブロック族の周囲は
 平行移動の位数によらずnoetherian AS-regularである。
 中心正則元の商についてのMF--特異圏比較とOrlov windowにより、
 右辺は \(A_{\mathbf h,m}/(W_{\mathbf h,m})\) の残余成分と同一視できる。
 この周囲の正則性は、残余圏そのもののpropernessや比較の本質的全射性とは別である。
\end{remark}

\subsection{Euler導分と中心支持への帰着}

\begin{lemma}[ブロックpotentialの零ホモトピー]\label{lem:block-euler}
 任意の有限自由factorization \(P\) of \(W_{\mathbf h,m}\) 上で、
 各 \(G_j:P\to P(w)\) と各 \(u_a^2:P\to P(2e_{r+a})\) は零ホモトピックである。
 この主張は \(\Gamma_s\)-次数付きでも通常の整数次数付きでも成り立つ。
\end{lemma}
\begin{proof}
 \(j\) 番目のブロックの次数を数え、他のブロックと \(u_a\) では零となるEuler導分を
 \(\delta_j\) と書く。Sklyanin関係式はそのブロックについて斉次なので、
 \(\delta_j\) は自由代数から商代数へ降りる。
 \(\delta_j(W_{\mathbf h,m})=3G_j\) である。
 factorizationの有限自由加群に斉次基底を選び、その微分行列 \(d\) の係数に
 \(\delta_j\) を作用させる。積のLeibniz則から
 \[
 d\,\delta_j(d)+\delta_j(d)\,d
 =\delta_j(d^2)
 =3G_j\,\Id_P.
 \]
 従って \(\frac13\delta_j(d)\) が所要のホモトピーである。
 \(\delta_j\) は代数の元の \(\Gamma_s\)-次数を変えない。
 \(3e_j=w\) という次数群の関係があっても、
 元のtensor factorsから定義した導分の存在は失われない。
 同じ計算を中心変数の導分 \(\partial_{u_a}\) に対して行えば
 \[
 d\,\partial_{u_a}(d)+\partial_{u_a}(d)\,d=3u_a^2\Id_P
 \]
 を得る。基底に依存するのはホモトピーの選択であり、零ホモトピックという結論ではない。
\end{proof}

\begin{proposition}[生成性の中心商への帰着]\label{prop:block-support}
 全ての \(h_j^6\ne-27\) とし、\(A=A_{\mathbf h,m}\) と置く。
 命題\ref{prop:tensor-sklyanin-ambient}により、この \(A\) はnoetherian AS-regularである。さらに
 \[
 R=A/(W_{\mathbf h,m}),\qquad
 C_j=S_{h_j}/(G_j),\qquad C=\bigotimes_{j=1}^r C_j
 \]
 と置く。\(C\) は \(R\) の商とみなす。
 次の等式が、冪等完備化した \(\Gamma_s\)-次数付き特異圏で成り立つ。
 \[
 \Dsg^{\Gamma_s}(R)
 =\thick\!\left(
 \operatorname{Res}\bigl(D^b(\operatorname{gr}^{\Gamma_s}C)\bigr)
 \right).
 \tag{\(\ddagger\)}\label{eq:block-support}
 \]
 また、\(C\) は \(R\)-perfectなので、restrictionは特異圏の間の関手にも降りる。
\end{proposition}
\begin{proof}
 各 \(G_j\) は中心正則元である。\(S_{h_j}/G_jS_{h_j}\) の斉次基底を持ち上げた
 ベクトル部分空間を \(V_j\subset S_{h_j}\) とすると、次数に関する帰納法と正則性から
 \[
 V_j\otimes_\CC\CC[t_j]\xrightarrow{\sim}S_{h_j},\qquad
 v\otimes t_j^b\longmapsto vG_j^b
 \]
 は中心 \(\CC[t_j]\)-加群の同型である。
 従って \(A\) は中心多項式環
 \(\CC[t_1,\ldots,t_r,u_1,\ldots,u_m]\) 上自由である。
 \(W\) はこの環の \(\sum_jt_j+\sum_au_a^3\) に対応するので正則であり、
 \(R\) は \(\CC[t_1,\ldots,t_{r-1},u_1,\ldots,u_m]\) 上自由となる。
 よって
 \[
 u_1,\ldots,u_m,G_1,\ldots,G_{r-1}
 \]
 は \(R\) の中心正則列で、その商は \(C\) である。
 Koszul分解は \(C\) の有限自由 \(R\)-分解を与える。
 これでrestrictionがperfect対象をperfect対象へ送ることも従う。

 まず冪等完備化前の \(\Dsg^{\Gamma_s}(R)\) の対象 \(N\) を取る。
 中心正則元によるAS-regular代数の商に対するMF比較
 \cite[Theorem 5.8]{CCKM}を使えば、\(N\) はfinite-free factorizationのcokernelのshiftで表せる。
 この比較は通常の整数次数で記述されているが、
 \(\Gamma_s\simeq\ZZ\oplus\Gamma_{s,\mathrm{tors}}\) と標数零を用い、
 各有限自由分解とホモトピーを斉次成分ごとに取る同じ構成で
 \(\Gamma_s\)-次数を保った比較を得る。
 補題\ref{lem:block-euler}より、\(N\) は
 \[
 K(G_1,\ldots,G_{r-1},u_1^2,\ldots,u_m^2;N)
 \]
 のshiftと内部twistの直和因子となる。
 ここで \(K\) は通常の中心Koszul複体である。
 この複体の各cohomologyは \(G_1,\ldots,G_{r-1},u_1^2,\ldots,u_m^2\) に殺される。
 さらに \((u_1,\ldots,u_m)^{m+1}\subset(u_1^2,\ldots,u_m^2)\) なので、
 各cohomologyの有限filtrationの商は全 \(u_a\) と全 \(G_j\) に殺される \(C\)-加群である。
 \(G_r\) の消滅は \(W=0\) から従う。
 noetherian性によりこれらは有限生成であり、有限個のtruncation triangleとfiltrationから
 未完備の特異圏の各対象が右辺に入る。
 最後に直和因子閉包を取れば、\eqref{eq:block-support} を得る。
\end{proof}

\subsection{必要な生成性の正確な内容}

命題\ref{prop:block-support}は、全ての \(C\)-加群のrestrictionが
\(C_1,\ldots,C_r\) の外部積から生成されることまでは述べていない。
十分条件の一つは
\[
 D^b(\operatorname{gr}^{\Gamma_s}C)
 =\thick\left\{
 \operatorname{coarsen}(M_1\boxtimes\cdots\boxtimes M_r)(\gamma)
 :M_j\in D^b(\operatorname{gr}^{\ZZ}C_j),\ \gamma\in\Gamma_s
 \right\}
 \tag{G}\label{eq:block-generation}
\]
である。\(\gamma\) の全shiftは、中心変数が零となった後にも残る次数類を含めるために必要である。
以下で用いるsyzygyへの置換を正当化するため、
\(A_{\ne i}=(\bigotimes_{j\ne i}S_{h_j})[u_1,\ldots,u_m]\) とし、
他因子の外部積に全ての \(u_a\) が零で作用する \(A_{\ne i}\)-加群を \(M_{\ne i}\) とする。
\(S_{h_i}\) の \(\CC[G_i]\) 上の自由性と \(W=0\) による \(G_i\) の消去から、
\(R\) は \(A_{\ne i}\) 上自由であり、
\[
 \operatorname{Res}(C_i\boxtimes M_{\ne i})
 \simeq M_{\ne i}\otL_{A_{\ne i}}R.
\]
各 \(S_{h_j}\) 上の有限自由分解と \(u_a\) のKoszul分解を外部tensorすると、
\(M_{\ne i}\) は \(A_{\ne i}\)-perfectであり、上のrestrictionは \(R\)-perfectである。
従って一因子が \(C_i\)-perfectである項はrestriction後に零となり、
各因子を有限syzygyへ置き換えてよい。有界複体の場合もtruncation triangleから従う。
各因子の有限syzygyを取りMFで表した後、外部tensor productの有限ambient分解を
\(W\) の零ホモトピーで周期化すると、そのrestrictionの安定化は外部積factorizationになる
（全体のcohomological shiftは生成性に影響しない）。
したがって (G) が成立すれば、命題\ref{prop:block-support}と
補題\ref{lem:multiblock-ff}により非可換側の外部積もMorita同値となり、
命題\ref{prop:multiblock-embedding}を圏同値へ強められる。
より弱く、(G)の等式をrestriction後の \(\Dsg^{\Gamma_s}(R)\) だけで示しても十分である。

しかし (G) を、multigraded tailsの可換な積の幾何だけから推論することはできない。
\(C\) は元来 \(\ZZ^r\)-次数を持つが、ここで必要なのは
\(\Gamma_s\)-次数付き加群である。\(\ZZ^r\to\Gamma_s\) の核は
\(3(e_i-e_1)\) からなる無限格子を含み、有限群降下ではない。
\(\Gamma_s\)-次数付き加群は一般には \(\ZZ^r\)-次数へ持ち上がらない。
最も単純には、\(\CC[a,b]\) に \(\deg a=e_1\)、\(\deg b=e_2\)、
\(3e_1=3e_2\) を入れると、\(\CC[a,b]/(a^3-b^3)\) は \(\Gamma_2\)-次数付きだが
\(\ZZ^2\)-次数へ持ち上がらない。そのannihilatorが \(\ZZ^2\)-斉次でないためである。
この例は (G) 自体を反証するものではなく、個々の対象を細分次数へ持ち上げる証明法を反証する。

さらにmultigraded tailsでは各座標軸上のtorsionが消える。
生成性を完成させるには、これらのtorsionと上記の非liftable加群を、
restriction後も含めて制御する必要がある。
第\ref{sec:finite-center}節では、高々一つを除く因子の有限中心性でこの両方を扱い、
一つの任意位数ブロックを許して本質的全射性を証明する。
二つ以上の無限位数因子の一般の場合は未解決として保持する。

\begin{remark}[別の十分条件]
 \(\mathcal A\to\cD\) をsmooth proper圏からのdg完全忠実関手とする。
 \(\cD\) がproperでSerre関手を持ち、埋込みがSerre関手と両立すれば、
 \(\mathcal A\) のsaturated性から左右随伴が存在し、その直交成分は左右一致する。
 したがって \(\cD\) に非自明な直交直和分解がなければ埋込みは同値である。
 これは上の生成性とは別の完成条件となる。
 ただし複数ブロックの全MFについて、必要なpropernessと直交分解不能性を
 ここで証明したわけではない。
\end{remark}

\section{一方有限中心の外部積生成性と多ブロックの可換モデル}\label{sec:finite-center}

この節では生成条件 \eqref{eq:block-generation} を、高々一つを除くブロックが
有限位数である場合に証明する。任意位数の一因子は係数代数として残し、
有限中心の因子を一つずつ追加する。
非可換係数付きBeilinson分解、中心のNoether正規化、分離的な稠密開集合と
その補集合の帰納法を用いる。
\(\ZZ^r\)-次数へ持ち上がらない対象と軸上のtorsionも初めから含む。
有限中心の場合の基礎的な局所化・生成法については \cite{ELS} と比較されたい。

\subsection{斉次中心局所化と稠密開集合}

\(\operatorname{gr}^{\Lambda}D\) は有限生成の \(\Lambda\)-次数付き右 \(D\)-加群の圏とする。
ここで \(\Lambda\) は有限生成アーベル群であり、torsionを持っていてよい。
全てのthick閉包は、有限直和、shift、cone、直和因子に関する閉包を表す。

\begin{lemma}[斉次中心局所化]\label{lem:graded-localization}
 \(D\) をnoetherianな \(\Lambda\)-次数付き代数、
 \(f\in Z(D)\) を斉次元とする。
 有界導来圏の局所化は
 \[
 D^b(\operatorname{gr}^{\Lambda}D)\big/
 D^b_{f\text{-tor}}(\operatorname{gr}^{\Lambda}D)
 \simeq D^b(\operatorname{gr}^{\Lambda}D_f)
 \]
 を与える。左辺の核は、各cohomologyが \(f\) の冪に殺される対象からなる。
 この核は \(D/(f)\) 上の有限生成 \(\Lambda\)-次数付き加群のrestrictionで生成される。
\end{lemma}
\begin{proof}
 局所化は中心元についての平坦な局所化であるからcohomologyと可換である。
 有限生成加群が局所化後に零なら、その有限個の生成元を同時に殺す \(f\) の冪がある。
 \(f\)-adic filtrationと有界truncation triangleにより核についての主張を得る。

 局所化後の有限生成次数付き加群は、有限個の斉次生成元の
 \(D\)-spanを取ることによって、局所化前の有限生成次数付き加群から得られる。
 射については有限生成加群の有限生成自由分解の必要な有限部分を使う。
 中心局所化と \(\operatorname{Ext}^q\) は可換であり、
 次数零の射は \(f\) の冪を分母としてclearできる。
 同じことは有界複体にも有限個のtruncationによって成立する。
 次数を保つ分母の射は \(f^N:M(-N\deg f)\to M\) と書く。
 そのconeは \(f\)-torsionである。
 従って上記Verdier商は局所化後の圏へ完全忠実に写り、
 加群とその有限個のconeを全て含むので本質的全射である。
 この議論で次数を細分する必要はない。
\end{proof}

\begin{lemma}[有限中心代数の良い稠密開集合]\label{lem:generic-separable}
 \(k\) を標数零の体とする。
 \(D\) を連結非負次数付き \(k\)-代数で、有限生成の可換次数付き中心部分代数
 \(Z\subset Z(D)\) 上有限なものとする。
 \(Z\) がreducedであると仮定する。
 ある斉次元 \(f\in Z\) と斉次両側イデアル \(I\subset D\) が存在し、次を満たす。
 \begin{enumerate}
 \item \(f\) は \(Z\) の全てのminimal primeを避ける。
       特に \(Z/(f)\ne0\) なら \(\dim Z/(f)<\dim Z\) である。
 \item \(Z_f\) は \(k\) 上smoothな有限型可換環である。
 \item \(I_f\) は冪零であり、\(Q_f=(D/I)_f\) は
       \(Z_f\) 上有限射影かつ分離的な代数である。
 \end{enumerate}
 \(\dim Z=0\) の場合には \(Z=k\)、\(f=1\) と取れる。
\end{lemma}
\begin{proof}
 \(K\) を \(Z\) の全商環とする。
 \(Z\) はreduced noetherianなので、\(K\) はminimal primeの分数体の有限積である。
 \(D_K\) は \(K\) 上有限次元であり、そのJacobson radical \(J\) は冪零である。
 \(I\) を \(D\to D_K\) による \(J\) の逆像とする。
 次数を与える \(\mathbb G_m\)-作用はminimal primeの集合と \(J\) を保つので、
 \(I\) は斉次である。
 \(I^N\) は十分大きい \(N\) で \(K\) 上零となり、
 有限生成性から、ある全商環の分母を反転すれば \(I^N=0\) となる。

 \(Q_K=D_K/J\) は半単純である。
 標数零なので、各simple factorの中心は対応する成分体の有限分離拡大であり、
 \(Q_K\) は \(K\) 上分離的である。
 generic freenessを用いて、ある稠密開集合上で \(Q\) を有限射影とする。
 分離性は、積の写像 \(Q\otimes_ZQ^{\mathrm{op}}\to Q\) の
 \(Q\)-双加群としてのsection、すなわち分離冪等元の存在で表せる。
 generic fiberにある分離冪等元と、それが満たす有限個の等式の分母を払えば、
 さらに小さい稠密開集合上でこの分離冪等元が存在する。
 また \(Z\) のsmooth locusは全てのgeneric pointを含む。

 以上の良い開集合は \(\mathbb G_m\)-作用で不変に取れる。
 実際、\(I^N\) の消滅、局所自由性、分離性、smoothnessは全て
 次数作用によって保たれる開条件である。
 補集合のイデアルは斉次で、どのminimal primeにも含まれない。
 \(Z_0=k\) と \(k\) の無限性により、そのイデアル内から
 全てのminimal primeを避ける一つの斉次元 \(f\) を選べる。
 必要なら候補を冪にして次数を揃え、一般の線型結合を取ればよい。
 \(\dim Z>0\) のとき零次元の既約成分はない
 （全ての既約成分は連結な錐の頂点を含む）ので、この操作に例外はない。
 \(Z=k\) の場合は有限次元代数のradicalと半単純商を用いればよい。
\end{proof}

\begin{lemma}[分離代数と射影分解]\label{lem:separable-perfect}
 \(Z\) をsmooth affineな有限型 \(k\)-代数、
 \(Q\) を \(Z\) 上有限射影かつ分離的な代数とする。
 このとき \(Q\) は有限global dimensionを持つ。
 特に \(Q\) に任意のアーベル群次数を与えたとき、
 \(D^b(\operatorname{gr}^{\Lambda}Q)\) は \(Q(\lambda)\)
 \((\lambda\in\Lambda)\) で生成される。
\end{lemma}
\begin{proof}
 任意の右 \(Q\)-加群 \(M\) は、分離冪等元により
 \(M\otimes_ZQ\to M\) の \(Q\)-線型sectionを持つ。
 \(Z\) は有限global dimensionを持ち、\(Q\) は \(Z\)-flatなので、
 \(M\) の \(Z\)-射影分解を \(-\otimes_ZQ\) で送ると
 \(M\otimes_ZQ\) の有限 \(Q\)-射影分解を得る。
 従ってその直和因子 \(M\) も有限射影次元を持つ。

 次数付き有限生成加群に有限生成斉次自由分解を作ると、十分高いsyzygyは
 ungradedには射影的である。
 自由加群からこのsyzygyへの次数零の全射のungraded sectionを取り、
 sectionの次数零成分を抜き出すと次数付きsectionになる。
 よってsyzygyは次数付き有限自由加群の直和因子であり、
 元の加群は有限個の \(Q(\lambda)\) から生成される。
\end{proof}

\subsection{非可換係数の多項式拡大と一方有限中心の生成定理}

ここでは係数代数の有限中心性を仮定しない。
これが一つの無限位数ブロックを許すための要点である。
以下、\(k\) は標数零の体、\(\Lambda\) は有限生成アーベル群とする。

\begin{lemma}[非可換係数を持つ多項式拡大]\label{lem:polynomial-coefficient}
 \(A\) を左右noetherianな \(\Lambda\)-次数付き \(k\)-代数とする。
 \(P=k[t_1,\ldots,t_d]\) の各変数に任意の \(\Lambda\)-次数を与え、
 \(A\otimes_kP\) には和の次数を与える。このとき
 \[
 D^b(\operatorname{gr}^{\Lambda}(A\otimes_kP))
 =\thick\{M\otimes_kP:\ M\in\operatorname{gr}^{\Lambda}A\}.
 \]
 \(P\) の斉次元 \(f\) を反転した \(P_f\) についても同じ結論が成立する。
 \(A\) が有限global dimensionを持つという仮定はない。
\end{lemma}
\begin{proof}
 中心多項式拡大のHilbert基底定理により \(A\otimes P\) は左右noetherianである。
 \(d=0\) は自明なので \(d>0\) とする。
 有限生成 \(\Lambda\)-次数付き \(A\otimes P\)-加群 \(T\) の有限斉次表示を取る。
 射影用の独立な整数次数を導入し、\(t_0,t_1,\ldots,t_d\) の射影次数を全て1、
 \(A\) の射影次数を0とする。\(\Lambda\)-次数は
 \(\deg_\Lambda(t_0)=0\)、\(\deg_\Lambda(t_i)=\deg_\Lambda(t_i\in P)\) とする。
 表示の各関係を \(t_0\) で斉次化すると、\(\PP^d_k\) 上の
 \(\mathcal A=A\otimes_k\OO_{\PP^d}\)-加群 \(\cF\) が得られる。
 これは \(\mathcal A\) 上有限表示で、\(t_0\ne0\) への制限が \(T\) である。
 各関係の斉次化では \(\Lambda\)-次数を変えていない。
 以後の \(\Lambda\)-次数は対角化可能群 \(D(\Lambda)\) による線型化としても解釈できる。

 必要なcohomologyの有限性を確認する。
 \(\mathcal A\) 上のcoherent加群には有限個の
 \(\mathcal A(-a)(\lambda)\) からの全射があり、核もcoherentである。
 これは中心変数による有限斉次表示とnoetherian性から従う。
 標準affine被覆の\v Cech複体は長さ \(d\) なので、cohomologyは次数 \(0,\ldots,d\)
 に限られる。一方、\(H^q(\PP^d,\mathcal A(a))
 =A\otimes_kH^q(\PP^d,\OO(a))\) は有限生成 \(A\)-加群である。
 上記全射とその核の長完全列を用い、\(q=d,d-1,\ldots,0\) の順に帰納すれば、
 全てのcoherent \(\mathcal A\)-加群の \(H^q\) が有限生成 \(A\)-加群と分かる。
 この議論は線型化と両立するので、\(\cF\) に有限階数線束をテンソルしたものの
 導来大域切断は \(D^b(\operatorname{gr}^{\Lambda}A)\) に属する。
 ここでcoherentは \(\mathcal A\) 上の意味であり、\(\OO\) 上有限とは仮定しない。

 Beilinsonの対角分解 \cite{Beilinson} を、因子の順序を適切に選んで
 \[
 0\longrightarrow\Omega^d(d)\boxtimes\OO(-d)
 \longrightarrow\cdots\longrightarrow
 \Omega^1(1)\boxtimes\OO(-1)
 \longrightarrow\OO\boxtimes\OO
 \longrightarrow\OO_\Delta\longrightarrow0
 \]
 と書く。これは \(D(\Lambda)\)-線型化を保つ分解である。
 \(\mathcal A\) の係数を保ったまま \(\cF\) に作用させると、\(\cF\) は
 \[
 R\Gamma(\PP^d,\cF\otimes_{\OO}\Omega^i(i))\otimes_k\OO(-i),
 \qquad 0\le i\le d,
 \]
 の有限個のconeから作られる。
 対角のKoszul分解は各射影に関してflatであるから、\(\cF\) の
 \(\OO\)-flatnessは必要ない。非可換な \(A\)-作用も全ての操作と可換である。
 係数複体は前段落により \(D^b(\operatorname{gr}^{\Lambda}A)\) に属する。
 \(t_0\ne0\) に制限すると、\(t_0\) の \(\Lambda\)-次数が0であるため
 \(\OO(-i)\) は次数を保って自明化される。
 よって \(T\) は \(M\otimes P\) から生成される。
 係数複体を有限個のcohomology加群に分けるには通常のtruncationを用いればよい。
 有界複体全体についても同じtruncationから結論する。

 \(P_f\) の場合は補題\ref{lem:graded-localization}による局所化を適用する。
 有限生成次数付き加群の有限表示は局所化前に持ち上がるので、
 局所化後の生成対象はちょうど \(M\otimes P_f\) である。
\end{proof}

\begin{lemma}[多項式中心上の分離的な稠密開集合]\label{lem:polynomial-separable}
 \(D\) を連結局所有限な非負次数付き \(k\)-代数で、有限型可換斉次中心部分代数
 \(Z\) 上有限なものとする。\(d=\dim Z\) と置く。
 正次数の斉次変数を持つ多項式部分代数 \(P=k[t_1,\ldots,t_d]\subset Z\)、
 非零斉次元 \(f\in P\)、斉次両側イデアル \(J\subset D\) を選んで、
 \(D\) が \(P\) 上有限、\(J_f\) が冪零、\(Q=(D/J)_f\) が
 \(P_f\) 上有限射影かつ分離的となるようにできる。
 分離冪等元は元の整数次数で次数0に取れる。
\end{lemma}
\begin{proof}
 斉次Noether正規化によって \(Z\) は正次数変数の多項式部分代数 \(P\) 上有限となる。
 \(Z\) のreducednessは必要ない。\(K=\operatorname{Frac}P\) とすると、
 \(D_K\) は有限次元 \(K\)-代数である。
 そのradicalの \(D\) への逆像を \(J\) とする。
 次数のスカラー作用は \(P\)、\(K\) とradicalを保つので \(J\) は斉次である。
 十分大きい \(a\) で \((J^a)_K=0\) だから、有限生成性によって
 ある非零元を反転すると \(J^a=0\) となる。
 半単純商 \(D_K/\operatorname{rad}(D_K)\) は標数零の体上分離的である。
 generic freenessと分離冪等元の有限個の等式の分母払いにより、
 ある稠密開集合上で \(D/J\) は有限射影かつ分離的となる。

 これらの条件の成立する開集合は次数作用に関して不変である。
 その補集合のイデアルは斉次であり、\(P\) は整域なので、そこから
 非零斉次元 \(f\) を一つ選べる。
 \(d=0\) の場合は \(P=k\)、\(f=1\) とし、有限次元代数のradicalを使う。
 最後に、分離冪等元を元の整数次数で分解し、次数0の成分を取る。
 \(\mu(e)=1\) と \((q\otimes1)e=e(1\otimes q)\) は斉次 \(q\) ごとに次数に分解されるから、
 次数0成分も同じ条件を満たす。これが必要な次数付き分離冪等元である。
\end{proof}

\begin{theorem}[一方だけが有限中心である外部積生成]\label{thm:one-sided-boxes}
 \(A\) を左右noetherianな \(\Lambda\)-次数付き \(k\)-代数とする。
 \(D\) を連結局所有限な非負次数付き \(k\)-代数で、
 有限型可換斉次中心部分代数 \(Z\) 上有限なものとする。
 任意の準同型 \(\beta:\ZZ\to\Lambda\) により \(D\) の次数を粗化し、
 \(A\otimes_kD\) には和の次数を与える。このとき
 \begin{equation}\label{eq:one-sided-boxes}
 \begin{split}
 D^b(\operatorname{gr}^{\Lambda}(A\otimes_kD))
 =\thick\{\,M\boxtimes\operatorname{coarsen}_\beta N:\qquad\qquad\qquad\quad\\
 M\in\operatorname{gr}^{\Lambda}A,\quad
 N\in\operatorname{gr}^{\ZZ}D\,\}.
 \end{split}
 \end{equation}
 特に \(A\) の有限中心性、有限global dimension、あるいは
 左辺の対象の細分次数へのliftは仮定しない。
\end{theorem}
\begin{proof}
 補題\ref{lem:polynomial-separable}の \(P\subset Z\) に対して
 \(A\otimes D\) は \(A\otimes P\) 上左右有限である。
 後者はnoetherianなので前者もnoetherianである。
 右辺のthick部分圏を \(\mathcal B\) とし、\(d=\dim Z\) に関して帰納する。
 \(d=0\) では \(D\) は有限次元であり、連結性から \(D_{>0}\) は冪零で商は \(k\) である。
 \(D_{>0}\)-adic filtrationの各層は有限生成 \(\Lambda\)-次数付き \(A\)-加群であるから、
 \(M\boxtimes k\) からの生成を得る。

 \(d>0\) とし、同補題の \(f,J,Q\) を取る。
 \(f=1\) と取れる場合も以下の開集合の議論はそのまま適用でき、閉集合側は空になる。
 \(J_f\)-adic filtrationにより、\(A\otimes D_f\) 上の全加群は
 \(A\otimes Q\) 上の有限生成加群から生成される。
 そのような加群 \(T\) は \(A\otimes P_f\) 上も有限生成である。
 次数0の分離冪等元 \(e=\sum_\nu x_\nu\otimes y_\nu\) により、乗法写像
 \[
 \operatorname{Res}T\otimes_{P_f}Q\longrightarrow T
 \]
 は \(A\otimes Q\)-線型かつ \(\Lambda\)-次数0のsection
 \(t\mapsto\sum_\nu tx_\nu\otimes y_\nu\) を持つ。
 \(A\) の作用は \(Q\) と可換なので、このsectionは \(A\)-線型でもある。

 補題\ref{lem:polynomial-coefficient}により \(\operatorname{Res}T\) は
 \(M\otimes P_f\) から生成される。
 \(Q\) は \(P_f\) 上射影なので \(-\otimes_{P_f}Q\) はexactであり、
 上のsectionと合わせて \(T\) は \(M\otimes Q\) から生成される。
 \(M\otimes Q\) は \(M\boxtimes(D/J)\) の局所化であるから、
 開集合上の全圏は \(\mathcal B\) の像で生成される。
 ここで \(A\otimes Q\) が有限global dimensionを持つとは主張していない。

 局所化の核は補題\ref{lem:graded-localization}により、
 \(A\otimes(D/(f))\) 上の有限生成次数付き加群から生成される。
 \(D/(f)\) は \(P/(f)\) の像という有限型可換斉次中心上有限である。
 \(f\) が非単元ならこの中心の次元は \(d-1\) 以下であるから、帰納法を適用できる。
 得られる外部積の第二因子は \(D/(f)\)-加群であり、元の \(D\)-加群としても
 元の整数次数を保っている。従って局所化の核は \(\mathcal B\) に含まれる。
 開集合と核の生成をVerdier局所化によって合わせれば \eqref{eq:one-sided-boxes} を得る。
 全てのconeと直和因子は粗化後の圏で取っており、対象の次数のliftは用いない。
\end{proof}

\subsection{任意の粗化次数と有限中心因子の反復}

\begin{theorem}[有限中心代数の外部積生成]\label{thm:finite-center-boxes}
 \(D_1,\ldots,D_r\) を連結局所有限な非負次数付き \(k\)-代数で、
 各 \(D_i\) は有限型可換斉次中心部分代数上有限とする。
 任意の準同型 \(\alpha:\ZZ^r\to\Lambda\) に対し、
 \(D=\bigotimes_iD_i\) に粗化次数を与えると
 \[
 D^b(\operatorname{gr}^{\Lambda}D)
 =\thick\{\operatorname{coarsen}_{\alpha}
 (M_1\boxtimes\cdots\boxtimes M_r)(\lambda):
 M_i\in\operatorname{gr}^{\ZZ}D_i,\ \lambda\in\Lambda\}.
 \tag{\(G_{\mathrm{fin}}\)}\label{eq:finite-center-generation}
 \]
 \(\alpha\) は単射・全射である必要はなく、無限の核を持っていてもよい。
\end{theorem}
\begin{proof}
 \(A=k\) から出発し、\(\beta_i(1)=\alpha(e_i)\) として
 定理\ref{thm:one-sided-boxes}を順に適用する。
 各段階の係数代数は有限型可換中心上有限なので左右noetherianである。
 \(k\) の有限生成 \(\Lambda\)-次数付き加群は有限個の \(k(\lambda)\) の和であり、
 反復によって表示の外部積を得る。
 複体の外部積は体 \(k\) 上exactなので、前段階のthick生成を次段階へ送れる。
\end{proof}

\begin{remark}[nonliftable対象と軸torsion]
 外部積のconeには粗化後の全ての次数零の射を用いる。
 例えば \(k[a,b]/(a^3-b^3)\) は \(3e_1=3e_2\) の次数では
 \(a^3-b^3\) による二つのfree shiftの間のconeで生成される。
 同じ加群が \(\ZZ^2\)-次数へ持ち上がる必要はない。
 また \(f\)-adic filtrationは正次数の中心元に殺される軸上の加群も保持する。
 multigraded tailsへの移行によってこれらを捨てる操作は用いない。
\end{remark}

\subsection{高々一つの無限位数ブロックへの帰結}

次の有限中心性は \cite[Corollary 2.3]{SmithTate} に含まれる既知の結果である。
必要な斉次中心を明示するため、証明を記す。

\begin{lemma}[有限位数THCRの有限中心性]\label{lem:finite-order-thcr}
 \(E\) をsmooth elliptic curve、\(L\) を正次数の線束とし、
 \(\sigma\in\operatorname{Aut}(E)\) が有限位数 \(q\) を持つとする。
 このとき \(B=B(E,L,\sigma)\) は有限型可換斉次中心部分代数上有限である。
\end{lemma}
\begin{proof}
 \(L=\mathcal O_E(D)\) となるdivisor \(D\) を選び、
 \(D_d=D+\sigma^*D+\cdots+(\sigma^{d-1})^*D\) と置く。
 通常のskew Laurent表示により、\(B\) を
 \[
 B_d=H^0(E,\mathcal O_E(D_d))t^d
 \subset k(E)[t,t^{-1};\sigma^*]
 \]
 と書ける。\(\sigma^q=1\) なので \(D_q\) は \(\sigma\)-不変であり、
 \(q\)-Veroneseは可換なsection ring
 \[
 R=B^{(q)}=\bigoplus_{n\ge0}H^0(E,\mathcal O_E(nD_q))t^{qn}
 \]
 になる。\(R\) 上では \(\sigma\) を係数に作用させ、\(t^q\) を固定する。
 \(Z=R^{\langle\sigma\rangle}\) は \(B\) の中心に含まれる。
 実際、\(a\in H^0(E,\mathcal O_E(nD_q))\) が \(\sigma\)-不変なら、
 任意の \(bt^d\in B_d\) に対して
 \[
 (at^{qn})(bt^d)=ab\,t^{qn+d}
 =b\sigma^{d*}(a)t^{qn+d}=(bt^d)(at^{qn})
 \]
 である。

 \(D_q\) は正次数なのでampleである。従って \(R\) は有限型可換代数であり、
 各 \(0\le d<q\) について
 \[
 \bigoplus_{n\ge0}H^0(E,\mathcal O_E(nD_q+D_d))t^{qn+d}
 \]
 は \(R\) 上有限なsection moduleである。
 これらの有限直和が \(B\) なので \(B\) は \(R\) 上有限である。
 有限群の不変式環について \(R\) は有限型環 \(Z\) 上有限であり、
 従って \(B\) も \(Z\) 上有限である。
\end{proof}




\begin{corollary}[一つの任意位数因子を含む生成条件]\label{cor:PI-generation}
 全ての \(j\) で \(h_j^6\ne-27\) とし、並べ替えた後の
 \(\sigma_j=t_{p_j}\ (j\ge2)\) のみ有限位数と仮定する。
 このとき \(C_j=S_{h_j}/(G_j)\)、\(C=\bigotimes_jC_j\) に対して、
 条件 \eqref{eq:block-generation} は成立する。
 中心変数を零とした後に残る \(\Gamma_s\) の全shiftも含む。
\end{corollary}
\begin{proof}
 \(C_1\) はnoetherianであり、その \(\Gamma_s\)-次数は
 単射 \(\ZZ e_1\hookrightarrow\Gamma_s\) によって与えられる。
 任意の有限生成 \(\Gamma_s\)-次数付き \(C_1\)-加群は
 \(\Gamma_s/\ZZ e_1\) の剰余類ごとの部分加群に分解する。
 有限個の生成元の次数が関わる剰余類は有限個だけであるから、
 各剰余類の代表によるshiftを戻せば、元の整数次数を持つ \(C_1\)-加群の有限和となる。
 ここで細分しているのは一因子上の加群だけであり、積の一般の加群ではない。

 \(j\ge2\) では補題\ref{lem:finite-order-thcr}によって \(C_j\) は有限中心上有限である。
 定理\ref{thm:one-sided-boxes}を
 \(A=C_1\)、\(D=C_2\)、\(\beta(1)=e_2\) に適用し、次に係数代数を
 \(C_1\otimes C_2\) として反復する。
 各段階で左右noetherian性は同定理の有限中心多項式拡大によって保証される。
 外部積は体上exactなので、前段階で得たthick生成をそのまま次段階へ送れる。
 これで全ての \(C_j\) について条件 \eqref{eq:block-generation} が得られる。
\end{proof}

以下では定理\ref{thm:multiblock-model}の仮定の下で
\[
 \begin{gathered}
 N=3r+m,\qquad A=\left(\bigotimes_{j=1}^rS_{h_j}\right)[u_1,\ldots,u_m],\\
 W=\sum_jG_{h_j}+P(u),\qquad R=A/(W)=B_{\mathbf h,P}.
 \end{gathered}
\]
と置く。\(m=0\) の場合は \(P=0\) とし、中心変数の因子を省く。

\begin{proof}[定理\ref{thm:multiblock-model}の証明]

 \emph{周囲の代数と中心元の正則性.}
 命題\ref{prop:tensor-sklyanin-ambient}により、\(A\) は左右noetherianな
 Koszul AS-regular代数で、global dimensionとGorenstein parameterは \(N\) である。
 \(W\) は中心正則元である。以下の支持帰着のために、中心上の自由性も具体的に用いる。

 \(C_j\) の斉次基底の持ち上げを \(V_j\subset S_{h_j}\) とすると、
 正則性と次数に関する帰納法により
 \[
 V_j\otimes\CC[t_j]\xrightarrow{\sim}S_{h_j},\qquad
 v\otimes t_j^a\longmapsto vG_j^a
 \]
 は中心 \(\CC[t_j]\)-加群としての同型である。
 従って \(A\) は中心多項式環
 \(Z_0=\CC[t_1,\ldots,t_r,u_1,\ldots,u_m]\) 上自由である。
 \(W\) は \(\sum_jt_j+P(u)\in Z_0\) の作用なので、左右正則である。
 その商 \(R\) は
 \(\CC[t_1,\ldots,t_{r-1},u_1,\ldots,u_m]\) 上自由であり、
 \[
 u_1,\ldots,u_m,G_1,\ldots,G_{r-1}
 \]
 は \(R\) の中心正則列で、商は \(C\) である。

 \emph{全生成性.}
 \(m>0\) では中心多項式環全体 \((T,P)=(\CC[u],P)\) を一つの因子とし、
 \[
 \Gamma=\ZZ^{r+1}/\langle3e_j-3e_0:1\le j\le r\rangle
 \]
 を用いる。\(\deg(u_a)=e_0\)、Sklyanin因子の次数は \(e_j\) である。
 \(m=0\) では \(e_0\) を省き、\(\Gamma=\Gamma_r\) とする。
 finite-free factorizationの微分を \(d\)、第 \(j\) ブロックのEuler導分を
 \(\delta_j\) とすれば、補題\ref{lem:block-euler}と同じ計算により
 \[
 d\delta_j(d)+\delta_j(d)d=3G_j\operatorname{id}.
 \]
 中心変数についても行列係数を微分することで
 \[
 d\,\partial_{u_a}(d)+\partial_{u_a}(d)\,d
 =\partial_{u_a}P\operatorname{id}
 \]
 を得る。これらは斉次のnull-homotopyである。

 \(J_P=(\partial_{u_1}P,\ldots,\partial_{u_m}P)\) とする。
 孤立性より、ある \(b\) について \((u_1,\ldots,u_m)^b\subset J_P\) である。
 命題\ref{prop:block-support}の証明で \(u_a^2\) を \(\partial_{u_a}P\) に置き換えると、
 各対象は \(G_1,\ldots,G_{r-1},\partial_{u_1}P,\ldots,\partial_{u_m}P\)
 に関するKoszul複体のshift・内部twistの直和因子になる。
 各cohomologyの \((u)\)-adic有限filtrationの商は全ての \(u_a\) に殺される。
 それらは \(G_1,\ldots,G_{r-1}\) にも殺され、\(W=0\) と \(P(0)=0\) から
 \(G_r\) にも殺される。従って
 \[
 \Dsg^\Gamma(R)
 =\thick\operatorname{Res}D^b(\operatorname{gr}^\Gamma C).
 \]
 \(m=0\) では中心変数についてのKoszul複体とfiltrationを省けばよい。
 全ての \(\Gamma\)-次数付き \(C\)-加群が
 系\ref{cor:PI-generation}により外部積で生成される。

 この外部積のrestrictionのstabilizationを説明する。
 各 \(C_j=S_{h_j}/(G_j)\) の加群を十分高いsyzygyに置き換えて二項ambient自由分解で表し、
 \(m>0\) では \(T/(P)\) 上の残余体 \(k\) についても同じ操作をする。
 一つのhypersurface因子においてperfectな項は、restriction後に \(R\)-perfectとなる。
 この点を右加群の規約で明示する。選んだambientを \(A_i\)、potentialを \(f_i\)、
 その商を \(D_i\)、他のambientの積を \(A'\)、potentialの和を \(g'\) と書く。
 \(A_i\cong V_i\otimes\CC[f_i]\) と \(f_i=-g'\) の消去により、
 \(R=(A_i\otimes A')/(f_i+g')\) は \(A'\) 上左右とも自由である。
 他因子の外部積 \(M'\) は \(M'g'=0\) を満たし、各ambient上の有限自由分解の
 tensor productによって \(A'\)-perfectである。従って
 \[
 \begin{aligned}
 M'\otL_{A'}R&\simeq M'\otimes_{A'}R
       \xrightarrow{\sim}\operatorname{Res}(D_i\boxtimes M'),\\
 m'\otimes[a\otimes b]&\longmapsto\bar a\otimes m'b
 \end{aligned}
 \]
 は次数を保つ同型であり、右辺は \(R\)-perfectである。
 有限cone、shift、直和因子を取ることで、\(D_i\) を任意の \(D_i\)-perfect複体に
 置き換えても成立する。
 中心因子についても \(m>0\) では孤立性から \(P\ne0\) であり、
 \(T\cong V_0\otimes\CC[P]\) となるので同じ議論を適用できる。
 従ってsyzygyへの置換で失う項は特異圏で零である。

 二項ambient分解の微分を \(d_i\)、potential倍のnull-homotopyを \(s_i\) とする。
 外部tensor積上のKoszul符号付きの和 \(d=\sum d_i\)、\(s=\sum s_i\) は
 \(ds+sd=W\)、\(s^2=0\) を満たす。
 そのShamash周期化のtailはfactorizationの外部積に一致する。
 よって各boxのrestrictionのstabilizationは外部積関手のthick imageに入る。
 全 \(\Gamma\)-shiftも各因子のdegree shiftで実現されるので、
 補題\ref{lem:multiblock-ff}の完全忠実関手はMorita同値である。

 \emph{可換側との比較と整数次数への降下.}
 各Sklyanin因子に系\ref{cor:block}を用い、中心因子には恒等を用いる。
 可換側の外部積には \cite[Theorem 5.15, Corollary 5.18]{BFK} を適用する。
 命題\ref{prop:multiblock-embedding}と同様に、\(\Gamma\to\ZZ\) の有限核の双対による
 同変化と双対同変化を用いる。因子ごとのdegree shiftと \((3)=[2]\) のcoherenceは
 各比較で保存されるので、通常の整数次数へ降下すると
 \[
 \MF^{\ZZ}(A,W)\simeq
 \MF^{\ZZ}\!\left(\CC[x_j,y_j,z_j,u_a],
 \sum_j(x_j^3+y_j^3+z_j^3+h_j^2x_jy_jz_j)+P(u)\right)
 \tag{\(*\)}\label{eq:PI-model}
 \]
 を得る。

 \emph{射影模型と一意性.}
 中心正則元の商 \(R\) はnoetherian AS-Gorensteinで、injective dimensionは \(N-1\)、
 Gorenstein parameterは \(N-3\) である。
 \cite[Theorem 5.8]{CCKM} のMF--特異圏比較と
 \cite[Theorem 2.5]{Orlov} の残余windowを適用できる。
 \(N>3\) ならそのwindowは
 \[
 D^b(\qgr R)=\langle\mathcal K_R,\OO,\OO(1),\ldots,\OO(N-4)\rangle
 \]
 を与え、\(N=3\) なら例外対象はなく \(\mathcal K_R=D^b(\qgr R)\) である。
 degree shiftはwindowを通してrotationに対応する。
 可換potentialの勾配は各変数群に分離しており、各 \(F_{h_j}\) と \(P\) の
 アフィン臨界点は原点だけであるから、その和も原点以外に臨界点を持たない。
 従って \(Y_{\mathbf h,P}\) は滑らかで、同じOrlov比較は可換側を
 \((\Ku(Y_{\mathbf h,P}),\rho_{Y_{\mathbf h,P}})\) と同定する。
 \(N-2\ge3\) における射影同型類の一意性は
 \cite[Theorem 1.7]{LZ} による。
\end{proof}

\begin{lemma}[位数五の明示的パラメータ]\label{lem:order-five}
 \(h^6=5\) とする。原点 \(o=[1:-1:0]\) を持つ
 \(E_h=V(x^3+y^3+z^3+h^2xyz)\) 上で、\(p=[1:-1:-h]\) はちょうど位数五である。
 従って対応するSklyaninの平行移動 \(\sigma=t_p\) も位数五である。
\end{lemma}
\begin{proof}
 \(o\) における接線は \(3x+3y-h^2z=0\) であり、
 この直線上で \(F_h=(1+h^6/27)z^3\) となるので、\(o\) はinflection pointである。
 逆元は \([x:y:z]\mapsto[y:x:z]\) である。実際、この二点と \(o\) は一直線上にある。
 \(a=h^3\) とし、\(q=[a+1:a-1:-2h]\) と置く。
 \(F_h(q)=0\) であり、\(p\) における接線
 \[
 (3+a)x+(3-a)y+2h^2z=0
 \]
 は \(-q=[a-1:a+1:-2h]\) を通る。
 \(h\ne0\) では \(-q\ne p\) であるから、接線の加法則により \(q=2p\) である。
 \(q\) における接線の係数は
 \[
 [a^2+8a+3:a^2-8a+3:h^2(a^2+11)]
 \]
 であり、これを \(p\) で評価すると \(a(5-a^2)=h^3(5-h^6)\) となる。
 \(p\ne q\) は \(q=2p\) と \(p\ne o\) から従う。
 \(h^6=5\) では \(p\) がこの接線上にあるので \(2q+p=0\)、従って \(5p=0\) である。
 \(p\ne o\) かつ五は素数なので、位数はちょうど五である。
\end{proof}

\begin{remark}[二・三ブロックのsixfold・sevenfold]
 全ての \(h_j^6=5\) とする。
 \(r=2,m=2,P=u_1^3+u_2^3\) は二つの非可換ブロックを持つcubic sixfoldの
 rotation付き残余圏の完全な可換模型を与える。
 \(r=3,m=0\) は三つの非可換ブロックを持つcubic sevenfoldの完全な可換模型を与える。
 両者の方程式は定理\ref{thm:multiblock-model}の \(Y_{\mathbf h,P}\) であり、
 元の対象の一部分ではなく全残余圏を実現する。
\end{remark}

\begin{lemma}[無限位数の明示的パラメータ]\label{lem:order-infinite}
 \(h=2\) とする。楕円曲線
 \[
 E_2=V(x^3+y^3+z^3+4xyz),\qquad o=[1:-1:0]
 \]
 上の点 \(p_2=[1:-1:-2]\) は無限位数である。
\end{lemma}
\begin{proof}
 曲線と点は \(\mathbb Q\) 上定義されている。
 \(2^6+27=91\) は5,11のどちらでも割れないので、Hesseの勾配による判定から
 5と11で良い還元を持つ。
 \(z=1\)、\(z=0,y=1\)、\([1:0:0]\) の三つのchartに分けて列挙すると
 \[
 \#E_2(\mathbb F_5)=9,\qquad \#E_2(\mathbb F_{11})=12.
 \]
 実際、\(z=1\) の点数はそれぞれ8と11、\(z=0\) にはどちらも一点だけ存在する。
 この列挙は整数の合同計算で再現でき、\texttt{verify\_algebra.py} に全座標を記録する。
 良い還元では剰余標数と素なtorsionの還元写像は単射である
 \cite[Reduction of torsion points]{SnowdenDVR}。
 従って \(p_2\) の位数 \(N\) が有限なら、\(N\) の5と素な部分は9を割り、
 11と素な部分は12を割る。後者により5-primary部分はなく、前者により11-primary部分もなく、
 残りの素数について両方を比較すれば \(N\mid3\) となる。

 \(p_2\ne o\) である。一方、接線による倍点公式から
 \[
 2p_2=[9:7:-4],\qquad -p_2=[1:-1:2]
 \]
 であり、これらは異なる射影点である。
 従って \(3p_2\ne0\) であり、\(N\mid3\) は不可能である。
 \(p_2\) は無限位数である。
\end{proof}

\begin{remark}[真に無限位数を含むsixfoldとsevenfold]
 \(h_1=2\)、\(h_j=5^{1/6}\ (j\ge2)\) と取る。
 第一ブロックの平行移動は無限位数、他のブロックは位数5である。
 \((r,m)=(2,2)\)、\(P=u_1^3+u_2^3\) では可換sixfoldは
 \[
 V\bigl(x_1^3+y_1^3+z_1^3+4x_1y_1z_1
       +x_2^3+y_2^3+z_2^3+5^{1/3}x_2y_2z_2+u_1^3+u_2^3\bigr)
 \subset\PP^7
 \]
 となる。\((r,m)=(3,0)\) では第二・第三ブロックの係数を \(5^{1/3}\) とした
 cubic sevenfoldを得る。いずれも全残余圏とrotationの同値である。
 さらに他の有限位数パラメータを固定すれば、\(h_1^6\ne-27\) を満たす
 全ての \(h_1\in\CC\) について同じ結論が成り立つ。
 これは各ファイバーでの主張であり、同値の族をこのパラメータ直線全体に
 代数的に降下させたという主張ではない。
\end{remark}

\subsection{証明範囲と残る課題}

有限位数条件は \(j\ge2\) のみに用いる。
従ってより一般に、\(C_1\) がnoetherian、\(j\ge2\) の各 \(C_j\) が
有限型可換斉次中心上有限で、周囲が命題\ref{prop:block-support}の仮定を満たせば、
同じ生成と安定化の議論を適用できる。
定理\ref{thm:one-sided-boxes}では任意のnoetherianな係数代数を保てる一方、
新たに加える因子の有限中心性は必要である。
従って二つ以上の独立な無限位数THCRを持つ場合は、この証明だけでは解けない。
これはその場合の非存在を意味せず、条件 (G) が未解決であるという意味である。

次の焦点は、二つの無限位数因子の条件 (G) を証明することである。
それが得られれば、残りの有限中心因子は定理\ref{thm:one-sided-boxes}で追加できる。
必要なのは、THCRに対しても任意のnoetherian係数を許す外部積生成、
または粗化次数を保った対角の有限構成である。
可換ambientのThom--Sebastiani定理を非可換ambientへ無条件に適用したわけではない。
先行研究上の新規性は未確認である。

\section{適用範囲と残る問題}\label{sec:scope}

\subsection{例外パラメータでpropernessが失われる}

\begin{proof}[命題\ref{prop:bad-intro}の証明]
 \(h_0^6=-27\) とする。\(E_{h_0}\) は三本の直線からなる三角形である。
 まず
 \[
 S_{h_0}/(G_{h_0})\simeq B(E_{h_0},L,\sigma_{h_0})                 \tag{6}\label{eq:triangle}
 \]
 を確認する。楕円曲線の場合だけの公式を、この特殊値へ無条件に適用しない。
 \(h_0\) の近傍で、関係式を多重線形化した行列
 \[
 M_h(x,y,z)=
 \begin{pmatrix}-y&x&-hz\\-hx&-z&y\\z&-hy&-x\end{pmatrix}
 \]
 は \(E_h\) 上でrank二である。実際その行列式は
 \(-h(x^3+y^3+z^3+h^2xyz)\) であり、rank一なら小行列式から
 \(yz+hx^2=yz-hx^2=0\) と巡回的な式を得て、\(x=y=z=0\) となる。
 核は相対的な射 \(\sigma_h\) を与える。グラフを転置すると \(h\) を \(-h\) に
 変える同じ関係式になるので、この射は相対自己同型である。

 \(\mathcal L_d=\bigotimes_{j=0}^{d-1}(\sigma^j)^*\OO_E(1)\) は、
 滑らかなファイバーでは次数 \(3d\)、三角形では各成分次数 \(d\) を持つ。
 \(d>0\) なら \(H^1(E_h,\mathcal L_d)=0\)、\(h^0=3d\) であり、
 相対切断加群はlocally freeで底変更と両立する。
 自然な次数付き写像 \(S_h\to\bigoplus_dH^0(E_h,\mathcal L_d)\) の下で、
 \(G_h\) は一般の滑らかなファイバーで零なので、局所自由性から恒等的に零となる。
 三角形のtwisted homogeneous coordinate ringが次数一で生成されることを確認する。
 \(L\) のbasepoint-free pencilに対するpencil trickを使うと、
 多次数 \((d,d,d)\) の \(M\) について
 \(H^0(L)\otimes H^0(M)\to H^0(LM)\) は
 \(H^1(ML^{-1})=0\) なら全射である。
 \(d\ge2\) では各成分上の正値性からこの消滅が従う。
 \(d=1\) で \(ML^{-1}\) が非自明なら、多次数零の直線束の切断は
 三角形の三成分の接合条件により零で、Euler数も零なので同じ消滅が成り立つ。
 \(M\simeq L\) なら、平面三次曲線の斉次イデアルに二次式がないことから全射である。
 従って \(S_{h_0}/(G_{h_0})\to B(E_{h_0},L,\sigma_{h_0})\) は全射である。
 補題\ref{lem:center}の正則性から左辺のHilbert係数は \(1,3d\) であり、
 右辺の係数と一致する。よってこの全射は同型で、\eqref{eq:triangle}が従う。
 \(\mathcal L_d\) の各成分上の次数が \(d\) とともに増大するので
 \(\sigma\)-amplenessも成り立ち、tailsは \(\operatorname{coh}E_{h_0}\) と同値になる。

 Gorenstein parameter零のOrlov同値と \cite[Theorem 5.8]{CCKM}から
 \[
 \MF^\ZZ(S_{h_0},G_{h_0})\simeq D^b(\operatorname{coh}E_{h_0}).
 \]
 三角形の節点 \(p\) のskyscraper sheafは、任意に高い次数で非零の自己Extを持つ。
 節点の完備局所環 \(\CC[[s,t]]/(st)\) の剰余体の最小自由分解は無限であり、
 局所大域Extスペクトル系列では高次Extの支持が \(p\) だけなので、
 この非零性は大域自己Extにも残る。従ってこのdg圏はproperでない。

 \(m=0\) の場合の非propernessはこれで示された。以下では \(m\ge1\) とする。
 \(P=\sum_i u_i^3\) とし、第\ref{sec:TS}節の二ブロック次数 \(\Gamma\) と
 Koszul factorization \(K_P\) を使う。
 \(K_P\) の全微分行列は \((u_1,\ldots,u_m)\) に属する。
 もし恒等写像が零ホモトピーなら、係数を原点で評価して非零自由加群上の恒等写像が
 零となるので矛盾する。従って \(H^0\REnd(K_P)\ne0\) である。
 節点の対象を \(N\) とすると、外部積の完全忠実性と体上のKünneth公式により
 \[
 H^r\REnd(N)\otimes H^0\REnd(K_P)
 \subset H^r\REnd(N\boxtimes K_P)
 \]
 が直和因子として残る。よって \(\Gamma\)-次数の全体圏はnonproperである。
 もし通常の整数次数付き圏がproperなら、その有限群 \(H_\Gamma\simeq\mu_3\) による
 同変圏も、Hom複体の不変部分が直和因子となることからproperとなる。
 これは式\eqref{eq:eq}に反する。

 さらに、この議論は単なるHom総次元の非有限性より強い。
 \(\Gamma\)-次数の一つの対象に、任意に高い正次数の自己Homが残っている。
 有限群同変化から忘却すると、同変Homは通常のHomの直和因子なので、
 通常次数でも同じ正次数非消滅を持つ一つの対象が得られる。
 \(m>0\) ではOrlovの完全忠実windowによりこれを \(D^b(\qgr B_{h_0,m})\) へ送れ、
 \(m=0\) では上の三角形との同値を用いる。
 もし \(\qgr B_{h_0,m}\) のhomological dimensionが有限なら、
 その有界導来圏の任意の固定された二対象の間の正次数Homは十分高次で消える。
 これは有界複体のcohomologyとExtからなる有限スペクトル系列で分かる。
 今得た対象はこれに反するから、homological dimensionは無限である。

 一方、周囲の \(S_{h_0}[u]\) はAS-regularであり、三次の中心元の正則性は
 \(m>0\) なら最後の中心変数についてmonicであること、\(m=0\) なら
 補題\ref{lem:center}から従う。従ってHilbert seriesは依然として
 \((1-t^3)/(1-t)^{m+3}\) であり、degree shiftの三乗も \([2]\) である。
 滑らかな射影多様体のadmissible成分はproperなので、求める非存在が従う。
\end{proof}

この例は、量子化を構成する開集合をさらに縮めて例外値を除くことを妨げない。
従って、零を含む十分小さな開集合での存在問題に対する反例ではない。

\subsection{量子化原稿への適用}

三次の中心元によるAS-regularな周囲の商
\[
 B_\hbar=A_\hbar/(\Omega_\hbar),\qquad B_0=\CC[x]/(F)
\]
を、次数ごとにflatな族として構成できていれば、
その形式完備化には定理\ref{thm:main}が適用できる。
この適用にはPoisson構造の分類を再度用いず、各Artin次数でのgraded flatnessだけを使う。
従って、量子化の構成自体の証明と、本稿の相対残余圏・復元の証明は独立に検証できる。
第\ref{pthm:quantization}定理への適用では、そこで構成する次数ごとのflatnessを用いる。
代数の量子化と圏の復元は異なる段階であり、それぞれの仮定を区別する。

形式的な可換モデルは、半古典極限の固定された \(X\) と同じとは限らない。
定理\ref{thm:global}の代表式では \(h^2xyz\) が現れる。
このモデルを初期の \(F\) と取り違えない。

\subsection{一般の有限パラメータと次の核心}

定理\ref{thm:etale}により、smooth properな代数的pairの族が与えられれば、
可換cubic点のétale近傍で存在問題は解ける。
一般の量子化族に残る最初の課題は、その相対saturationの構成である。
各Artin切断のsmooth proper性だけから、有限表示の代数的模型の存在を結論しない。

独立な複数ブロックでは、命題\ref{prop:multiblock-embedding}で得た完全忠実な像が
全MF圏を生成するかが、二つ以上の無限位数因子の場合の核心である。
命題\ref{prop:block-support}と条件\eqref{eq:block-generation}は、その具体的な帰着を与える。
高々一つの無限位数因子を持つ場合は定理\ref{thm:multiblock-model}で解決した。
中心三次式内の混合項は定理\ref{thm:global}と定理\ref{thm:multiblock-model}に含まれるが、
非可換ブロックと中心変数の混合項は含まれない。
また、étale局所モデルの元の底への降下と、実現開部分の外のファイバーの分類は
別の大域問題として残る。


\section{相対saturationへの二つの還元}
\label{sec:relative-saturation-reduction}
本節では滑らかな可換cubicを含むPoisson量子化族の相対saturationを、
全超曲面圏の対角の有限証明書へ還元する。
相対Beilinson仮定(RB1)--(RB2)、相対例外列、相対Orlov比較は
元論文の量子化定理からは自動的に従わない。
gluingに関する以下の命題は既知の一般原理の相対版であり、
新規な分類定理とは主張しない\cite[Theorem 3.25]{OrlovGlue}。

\subsection{例外列を除く前後の相対smoothness}
\begin{proposition}[相対smoothnessの同値]
\label{prop:saturation-reduction}
$R$ を可換Noether環、$\mathcal C$ を小さいidempotent completeな
$R$-linear dg圏とし、
$\RHom_{\mathcal C}(X,Y)\in\Perf(R)$ が任意の対象$X,Y$に対し成り立つとする。
有限個の相対例外対象 $E_1,\ldots,E_r$ と
dg半直交分解
\[
\mathcal C=\langle\mathcal K,E_1,\ldots,E_r\rangle
\]
が与えられ、さらに$\mathcal K$がcompact generatorを持つならば、
\[
\mathcal C\text{ は }R\text{-smooth}
\quad\Longleftrightarrow\quad
\mathcal K\text{ は }R\text{-smooth}.
\]
ゆえに相対propernessが既に分かっていれば、
両者の相対saturationは同値な条件である。
\end{proposition}
\begin{proof}
例外対象を一つずつ追加する帰納法により一個の場合に帰着する。
$\mathcal K$ のcompact generatorを$G$、例外対象を$E$とする。
この半直交分解のdg enhancementは$\mathcal K$と$\Perf(R)$を
\[
M=\RHom_{\mathcal C}(G,E)
\]
により三角dg gluingしたものとして表される（右左の規約によるoppositeを適宜取る）。
$\Gamma=\REnd_{\mathcal K}(G)$ と置けば、propernessより
$M$は$R$-perfectである。
$\mathcal K$がsmoothならば
$\Gamma\in\Perf(\Gamma^{\mathrm{op}}\otL_R\Gamma)$。
恒等式 $M\simeq M\otL_\Gamma\Gamma$ に対角の有限cell表示を適用すると、
その像は $M\otL_R\Gamma$ の有限回のshift、cone、直和因子から構成される。
$M$が$R$-perfectだから$M$は$\Gamma$-perfect、従って
gluing bimoduleがperfectとなる。三角dg圏の対角の標準triangleにより
$\mathcal C$はsmoothとなる。
逆に$\mathcal C$がsmoothなら、同じgluingの対角triangleを
両成分へ制限して$\mathcal K$の対角がperfectである。
以上は有限cell・導来テンソル・直和因子のみを使うので
$R$上でも成立する。体上の定理は\cite[Theorem 3.25]{OrlovGlue}。
\end{proof}

\begin{corollary}[必要な対角証明書は一つでよい]
\label{cor:only-ambient-diagonal}
量子化 $B_R=A_R/(\Omega_R)$ が研究ノートの相対Beilinson仮定
(RB1)--(RB2)と相対例外列を満たし、
$E=\bigoplus_{i=0}^{N-1}\mathcal O_{B_R}(i)$ がcompact generatorであるとする。
$\Lambda_R=\REnd(E)$は$R$-perfectであり、
その残余圏$\mathcal K_R$は、次の条件と同値に$R$-smoothである：
\[
\Lambda_R\in
\Perf(\Lambda_R^{\mathrm{op}}\otL_R\Lambda_R).
\tag{S}
\]
残余圏のEnd代数$\Gamma_R$と$\Lambda_R$の二つの対角を
別々に持ち上げる必要はない。
\end{corollary}
\begin{proof}
研究ノートの相対proper模型によって$\Perf(\Lambda_R)$はproperで、
$\mathcal K_R$と有限相対例外列による半直交分解を持つ。
命題\ref{prop:saturation-reduction}を適用する。
\end{proof}

\begin{remark}[閉ファイバーのsmoothnessだけでは足りない]
$R=\CC[\hbar]_{(\hbar)}$とし、
$\mathcal C_0\simeq\Perf(X)$が滑らかな可換cubicの圏でも、
その対角の有限cell分解を$R$へ持ち上げることは別の問題である。
$P\in\Perf(\Lambda_R^{\mathrm{op}}\otL_R\Lambda_R)$と
$P\to\Lambda_R$を構成し、$\hbar=0$で擬同型であることを証明できれば、
coneの$R$-perfect性とsupportの閉性により
零を含むZariski開近傍で(S)が成り立つ。
この有限証明書はまだ一般量子化族について構成されていない。
\end{remark}

\subsection{graded matrix factorizationの三次周期性}
\begin{lemma}[中心三次元の形式的周期性]
\label{lem:mf-three-period}
$A$を可換環$R$上の整数次数付き結合代数、
$\Omega\in Z(A)_3$を中心元とする。
有限射影次数付き加群のgraded matrix factorizationの
内部次数シフトを$(1)$とすると、自然なdg自己同値の同型
\[
(1)^3\simeq[2]
\]
があり、任意の基底変更と両立する。
\end{lemma}
\begin{proof}
因子分解を
\[
P_0\xrightarrow{d_0}P_1\xrightarrow{d_1}P_0(3),
\quad
d_1d_0=\Omega,\quad d_0(3)d_1=\Omega
\]
と書く。懸垂$[1]$は二項を入れ替え、片側を$(3)$移し、
二つの微分に負号を付ける。二回の懸垂では負号が消え、
両項および微分が一斉に$(3)$移る。
従って$[2]\simeq(3)=(1)^3$である。
次数移動と符号のみを用いるので基底変更とも可換である。
\end{proof}

\begin{corollary}[相対Orlov比較の下でのrotation]
\label{cor:relative-rotation-period}
$\mathcal K_R$と$\MF^{\mathrm{gr}}(A_R,\Omega_R)$の間に、
基底変更と両立する相対Orlov同値があり、
$(1)$が標準rotation$\rho_R$に対応すると仮定する。
すると$\rho_R$は自己同値で、相対的に
\[
\rho_R^3\simeq[2]
\]
が成立する。この段階では相対smoothnessを別途仮定しない。
\end{corollary}
\begin{proof}
補題\ref{lem:mf-three-period}を相対同値に沿って移す。
\end{proof}

\subsection{可換復元に向けた残る問題}
上の議論により、少なくとも次の四点を区別できる：
(RB1)--(RB2)を指定量子化族について証明すること、
全超曲面圏$\Perf(\Lambda_R)$の対角の有限証明書を可換点から持ち上げること、
相対Orlov比較でrotationを同定すること、
そして得られたsmooth proper pairが可換cubicのpairの代数的変形像に入ること。
前の三点が解けても最後の可換モデルの存在は別問題であり、
Zhang twistの障害を無視することはできない。


\begin{thebibliography}{99}
\bibitem{PHartshorne}
R.~Hartshorne,
\emph{Algebraic Geometry}, Graduate Texts in Mathematics 52,
Springer, 1977.
\href{https://doi.org/10.1007/978-1-4757-3849-0}{doi:10.1007/978-1-4757-3849-0}.

\bibitem{PMayanskiy}
E.~Mayanskiy,
Poisson cohomology of two Fano threefolds,
\emph{Journal of Algebra} \textbf{424} (2015), 21--45.
\href{https://doi.org/10.1016/j.jalgebra.2014.08.049}{doi:10.1016/j.jalgebra.2014.08.049};
\href{https://arxiv.org/abs/1303.5916}{arXiv:1303.5916}.
本文で用いる Grassmann 記述は Theorem~1。

\bibitem{PEtingofLie}
P.~Etingof,
\emph{Lie Groups and Lie Algebras}, MIT lecture notes, 2024 course version.
\href{https://live.ocw.mit.edu/courses/18-755-lie-groups-and-lie-algebras-ii-spring-2024/mit18_755_s24_lec_full.pdf}{MIT OpenCourseWare 公開原稿}.
本文では Proposition~16.9、\S19.4、Theorem~25.17、\S32.2 を参照する。

\bibitem{PMoriSmith}
I.~Mori and S.~P.~Smith,
The classification of 3-Calabi--Yau algebras with 3 generators and 3 quadratic relations,
\href{https://arxiv.org/abs/1502.07403v2}{arXiv:1502.07403v2} (2016).
Theorems~1.3, 1.4, 4.2 and Corollaries~1.6, 2.7, 2.9.

\bibitem{PBergerTaillefer}
R.~Berger and R.~Taillefer,
Poincar\'e--Birkhoff--Witt deformations of Calabi--Yau algebras,
\emph{Journal of Noncommutative Geometry} \textbf{1} (2007), 241--270.
\href{https://arxiv.org/abs/math/0610112}{arXiv:math/0610112}.
本文で用いる potential による PBW 変形は Theorem~1.1(i)。

\bibitem{PKontsevich}
M.~Kontsevich,
Deformation quantization of Poisson manifolds, I,
\href{https://arxiv.org/abs/q-alg/9709040}{arXiv:q-alg/9709040} (1997).
本文では \S4.6 および \S6 の affine formality を用いる。

\bibitem{PDolgushev}
V.~Dolgushev,
The Van den Bergh duality and the modular symmetry of a Poisson variety,
\href{https://arxiv.org/abs/math/0612288v3}{arXiv:math/0612288v3} (2007).
Theorem~3 および Proposition~1。

\bibitem{Beilinson}
A. A. Beilinson,
\emph{Coherent sheaves on \(\PP^n\) and problems of linear algebra},
Functional Analysis and Its Applications \textbf{12} (1978), 214--216.
\url{https://doi.org/10.1007/BF01681436}.
原著Section 4の対角Koszul分解を確認。
\url{https://www.mathnet.ru/eng/faa2009}.
\bibitem{DDL}
O. De Deken and W. Lowen,
\emph{Abelian and derived deformations in the presence of \(\ZZ\)-generating geometric helices},
Journal of Noncommutative Geometry \textbf{5} (2011), 477--505.
Theorem 4.5; Propositions 5.5, 5.8, 5.9; Lemma 5.14.
\url{https://doi.org/10.4171/JNCG/83}.
\bibitem{Keller}
B. Keller, \emph{Deriving DG categories},
Annales scientifiques de l'École Normale Supérieure \textbf{27} (1994), 63--102.
\url{https://doi.org/10.24033/asens.1689}.
\bibitem{Mathew}
A. Mathew, \emph{The Galois group of a stable homotopy theory},
Advances in Mathematics \textbf{291} (2016), 403--541.
Proposition 3.28, Remark 3.29, Proposition 3.35.
\url{https://arxiv.org/abs/1404.2156}.
\bibitem{Orlov}
D. Orlov,
\emph{Derived categories of coherent sheaves and triangulated categories of singularities},
arXiv:math/0503632, v3 (2009), Theorem 2.5.
\url{https://arxiv.org/abs/math/0503632}.
\bibitem{CCKM}
T. Cassidy, A. Conner, E. Kirkman and W. F. Moore,
\emph{Periodic free resolutions from twisted matrix factorizations},
Journal of Algebra \textbf{455} (2016), 137--163.
Proposition 5.7, Theorem 5.8.
\url{https://doi.org/10.1016/j.jalgebra.2016.01.037}.
\bibitem{CT}
A. Căldăraru and J. Tu, \emph{Curved \(A_\infty\) algebras and Landau--Ginzburg models},
New York Journal of Mathematics \textbf{19} (2013), 305--342.
Section 5、特にdg Orlov windowのTheorem 5.8.
\url{https://nyjm.albany.edu/j/2013/19-16v.pdf}.
\bibitem{BFK}
M. Ballard, D. Favero and L. Katzarkov,
\emph{A category of kernels for equivariant factorizations and its implications for Hodge theory},
Publications Mathématiques de l'IHÉS \textbf{120} (2014), 1--111.
Theorem 5.39のKoszul cohomologyによる公式、Theorem 5.15とCorollary 5.18のequivariant外部積.
\url{https://doi.org/10.1007/s10240-013-0059-9}.
確認した公刊本文 \url{https://www.numdam.org/article/PMIHES_2014__120__1_0.pdf}.
\bibitem{LZ}
X. Lin and S. Zhang,
\emph{Serre algebra, matrix factorization and categorical Torelli theorem for hypersurfaces},
arXiv:2310.09927, v2 (2024), Theorem 1.7.
\url{https://arxiv.org/abs/2310.09927}.
\bibitem{HR}
D. Huybrechts and J. V. Rennemo,
\emph{Hochschild cohomology versus the Jacobian ring, and the Torelli theorem for cubic fourfolds},
Algebraic Geometry \textbf{6} (2019), 76--99.
Lemma 2.19およびSections 2--3.
\url{https://doi.org/10.14231/AG-2019-005}.
\bibitem{Huy}
D. Huybrechts, \emph{The K3 category of a cubic fourfold},
Compositio Mathematica \textbf{153} (2017), 586--620.
Proposition 3.12およびSection 3.5.
\url{https://doi.org/10.1112/S0010437X16008137}.
\bibitem{JLLZ}
A. Jacovskis, X. Lin, Z. Liu and S. Zhang,
\emph{Infinitesimal categorical Torelli theorems for Fano threefolds},
arXiv:2203.08187, Theorem 4.4(2).
\url{https://arxiv.org/abs/2203.08187}.
\bibitem{Lurie}
J. Lurie, \emph{Derived Algebraic Geometry X: Formal Moduli Problems},
2011. Theorems 5.1.9, 5.3.16およびProposition 5.3.21.
\url{https://www.math.ias.edu/~lurie/papers/DAG-X.pdf}.
\bibitem{BKP}
A. Blanc, L. Katzarkov and P. Pandit,
\emph{Generators in formal deformations of categories},
Compositio Mathematica \textbf{154} (2018).
本文で用いるのはarXiv版のDefinition 3.3およびSection 3.3.
\url{https://arxiv.org/abs/1705.00655}.
\bibitem{KL}
B. Keller and W. Lowen,
\emph{On Hochschild cohomology and Morita deformations},
International Mathematics Research Notices (2009).
\url{https://webusers.imj-prg.fr/~bernhard.keller/publ/KellerLowenHochDefo.pdf}.
\bibitem{EG}
P. Etingof and V. Ginzburg,
\emph{Noncommutative del Pezzo surfaces and Calabi--Yau algebras},
Journal of the European Mathematical Society \textbf{12} (2010), 1371--1416.
Sections 3.5--3.6, Theorem 3.6.2, Proposition 8.1.1(i) およびその証明。
\url{https://doi.org/10.4171/JEMS/235}.
\href{https://arxiv.org/abs/0709.3593v4}{arXiv:0709.3593v4}.
\bibitem{Walton}
C. Walton, \emph{Representation theory of three-dimensional Sklyanin algebras},
Nuclear Physics B \textbf{860} (2012), 167--185.
Theorem 2.14および式(2.15)--(2.16).
\url{https://doi.org/10.1016/j.nuclphysb.2012.02.015}.
\bibitem{CGWZ}
K. Chan, J. Gaddis, R. Won and J. J. Zhang,
\emph{Ozone groups of Artin--Schelter regular algebras satisfying a polynomial identity},
arXiv:2312.13014, v2 (2024). Introductionの非退化Sklyanin代数の記述。
\url{https://arxiv.org/abs/2312.13014}.
\bibitem{Elagin}
A. Elagin, \emph{On equivariant triangulated categories},
arXiv:1403.7027, Theorem 4.2.
\url{https://arxiv.org/abs/1403.7027}.
\bibitem{ELS}
A. Elagin, V. A. Lunts and O. M. Schn\"urer,
\emph{Smoothness of Derived Categories of Algebras},
Moscow Mathematical Journal \textbf{20} (2020), no. 2, 277--309.
確認した一次資料はarXiv:1810.07626v2 (22 March 2019),
Theorem 4.4 (Verdier局所化), Proposition 4.10とTheorem 4.15 (稠密開集合と生成子),
Theorem 4.18 (外部積生成), Theorem 5.1 (有限中心代数の有界導来圏のsmoothness).
\url{https://arxiv.org/abs/1810.07626};
\url{https://www.mathjournals.org/mmj/2020-020-002/2020-020-002-003.html}.

\bibitem{TV}
B. Toën and M. Vaquié,
\emph{Moduli of objects in dg-categories},
Annales scientifiques de l'École Normale Supérieure (4) \textbf{40} (2007), 387--444.
Theorem 3.6; Corollaries 2.13, 3.15, 3.19, 3.21; Lemma 3.25 and proof of Corollary 3.24.
\url{https://arxiv.org/abs/math/0503269};
確認した著者版 \url{https://www.math.univ-toulouse.fr/~btoen/ttt.pdf}.
\bibitem{Stacks}
The Stacks Project Authors, \emph{The Stacks Project}.
Tag 02HX (Lemma 37.12.1, Noetherian底上の無限小smoothness判定),
Tag 055U (Lemma 37.38.6, smooth morphismのétale局所section).
\url{https://stacks.math.columbia.edu/tag/02HX};
\url{https://stacks.math.columbia.edu/tag/055U}.
参照日2026年11月4日。
\bibitem{Pirozhkov}
D. Pirozhkov, \emph{A categorical Torelli theorem for hypersurfaces},
Bulletin of the London Mathematical Society \textbf{56} (2024), 3075--3089.
\url{https://doi.org/10.1112/blms.13117}.
公刊版のabstractで、cubicに対する仮定 \(n>3\) を確認。
本稿の全次元の一意性には \cite[Theorem 1.7]{LZ} を用いる。
\bibitem{HuyUpdate}
D. Huybrechts, \emph{The K3 category of a cubic fourfold: an update},
Beiträge zur Algebra und Geometrie / Contributions to Algebra and Geometry (2025).
Section 9: Unknown deformations.
\url{https://doi.org/10.1007/s13366-025-00788-9}.

\bibitem{ChanMulti}
D. Chan, \emph{Twisted multi-homogeneous coordinate rings},
Journal of Algebra \textbf{223} (2000), 438--456.
Theorem 3.4; proof of Theorem 5.2, Claim 5.3 and Lemmas 5.4--5.6;
Corollaries 5.8--5.9.
\url{https://doi.org/10.1006/jabr.1999.7938}.
確認した著者版 \url{https://web.maths.unsw.edu.au/~danielch/paper/tbcr.ps}, pp.8--9,14--17.
\bibitem{SmithTate}
S. P. Smith and J. Tate,
\emph{The center of the 3-dimensional and 4-dimensional Sklyanin algebras},
K-Theory \textbf{8} (1994), 19--63.
Corollary 2.3 (p.~26); Theorem 3.7(c) (p.~38).
\url{https://doi.org/10.1007/BF00962089};
\url{https://sites.math.washington.edu/~smith/Research/tate.pdf}.
\bibitem{KeelerMulti}
D. S. Keeler, \emph{Noncommutative ampleness for multiple divisors},
Journal of Algebra \textbf{265} (2003), no. 1, 299--311.
Theorem 2.7, Corollary 2.8 (任意有限個のbimoduleのNC-ampleness);
Hypothesis 3.1, Theorem 3.3, Corollary 3.5 (二因子のnoetherian性).
番号はarXiv一次TeXソースで確認し、刊行版との番号一致は未照合。
\url{https://doi.org/10.1016/S0021-8693(03)00126-1};
\url{https://arxiv.org/abs/math/0210417};
\url{https://arxiv.org/e-print/math/0210417}.
\bibitem{WWY3}
C. Walton, X. Wang and M. Yakimov,
\emph{Poisson geometry of PI three-dimensional Sklyanin algebras},
Proceedings of the London Mathematical Society \textbf{118} (2019), 1471--1500.
Theorem 1.3(1),(2a), verified in arXiv:1704.04975v3, p.~3.
\url{https://doi.org/10.1112/plms.12220};
\url{https://arxiv.org/abs/1704.04975}.
\bibitem{ToenFinitude}
B. Toën, \emph{Finitude homotopique des dg-algèbres propres et lisses},
Proceedings of the London Mathematical Society (3) \textbf{98} (2009), 217--240.
Lemmas 3.13--3.14. Lemma 3.14の結論は部分前層のquasi-compactnessであり、一般の開性ではない。
\url{https://doi.org/10.1112/plms/pdn027};
確認した著者版 \url{https://www.math-beb.fr/wp-content/uploads/2026/10/finitesatdgalg.pdf}.

\bibitem{SnowdenDVR}
A. Snowden, \emph{Lecture 8: Elliptic curves over DVRs},
Math 679 lecture notes, University of Michigan, 2013.
Section ``Reduction of torsion points'', good reductionでのprime-to-residue-characteristic torsionの還元の単射性。
\url{https://websites.umich.edu/~asnowden/teaching/2013/679/L08.html}.
参照日2026年11月6日。
\bibitem{HMY}
J.-W. He, X.-C. Ma and Y. Ye, \emph{Generalized Kn\"orrer's Periodicity Theorem},
arXiv:2107.06438v2 (16 August 2021).
対象は中心二次元の商であり、本稿の三次の場合には直接適用しない。
\url{https://arxiv.org/abs/2107.06438}.

\bibitem{OrlovGlue}
D. Orlov,
\emph{Smooth and proper noncommutative schemes and gluing of DG categories},
Advances in Mathematics \textbf{302} (2016), 59--105.
Theorem 3.25 (smoothness of dg gluing; Proposition 3.24はsmoothnessからregularityのみ)。
\url{https://doi.org/10.1016/j.aim.2016.07.014};
\url{https://arxiv.org/abs/1402.7364}.

% Current verified additions, 2026-11-08.
% Insert inside the manuscript's thebibliography environment.

\bibitem{UeyamaIsolated}
K. Ueyama,
\emph{Noncommutative graded Gorenstein isolated singularities},
in \emph{Proceedings of the 45th Symposium on Ring Theory and Representation Theory}
(Shinshu University, September 7--9, 2012),
K. Koike (ed.), February 2013, 158--163.
Theorem 7, printed p.~161.
\url{https://ring-theory-japan.com/ring/oldmeeting/2012/report2012/proceedings.pdf}.

\bibitem{ZhangTwist}
J.~J. Zhang,
\emph{Twisted graded algebras and equivalences of graded categories},
Proceedings of the London Mathematical Society (3) \textbf{72} (1996), no.~2, 281--311.
Theorems 1.1 and 1.3.
\url{https://doi.org/10.1112/plms/s3-72.2.281}.
\url{https://deepblue.lib.umich.edu/bitstream/handle/2027.42/135651/plms0281.pdf}.

\bibitem{RRZ}
M. Reyes, D. Rogalski and J.~J. Zhang,
\emph{Skew Calabi--Yau algebras and homological identities},
Advances in Mathematics \textbf{264} (2014), 308--354.
Theorems 0.3 and 5.4; Example 5.5.
\url{https://doi.org/10.1016/j.aim.2014.07.010}.
\url{https://arxiv.org/abs/1302.0437}.

\bibitem{Griffiths}
P.~A. Griffiths,
\emph{On the periods of certain rational integrals.~I},
Annals of Mathematics (2) \textbf{90} (1969), no.~3, 460--495.
Section 8, especially Theorems 8.1 and 8.3 and formula (8.6).
\url{https://doi.org/10.2307/1970746}.
\url{https://www.cimat.mx/~luis/seminarios/Teoria-de-Hodge/period_integral_I.pdf}.

\bibitem{KuznetsovFCY}
A. Kuznetsov,
\emph{Calabi--Yau and fractional Calabi--Yau categories},
Journal f\"ur die reine und angewandte Mathematik \textbf{753} (2019), 239--267.
Lemmas 2.6, 3.13 and 3.19; Corollary 4.1.
\url{https://doi.org/10.1515/crelle-2017-0004}.

\bibitem{CHL}
C.-H. Cho, H. Hong and S.-C. Lau,
\emph{Noncommutative homological mirror functor},
Memoirs of the American Mathematical Society \textbf{271} (2021), no.~1326.
Chapter 5, especially Theorems 5.15 and 5.16.
\url{https://doi.org/10.1090/memo/1326}.
\url{https://arxiv.org/abs/1512.07128}.

\bibitem{Lee}
S. Lee,
\emph{Noncommutative homological mirror symmetry of elliptic curves},
Kyoto Journal of Mathematics \textbf{61} (2021), no.~3, 723--743.
Theorem 6.3 and the corollary in the introduction.
\url{https://arxiv.org/abs/2004.09145}.

\bibitem{CisinskiTabuada}
D.-C. Cisinski and G. Tabuada,
\emph{Lefschetz and Hirzebruch--Riemann--Roch formulas via noncommutative motives},
Journal of Noncommutative Geometry \textbf{8} (2014), 1171--1190.
Section 4 and Example 5.6.
\url{https://doi.org/10.4171/JNCG/183}.
\url{https://www.math.univ-toulouse.fr/~dcisinsk/ncLHRR.pdf}.

\bibitem{MathewHH}
A. Mathew,
\emph{Kaledin's degeneration theorem and topological Hochschild homology},
Geometry \& Topology \textbf{24} (2020), no.~6, 2675--2708.
Theorem 1.3.
\url{https://doi.org/10.2140/gt.2020.24.2675}.
\url{https://arxiv.org/abs/1710.09045}.

\bibitem{PerryIHC}
A. Perry,
\emph{The integral Hodge conjecture for two-dimensional Calabi--Yau categories},
Compositio Mathematica \textbf{158} (2022), no.~2, 287--333.
Lemmas 3.4 and 3.6; Theorem 3.5.
\url{https://doi.org/10.1112/S0010437X22007266}.

\bibitem{MSNE12} I. Mori and S. P. Smith,
\emph{The classification of 3-Calabi--Yau algebras with 3 generators and 3 quadratic relations},
Math. Z. \textbf{287} (2017), 215--241.
Lemma 4.1, Theorem 4.2, Corollary 4.3; arXiv v2, pp. 15--17.
\url{https://arxiv.org/abs/1502.07403}; doi:10.1007/s00209-016-1824-5.
\bibitem{LBSVNE12} L. Le Bruyn, S. P. Smith and M. Van den Bergh,
\emph{Central extensions of three dimensional Artin--Schelter regular algebras},
Math. Z. \textbf{222} (1996), 171--212.
Theorem 2.6, Corollary 2.7, Remark 2.8, Theorem 3.1.3.
\url{https://sites.math.washington.edu/~smith/Research/Central-extns.pdf}.
\bibitem{BTNE12} R. Berger and R. Taillefer,
\emph{Poincar\'e--Birkhoff--Witt deformations of Calabi--Yau algebras},
J. Noncommut. Geom. \textbf{1} (2007), 241--270.
Theorems 3.1 and 3.6, especially the potential bimodule complex in the proof of Theorem 3.6.
\url{https://arxiv.org/abs/math/0610112}; doi:10.4171/JNCG/6.
\bibitem{BergNE12} G. M. Bergman,
\emph{The diamond lemma for ring theory}, Adv. Math. \textbf{29} (1978), 178--218, Theorem 1.2.
\url{https://doi.org/10.1016/0001-8708(78)90010-5}.
\bibitem{LP-AZ} M. Artin and J. J. Zhang,
\emph{Noncommutative projective schemes}, Adv. Math. \textbf{109} (1994), 228--287.
\url{https://doi.org/10.1006/aima.1994.1087}.
\bibitem{LP-BS} J. Burke and G. Stevenson,
\emph{The derived category of a graded Gorenstein ring},
Commutative Algebra and Noncommutative Algebraic Geometry, Vol.~2: Research Articles,
Cambridge University Press (2015), 93--124.
\url{https://doi.org/10.1017/9781316570449.006}.
本稿で確認した版はarXiv:1507.00830v1 (2015年7月3日)。
Definition 5.1、Section 4、Lemmas 6.1--6.2、Theorem 6.4。
\url{https://arxiv.org/abs/1507.00830}.
\bibitem{LP-Keller} B. Keller, \emph{On differential graded categories},
International Congress of Mathematicians, Vol. II (2006), 151--190, Theorem 3.8.
\url{https://arxiv.org/abs/math/0601185}.
\bibitem{LP-Orlov} D. Orlov,
\emph{Smooth and proper noncommutative schemes and gluing of DG categories},
Adv. Math. \textbf{302} (2016), 59--105, Section 3.4, Proposition 3.31.
\url{https://arxiv.org/abs/1402.7364}.
\end{thebibliography}

\end{document}
